% Énergie électrique et puissance électrique — PCT 3ème — Cours signé OnBuch+
% Programme officiel MINESEC. Compiler avec : tectonic N14-energie-puissance-electrique.tex
\newcommand{\DOCMATIERE}{PCT}
\newcommand{\DOCNIVEAU}{3ème}
\newcommand{\DOCMODULE}{Module 2 : Électricité}
\newcommand{\DOCLECON}{N14}
\newcommand{\DOCTITRE}{Énergie électrique et puissance électrique}
% OnBuch+ — préambule « cours signés » ENRICHI (compilé avec tectonic / XeLaTeX)
% Système visuel « Pop » d'OnBuch+ : fond crème (jamais blanc), bordures encre
% épaisses, ombres « dures » décalées, titres Archivo Black en capitales.
% Version MATHÉMATIQUES : graphes (pgfplots/tikz), boîtes pédagogiques
% (définition, propriété, méthode, attention, le savais-tu, exercice résolu,
% exercice, TP, bilan), page de garde et signature OnBuch+ ancrée en bas.
%
% Macros attendues dans main.tex AVANT \input{preamble} :
%   \DOCMATIERE  \DOCNIVEAU  \DOCMODULE  \DOCLECON (numéro)  \DOCTITRE
\documentclass[11pt]{article}

\usepackage{fontspec}
\usepackage[french]{babel}
\usepackage[a4paper,margin=2cm,top=3.4cm,bottom=2.8cm,headheight=1.4cm,footskip=1.3cm]{geometry}
\usepackage{amsmath,amssymb,amsfonts}
\usepackage{mathtools} % \xmapsto, \coloneqq... souvent utilisés par les modèles
\usepackage{cancel} % simplifications barrées (bilans de forces)
% \abs{z} (module, valeur absolue) et \norm{u} (norme), utilisés par les modèles
\providecommand{\abs}{}\let\abs\relax\DeclarePairedDelimiter{\abs}{\lvert}{\rvert}
\providecommand{\norm}{}\let\norm\relax\DeclarePairedDelimiter{\norm}{\lVert}{\rVert}
\usepackage[table]{xcolor} % [table] : fournit aussi \rowcolors (alternance de lignes)
\usepackage{pagecolor}
\usepackage{tikz}
\usetikzlibrary{arrows.meta,positioning,calc,shapes.geometric,shapes.misc,shapes.arrows,decorations.pathmorphing,decorations.pathreplacing,decorations.markings,patterns,shadows,mindmap,trees,fit,backgrounds,matrix,intersections,angles,quotes}
\usepackage{pgfplots}
\usepackage[version=4]{mhchem} % équations nucléaires \ce{^{235}_{92}U}
\usepackage[european,siunitx]{circuitikz} % schémas électriques
\pgfplotsset{compat=1.18}
\usepgfplotslibrary{fillbetween,groupplots} % aires entre courbes (calcul intégral)
\usepackage{enumitem}
\usepackage{fancyhdr}
% [most] charge toutes les bibliothèques tcolorbox (listings, minted, raster,
% documentation...) et gonfle inutilement la mémoire TeX sur un document long
% (~300 boîtes) : on ne charge que ce qui est réellement utilisé.
\usepackage[skins,breakable]{tcolorbox}
\usepackage{graphicx}
\usepackage{colortbl}
\usepackage{booktabs}
\usepackage{tabularx}
\usepackage{multirow}
\usepackage{diagbox} % case d'en-tête diagonale des tableaux à double entrée
% Colonnes extensibles centrée / gauche / droite (C, L, R), souvent utilisées par les modèles
\newcolumntype{C}{>{\centering\arraybackslash}X}
\newcolumntype{L}{>{\raggedright\arraybackslash}X}
\newcolumntype{R}{>{\raggedleft\arraybackslash}X}
\usepackage{array}
\usepackage{multicol}
\usepackage{siunitx}
% parse-numbers=false : accepte aussi \SI{2,5 \times 10^{-3}}{...}
\sisetup{locale=FR,output-decimal-marker={,},group-minimum-digits=4,parse-numbers=false}
\DeclareSIUnit\ohm{\ensuremath{\symup{\Omega}}} % U+2126 absent de la police
\DeclareSIUnit\division{div} % réglages d'oscilloscope (ms/div, V/div)
\DeclareSIUnit\tr{tr} % tours (vitesse de rotation, ex. \SI{1500}{\tr\per\minute})
\DeclareSIUnit\molar{M} % concentration molaire (ex. \SI{3}{\molar}, \SI{1}{\micro\molar})
\DeclareSIUnit\an{an} % année (le modèle confond parfois avec \year, un compteur TeX) — ex. \SI{5}{\kilo\meter\per\an}
\DeclareSIUnit\FCFA{FCFA} % franc CFA (ex. \SI{500}{\FCFA\per\kilo\gram})
\DeclareSIUnit\degreeBrix{\textdegree Brix} % teneur en sucre d'un sirop/jus (réfractométrie)
\DeclareSIUnit\month{mois} % le modèle utilise parfois \month, un compteur TeX, comme unité de durée
\DeclareSIUnit\microliter{\micro\liter} % le modèle écrit parfois \microliter en un seul token
\DeclareSIUnit\million{millions} % dénombrement (ex. \SI{15}{\million\per\mL} spermatozoïdes)
\usepackage{qrcode}
% \degree : commande fréquemment utilisée par les modèles pour un angle ou une
% température en dehors de \SI{}{} (non fournie par les paquets chargés).
\providecommand{\degree}{^{\circ}}
% \qed (fin de démonstration), utilisé par les modèles sans amsthm
\providecommand{\qed}{\hfill$\square$}
% \barre : double barre des valeurs interdites dans un tableau de variations (tkz-tab)
\providecommand{\barre}{\ensuremath{\|}}
% environnement proof (amsthm non chargé) utilisé par les modèles
\ifdefined\proof\else\newenvironment{proof}[1][Démonstration]{\par\noindent\textit{#1.}\ }{\qed\par}\fi
\usepackage{lastpage}
\usepackage{hyperref}
\hypersetup{hidelinks}

% ============================================================
% Palette « Pop » OnBuch+ — mêmes hex que la classe P (plus_ui.dart)
% ============================================================
\definecolor{poporange}{HTML}{F59321}
\definecolor{popdark}{HTML}{E07A0C}
\definecolor{popcream}{HTML}{FFF6E8}
\definecolor{popink}{HTML}{1C1714}
\definecolor{poppurple}{HTML}{7A5AE0}
\definecolor{popgreen}{HTML}{2E9E6A}
\definecolor{poppink}{HTML}{E0457B}
\definecolor{popblue}{HTML}{2F7DE1}
\definecolor{popgold}{HTML}{C98A12}
\definecolor{popmuted}{HTML}{8A7F76}
\definecolor{popline}{HTML}{EADFD2}
\definecolor{popgray}{HTML}{8A7F76}
\colorlet{popgrey}{popgray}
% Alias tolérants (le modèle nomme parfois les couleurs différemment)
\colorlet{popwhite}{white}
\colorlet{popblack}{popink}
\colorlet{popred}{poppink}
\colorlet{popyellow}{popgold}
% Teintes claires (fonds de tableaux/encadrés) — le modèle nomme parfois
% les couleurs avec un suffixe L (ex. popblueL) sans les avoir définies.
\colorlet{poporangeL}{poporange!20}
\colorlet{popdarkL}{popdark!20}
\colorlet{poppurpleL}{poppurple!20}
\colorlet{popgreenL}{popgreen!20}
\colorlet{poppinkL}{poppink!20}
\colorlet{popblueL}{popblue!20}
\colorlet{popgoldL}{popgold!20}
\colorlet{popredL}{popred!20}
\colorlet{popgrayL}{popgray!20}
% Teintes claires pour fonds de tableaux / zones
\colorlet{popblueL}{popblue!10!white}
\colorlet{poporangeL}{poporange!14!white}
\colorlet{popgreenL}{popgreen!12!white}
\colorlet{poppurpleL}{poppurple!12!white}
\colorlet{poppinkL}{poppink!12!white}
\colorlet{popgoldL}{popgold!14!white}

\pagecolor{popcream}

% ============================================================
% Typographie
% ============================================================
\defaultfontfeatures{Ligatures=TeX}
\setmainfont{PlusJakartaSans-Medium.ttf}[Path=../../../fonts/,BoldFont=PlusJakartaSans-Bold.ttf]
% Maths en sans-serif (Fira Math) avec chiffres et lettres droites de Plus
% Jakarta, pour que les formules se fondent dans le texte.
\usepackage[math-style=ISO,bold-style=ISO]{unicode-math}
\setmathfont{FiraMath-Regular.otf}[Path=../../../fonts/]
\setmathfont{PlusJakartaSans-Medium.ttf}[Path=../../../fonts/,range={up/num,up/{Latin,latin},\mathup}]
\usepackage{icomma} % virgule décimale sans espace en mode math
% Caractères mathématiques Unicode absents des polices de texte (Plus Jakarta,
% Archivo Black) : rendus par leur équivalent mathématique, en texte comme en
% formule — sinon ils s'affichent en carré vide (ex. « Arithmétique dans ℤ »).
\usepackage{newunicodechar}
\newunicodechar{ℝ}{\ensuremath{\mathbb{R}}}
\newunicodechar{ℤ}{\ensuremath{\mathbb{Z}}}
\newunicodechar{ℂ}{\ensuremath{\mathbb{C}}}
\newunicodechar{ℕ}{\ensuremath{\mathbb{N}}}
\newunicodechar{ℚ}{\ensuremath{\mathbb{Q}}}
\newunicodechar{𝔻}{\ensuremath{\mathbb{D}}}
\newunicodechar{α}{\ensuremath{\alpha}}
\newunicodechar{β}{\ensuremath{\beta}}
\newunicodechar{γ}{\ensuremath{\gamma}}
\newunicodechar{δ}{\ensuremath{\delta}}
\newunicodechar{ε}{\ensuremath{\varepsilon}}
\newunicodechar{θ}{\ensuremath{\theta}}
\newunicodechar{λ}{\ensuremath{\lambda}}
\newunicodechar{μ}{\ensuremath{\mu}}
\newunicodechar{σ}{\ensuremath{\sigma}}
\newunicodechar{τ}{\ensuremath{\tau}}
\newunicodechar{φ}{\ensuremath{\varphi}}
\newunicodechar{ψ}{\ensuremath{\psi}}
\newunicodechar{ω}{\ensuremath{\omega}}
\newunicodechar{Σ}{\ensuremath{\Sigma}}
% majuscules produites par \MakeUppercase dans les titres (α-aminés -> Α)
\newunicodechar{Α}{\ensuremath{\alpha}}
\newunicodechar{Β}{\ensuremath{\beta}}
\newunicodechar{Γ}{\ensuremath{\Gamma}}
\newunicodechar{Δ}{\ensuremath{\Delta}}
\newunicodechar{Ε}{\ensuremath{\varepsilon}}
\newunicodechar{Θ}{\ensuremath{\Theta}}
\newunicodechar{Λ}{\ensuremath{\Lambda}}
\newunicodechar{Μ}{\ensuremath{\mu}}
\newunicodechar{Τ}{\ensuremath{\tau}}
\newunicodechar{Φ}{\ensuremath{\Phi}}
\newunicodechar{Ψ}{\ensuremath{\Psi}}
\newunicodechar{Ω}{\ensuremath{\Omega}}
\newunicodechar{↦}{\ensuremath{\mapsto}}
\newunicodechar{∈}{\ensuremath{\in}}
\newunicodechar{∉}{\ensuremath{\notin}}
\newunicodechar{∧}{\ensuremath{\wedge}}
\newunicodechar{⇔}{\ensuremath{\Leftrightarrow}}
\newunicodechar{⟺}{\ensuremath{\Longleftrightarrow}}
\newunicodechar{⇒}{\ensuremath{\Rightarrow}}
\newunicodechar{⟹}{\ensuremath{\Longrightarrow}}
\newunicodechar{∘}{\ensuremath{\circ}}
\newunicodechar{∪}{\ensuremath{\cup}}
\newunicodechar{∩}{\ensuremath{\cap}}
\newunicodechar{≡}{\ensuremath{\equiv}}
\newunicodechar{⊂}{\ensuremath{\subset}}
\newunicodechar{⊥}{\ensuremath{\perp}}
\newunicodechar{∃}{\ensuremath{\exists}}
\newunicodechar{∀}{\ensuremath{\forall}}
\newunicodechar{∛}{\ensuremath{\sqrt[3]{\,}}}
\newunicodechar{∜}{\ensuremath{\sqrt[4]{\,}}}
\newunicodechar{⃗}{\ensuremath{{}^{\to}}}
\newunicodechar{ⁿ}{\ifmmode^{n}\else\textsuperscript{\ensuremath{n}}\fi}
\newunicodechar{ˣ}{\ifmmode^{x}\else\textsuperscript{\ensuremath{x}}\fi}
\newunicodechar{ᵇ}{\ifmmode^{b}\else\textsuperscript{\ensuremath{b}}\fi}
\newunicodechar{ʳ}{\ifmmode^{r}\else\textsuperscript{\ensuremath{r}}\fi}
\newunicodechar{ᵘ}{\ifmmode^{u}\else\textsuperscript{\ensuremath{u}}\fi}
\newunicodechar{ʸ}{\ifmmode^{y}\else\textsuperscript{\ensuremath{y}}\fi}
\newunicodechar{ᵃ}{\ifmmode^{a}\else\textsuperscript{\ensuremath{a}}\fi}
\newunicodechar{ᵉ}{\ifmmode^{e}\else\textsuperscript{\ensuremath{e}}\fi}
\newunicodechar{ᵏ}{\ifmmode^{k}\else\textsuperscript{\ensuremath{k}}\fi}
\newunicodechar{ᵖ}{\ifmmode^{p}\else\textsuperscript{\ensuremath{p}}\fi}
\newunicodechar{ᴺ}{\ifmmode^{N}\else\textsuperscript{\ensuremath{N}}\fi}
\newunicodechar{ᶜ}{\ifmmode^{c}\else\textsuperscript{\ensuremath{c}}\fi}
\newunicodechar{ʰ}{\ifmmode^{h}\else\textsuperscript{\ensuremath{h}}\fi}
\newunicodechar{ᵠ}{\ifmmode^{q}\else\textsuperscript{\ensuremath{q}}\fi}
\newunicodechar{ₙ}{\ifmmode_{n}\else\textsubscript{\ensuremath{n}}\fi}
\newunicodechar{ₐ}{\ifmmode_{a}\else\textsubscript{\ensuremath{a}}\fi}
\newunicodechar{ᵢ}{\ifmmode_{i}\else\textsubscript{\ensuremath{i}}\fi}
\newunicodechar{ᵣ}{\ifmmode_{r}\else\textsubscript{\ensuremath{r}}\fi}
\newunicodechar{ₖ}{\ifmmode_{k}\else\textsubscript{\ensuremath{k}}\fi}
\newunicodechar{ᵦ}{\ifmmode_{\beta}\else\textsubscript{\ensuremath{\beta}}\fi}
\newunicodechar{ₓ}{\ifmmode_{x}\else\textsubscript{\ensuremath{x}}\fi}
\newfontfamily\popdisplay{ArchivoBlack-Regular.ttf}[Path=../../../fonts/]
\newfontfamily\poplogo{SpaceGrotesk-Bold.ttf}[Path=../../../fonts/]
\newfontfamily\popsemi{PlusJakartaSans-SemiBold.ttf}[Path=../../../fonts/]
\newfontfamily\popxbold{PlusJakartaSans-ExtraBold.ttf}[Path=../../../fonts/]
\newfontfamily\popxboldls{PlusJakartaSans-ExtraBold.ttf}[Path=../../../fonts/,LetterSpace=3.0]

\setlist[enumerate]{leftmargin=1.7em,itemsep=5pt,topsep=4pt}
\setlist[itemize]{leftmargin=1.7em,itemsep=4pt,topsep=4pt,label={\color{poporange}\textbullet}}
\setlength{\parindent}{0pt}
\setlength{\parskip}{5pt}
\renewcommand{\arraystretch}{1.25}

% pgfplots : style Pop par défaut
\pgfplotsset{
  every axis/.append style={axis line style={popink,line width=1pt},tick style={popink},
    grid style={popline,line width=0.5pt},label style={font=\small},tick label style={font=\footnotesize}},
  popcurve/.style={poporange,line width=1.8pt,no markers,smooth},
}
% Même style disponible dans un \draw TikZ simple (hors pgfplots)
\tikzset{popcurve/.style={poporange,line width=1.8pt}}
\tikzset{popgrid/.style={popline,very thin}}
\colorlet{popcurve}{poporange} % le nom du style est parfois utilisé comme couleur

% ============================================================
% Pill / squiggle
% ============================================================
\newcommand{\poppill}[3]{%
  \tikz[baseline=-0.55ex]\node[draw=popink,line width=1.2pt,rounded corners=2.6pt,%
    fill=#2,inner xsep=6.5pt,inner ysep=3pt]{%
    \popxboldls\fontsize{7}{7}\selectfont\color{#3}\MakeUppercase{#1}};%
}
\newcommand{\popsquiggle}[1]{%
  \tikz[baseline=-0.4ex]{
    \draw[#1,line width=2.6pt,line cap=round]
      plot [smooth] coordinates {(0,0.09) (0.34,0.01) (0.68,0.21) (1.02,0.01) (1.36,0.10)};
  }%
}
% Mot-clé surligné (marqueur orange sous le texte)
\newcommand{\cle}[1]{{\popxbold\color{popdark}#1}}

% ============================================================
% En-tête / pied
% ============================================================
\pagestyle{fancy}
\fancyhf{}
\renewcommand{\headrulewidth}{0pt}
\renewcommand{\footrulewidth}{0pt}
\lhead{\raisebox{-0.15ex}{\poplogo\fontsize{13}{13}\selectfont\color{popink}OnBuch\color{poporange}+}}
\rhead{\poppill{\DOCMATIERE\ \textbullet\ \DOCNIVEAU\ \textbullet\ Leçon \DOCLECON}{white}{popink}}
\lfoot{\popxbold\fontsize{7.6}{7.6}\selectfont\color{popmuted}\MakeUppercase{Cours signé OnBuch+}\ \textbullet\ {\popsemi\DOCTITRE}}
\rfoot{\tikz[baseline=-0.6ex]\node[rounded corners=8.5pt,fill=popink,minimum height=17pt,minimum width=17pt,inner xsep=5pt,inner sep=0]{\popxbold\fontsize{8}{8}\selectfont\color{popcream}\thepage\,/\,\pageref*{LastPage}};}

% ============================================================
% Titres
% ============================================================
\newcommand{\chaptitre}[2]{%
  \begin{tcolorbox}[enhanced,colback=poporange,colframe=popink,boxrule=2.6pt,arc=11pt,
      drop shadow=popink,left=15pt,right=15pt,top=12pt,bottom=11pt,before skip=3pt,after skip=18pt]
    \poppill{#2}{popink}{popcream}\par\vspace{9pt}
    {\raggedright\hyphenpenalty=10000\exhyphenpenalty=10000\popdisplay\fontsize{22}{24}\selectfont\color{popcream}\MakeUppercase{#1}\par}
  \end{tcolorbox}
}
% Section numérotée : pastille orange avec le numéro + titre Archivo
\newcounter{popsec}
\newcommand{\coursec}[1]{%
  \stepcounter{popsec}\setcounter{popsub}{0}%
  \par\vspace{16pt}\noindent
  \begin{minipage}{\linewidth}
  \tikz[baseline=-0.7ex]\node[draw=popink,line width=1.6pt,fill=poporange,rounded corners=4pt,minimum size=22pt,inner sep=0]{\popdisplay\fontsize{12}{12}\selectfont\color{popcream}\thepopsec};%
  \hspace{8pt}{\popdisplay\fontsize{15}{17}\selectfont\color{popink}\MakeUppercase{#1}}\par
  \vspace{4pt}\noindent{\color{poporange}\rule{36pt}{3.6pt}}
  \end{minipage}\par\nopagebreak\vspace{8pt}
}
\newcounter{popsub}
\newcommand{\courssub}[1]{%
  \stepcounter{popsub}%
  \par\vspace{9pt}\noindent{\popxbold\fontsize{11.5}{13}\selectfont\color{popink}{\color{poporange}\thepopsec.\thepopsub}\hspace{6pt}#1}\par\nopagebreak\vspace{4pt}
}

% ============================================================
% Boîtes pédagogiques — même recette : fond blanc, bordure épaisse colorée,
% ombre dure encre, badge plein. Titre optionnel [#1].
% ============================================================
\tcbset{popbase/.style={breakable,enhanced,colback=white,boxrule=2.3pt,arc=9pt,
  shadow={3pt}{-3pt}{0pt}{popink},left=14pt,right=14pt,top=11pt,bottom=11pt,before skip=11pt,after skip=15pt,
  fontupper=\popsemi\fontsize{10.6}{14.5}\selectfont\color{popink},
  fontlower=\popsemi\fontsize{10.6}{14.5}\selectfont\color{popink}}}
% Coche dessinée (le glyphe ✓ n'existe pas dans les polices du document)
\newcommand{\popcheck}[1]{\tikz[baseline=-0.5ex]\draw[#1,line width=1.6pt,line cap=round,line join=round](-0.13,0.0)--(-0.04,-0.09)--(0.13,0.1);}
\providecommand{\coche}{\popcheck{popgreen}}
\providecommand{\resultat}[1]{\par\smallskip\noindent\fcolorbox{popgreen}{popgreenL}{\parbox{\dimexpr\linewidth-2\fboxsep-2\fboxrule\relax}{\textbf{Résultat.} #1}}\par\smallskip}
\newcommand{\popboxhead}[3]{\poppill{#1}{#2}{white}\ifx\relax#3\relax\else\hspace{6pt}{\popxbold\fontsize{10.6}{12}\selectfont\color{popink}#3}\fi\par\vspace{7pt}}

\newtcolorbox{aretenir}[1][]{popbase,colframe=popgreen,before upper={\popboxhead{À retenir}{popgreen}{#1}}}
\newtcolorbox{exemplebox}[1][]{popbase,colframe=poporange,before upper={\popboxhead{Exemple}{poporange}{#1}}}
\newtcolorbox{definition}[1][]{popbase,colframe=popblue,before upper={\popboxhead{Définition}{popblue}{#1}}}
\newtcolorbox{propriete}[1][]{popbase,colframe=poppurple,before upper={\popboxhead{Propriété}{poppurple}{#1}}}
\newtcolorbox{methode}[1][]{popbase,colframe=popgold,before upper={\popboxhead{Méthode}{popgold}{#1}}}
\newtcolorbox{attention}[1][]{popbase,colframe=poppink,before upper={\popboxhead{Attention}{poppink}{#1}}}
\newtcolorbox{experience}[1][]{popbase,colframe=popblue,colback=popblueL,before upper={\popboxhead{Expérience}{popblue}{#1}}}
% « Le savais-tu ? » : carte violette pleine (contexte, culture, Cameroun)
\newtcolorbox{savaistu}[1][]{popbase,colback=poppurple,colframe=popink,
  fontupper=\popsemi\fontsize{10.4}{14.2}\selectfont\color{white},
  before upper={\poppill{Le savais-tu\,?}{white}{poppurple}\ifx\relax#1\relax\else\hspace{6pt}{\popxbold\color{white}#1}\fi\par\vspace{7pt}}}
% Exercice résolu : énoncé en haut, \tcblower puis la solution sur fond crème
\newcounter{popexo}\newcounter{popexr}
\newtcolorbox{exoresolu}[1][]{popbase,colframe=poporange,colbacklower=popcream,
  segmentation style={popink,line width=1.2pt,dashed},
  before upper={\stepcounter{popexr}\popboxhead{Exercice résolu \thepopexr}{poporange}{#1}},
  before lower={\poppill{Solution}{popink}{popcream}\par\vspace{6pt}}}
% Exercice (non résolu en place) — la correction est dans la partie Corrigés
\newtcolorbox{exercice}[1][]{popbase,colframe=popink,
  before upper={\stepcounter{popexo}\popboxhead{Exercice \thepopexo}{popink}{#1}}}
% Corrigé
\newtcolorbox{corrige}[1][]{popbase,colframe=popgreen,colback=popgreenL,
  before upper={\popboxhead{Corrigé}{popgreen}{#1}}}
% Objectifs / prérequis / bilan
\newtcolorbox{objectifs}{popbase,colframe=popink,colback=poporangeL,
  before upper={\popboxhead{Objectifs de la leçon}{popink}{}}}
\newtcolorbox{prerequis}{popbase,colframe=popink,colback=popblueL,
  before upper={\popboxhead{Prérequis}{popblue}{}}}
\newtcolorbox{bilan}{popbase,colframe=popink,colback=white,boxrule=2.8pt,
  before upper={\popboxhead{Fiche bilan}{poporange}{L'essentiel à connaître pour l'examen}}}
% Figure encadrée avec légende
\newtcolorbox{popfigure}[1][]{enhanced,breakable=false,colback=white,colframe=popline,boxrule=1.4pt,arc=8pt,
  left=10pt,right=10pt,top=10pt,bottom=8pt,before skip=10pt,after skip=12pt,halign=center,
  lower separated=false,halign lower=center,
  fontlower=\popsemi\fontsize{9}{11.5}\selectfont\color{popmuted},
  #1}
\newcommand{\legende}[1]{\par\vspace{6pt}{\popsemi\fontsize{9}{11.5}\selectfont\color{popmuted}#1}}

% ============================================================
% Signature OnBuch+
% ============================================================
\newcommand{\poplogoinline}[1]{{\poplogo\fontsize{#1}{#1}\selectfont\color{popink}OnBuch\color{poporange}+}}
\newcommand{\signatureonbuch}{%
  \begin{tcolorbox}[enhanced,colback=popink,colframe=popink,boxrule=2.2pt,arc=10pt,
      drop shadow=poporange,left=14pt,right=14pt,top=10pt,bottom=10pt,before skip=0pt,after skip=0pt]
    \tikz[baseline=-0.7ex]\node[circle,fill=popgreen,draw=popcream,line width=1.3pt,minimum size=22pt,inner sep=0]{\popcheck{white}};%
    \hspace{9pt}\begin{minipage}[c]{0.62\linewidth}
      {\popxbold\fontsize{12}{13}\selectfont\color{popcream}Signé\ \textbullet\ {\poplogo OnBuch\color{poporange}+}}\par\vspace{2pt}
      {\raggedright\popsemi\fontsize{8}{10.5}\selectfont\color{popline}\MakeUppercase{Équipe pédagogique OnBuch+}\\\MakeUppercase{Conforme au programme officiel MINESEC}\par}
    \end{minipage}\hfill
    {\poplogo\fontsize{22}{22}\selectfont\color{popcream}OnBuch\color{poporange}+}
  \end{tcolorbox}%
}

% ============================================================
% Page de garde
% #1 titre · #2 matière · #3 niveau · #4 module · #5 n° leçon · #6 durée ·
% #7 description / objectifs (texte ou itemize)
% ============================================================
\newcommand{\pagegarde}[7]{%
  \thispagestyle{empty}
  \newgeometry{a4paper,margin=1.7cm,top=1.6cm,bottom=1.4cm}%
  \begin{tikzpicture}[remember picture,overlay]
    \fill[poporange!18] ([xshift=-3.2cm,yshift=-3.6cm]current page.north east) circle (4.6cm);
    \fill[poppurple!14] ([xshift=2.4cm,yshift=7.6cm]current page.south west) circle (3.8cm);
  \end{tikzpicture}%
  {\poplogo\fontsize{30}{30}\selectfont\color{popink}OnBuch\color{poporange}+}\hfill
  \raisebox{0.6ex}{\poppill{Cours signé \textbullet\ Premium}{popink}{popcream}}\par
  \vspace{1pt}\popsquiggle{poppurple}\par
  \vspace{8pt}
  \begin{tcolorbox}[enhanced,colback=poporange,colframe=popink,boxrule=2.8pt,arc=13pt,
      drop shadow=popink,left=18pt,right=18pt,top=12pt,bottom=13pt,after skip=10pt]
    \poppill{#2 \textbullet\ #3}{popink}{popcream}\hspace{5pt}\poppill{#4}{white}{popink}\par\vspace{8pt}
    {\popdisplay\fontsize{14}{14}\selectfont\color{popink}LEÇON #5}\par\vspace{4pt}
    {\raggedright\hyphenpenalty=10000\exhyphenpenalty=10000\popdisplay\fontsize{25}{27}\selectfont\color{popcream}\MakeUppercase{#1}\par}
    \vspace{8pt}
    {\popxbold\fontsize{9.5}{9.5}\selectfont\color{popink}\MakeUppercase{Durée conseillée : #6}}
  \end{tcolorbox}
  \begin{tcolorbox}[enhanced,colback=white,colframe=popink,boxrule=2.2pt,arc=10pt,
      drop shadow=popink,left=16pt,right=16pt,top=10pt,bottom=10pt,after skip=8pt]
    \poppill{Dans cette leçon}{poppurple}{white}\par\vspace{5pt}
    \popsemi\fontsize{9.3}{12.6}\selectfont\color{popink}#7
  \end{tcolorbox}
  \noindent\begin{minipage}[t]{0.487\linewidth}
    \begin{tcolorbox}[enhanced,colback=popgreen,colframe=popink,boxrule=2.2pt,arc=10pt,
        drop shadow=popink,left=11pt,right=11pt,top=10pt,bottom=10pt,height=3.1cm,valign=center]
      \tcbox[colback=white,colframe=popink,boxrule=1.3pt,arc=5pt,boxsep=0pt,nobeforeafter,
        left=3pt,right=3pt,top=3pt,bottom=3pt,valign=center]{\qrcode[height=1.9cm]{https://play.google.com/store/apps/details?id=com.luvvix.onbuch&hl=fr}}%
      \hspace{8pt}\begin{minipage}[c]{0.5\linewidth}
        {\popdisplay\fontsize{11}{12.5}\selectfont\color{white}\MakeUppercase{Télécharge OnBuch}}\par\vspace{4pt}
        {\raggedright\popsemi\fontsize{8.4}{11}\selectfont\color{white}Gratuit \textbullet\ Google Play\\Tuteur IA, annales, concours, résultats\par}
      \end{minipage}
    \end{tcolorbox}
  \end{minipage}\hfill
  \begin{minipage}[t]{0.487\linewidth}
    \begin{tcolorbox}[enhanced,colback=poppurple,colframe=popink,boxrule=2.2pt,arc=10pt,
        drop shadow=popink,left=11pt,right=11pt,top=10pt,bottom=10pt,height=3.1cm,valign=center]
      \tikz[baseline=-0.7ex]\node[circle,fill=white,draw=popink,line width=1.3pt,minimum size=1.6cm,inner sep=0]{\popdisplay\fontsize{24}{24}\selectfont\color{poppurple}+};%
      \hspace{8pt}\begin{minipage}[c]{0.6\linewidth}
        {\popdisplay\fontsize{11}{12.5}\selectfont\color{white}\MakeUppercase{Passe à OnBuch+}}\par\vspace{4pt}
        {\raggedright\popsemi\fontsize{8.4}{11}\selectfont\color{white}Tous les cours signés, Tuteur IA illimité, fiches PDF — onglet OnBuch+ de l'app\par}
      \end{minipage}
    \end{tcolorbox}
  \end{minipage}
  \vspace{8pt}
  \popcontenu
  \vfill
  \signatureonbuch
  \restoregeometry
  \newpage
}

% Rangée « Ce cours contient » (page de garde)
\newcommand{\poptile}[3]{% #1 couleur #2 gros texte #3 légende
  \begin{tcolorbox}[enhanced,colback=#1,colframe=popink,boxrule=2pt,arc=9pt,shadow={2.5pt}{-2.5pt}{0pt}{popink},
      left=6pt,right=6pt,top=9pt,bottom=9pt,halign=center,valign=center,height=1.75cm,width=\linewidth,nobeforeafter]
    {\popdisplay\fontsize{12.5}{13}\selectfont\color{white}\MakeUppercase{#2}}\par\vspace{3pt}
    {\centering\popsemi\fontsize{7.8}{9.5}\selectfont\color{white}#3\par}
  \end{tcolorbox}}
\newcommand{\popcontenu}{%
  \par\noindent{\popxboldls\fontsize{7.5}{7.5}\selectfont\color{popmuted}CE COURS SIGNÉ CONTIENT}\par\vspace{7pt}
  \noindent\begin{minipage}[t]{0.235\linewidth}\poptile{popblue}{Cours}{complet, illustré, pas à pas}\end{minipage}\hfill
  \begin{minipage}[t]{0.235\linewidth}\poptile{poporange}{Méthodes}{et exercices résolus}\end{minipage}\hfill
  \begin{minipage}[t]{0.235\linewidth}\poptile{poppink}{Exercices}{type examen + corrigés}\end{minipage}\hfill
  \begin{minipage}[t]{0.235\linewidth}\poptile{popgreen}{Fiche bilan}{l'essentiel à retenir}\end{minipage}\par}

% Sommaire
\newtcolorbox{sommaire}{enhanced,colback=white,colframe=popink,boxrule=2.2pt,arc=10pt,
  drop shadow=popink,left=18pt,right=18pt,top=13pt,bottom=13pt,before skip=6pt,after skip=20pt,
  before upper={\poppill{Sommaire}{poppurple}{white}\par\vspace{9pt}\popsemi\fontsize{10.6}{14.5}\selectfont\color{popink}}}

% Signature de fin de cours — poussée tout en bas de la dernière page
\newcommand{\findecours}{%
  \par\vfill\nopagebreak
  \begin{center}\popsquiggle{poporange}\end{center}\vspace{4pt}
  \signatureonbuch
}
\AtBeginDocument{\def\llbracket{\mathopen{[\mkern-3.5mu[}}\def\rrbracket{\mathclose{]\mkern-3.5mu]}}} % intervalles d'entiers (glyphe ⟦ absent de la police)
% Glyphes absents de Fira Math : redéfinis à partir de glyphes disponibles
\AtBeginDocument{%
  \def\setminus{\mathbin{\backslash}}\let\smallsetminus\setminus
  \def\vdots{\mathord{\vbox{\baselineskip3.2pt\lineskiplimit0pt\kern4pt\hbox{.}\hbox{.}\hbox{.}}}}%
  \def\ddots{\mathinner{\mkern1mu\raise6.5pt\hbox{.}\mkern2mu\raise3.6pt\hbox{.}\mkern2mu\raise0.7pt\hbox{.}\mkern1mu}}%
  \def\triangle{\mathord{\scalebox{0.9}{$\symup{\Delta}$}}}%
  \def\nabla{\mathord{\raisebox{\depth}{\scalebox{1}[-1]{$\symup{\Delta}$}}}}%
  \def\checkmark{\popcheck{popdark}}%
  \def\star{\mathbin{\ast}}\def\bigstar{\mathord{\ast}}%
  \def\diamond{\mathbin{\scalebox{0.6}{$\Diamond$}}}%
  \def\bigcup{\mathop{\mathchoice{\raisebox{-0.3em}{\scalebox{1.7}{$\cup$}}}{\scalebox{1.25}{$\cup$}}{\cup}{\cup}}\displaylimits}%
  \def\bigcap{\mathop{\mathchoice{\raisebox{-0.3em}{\scalebox{1.7}{$\cap$}}}{\scalebox{1.25}{$\cap$}}{\cap}{\cap}}\displaylimits}%
  \def\lesseqqgtr{\mathrel{\lessgtr}}\def\bigoplus{\mathop{\oplus}\displaylimits}%
}
\providecommand{\ssi}{si et seulement si} % abréviation fréquente
\AtBeginDocument{\providecommand{\wideparen}{\overparen}} % arc de cercle

%% FIGFIX-BEGIN v5 — correctifs globaux des figures (docs/onbuch-generation/tools/figfix)
\usepackage{adjustbox}\usepackage{etoolbox}\tcbuselibrary{raster}
% toute figure TikZ/pgfplots plus large que la ligne est réduite à la largeur disponible
\BeforeBeginEnvironment{tikzpicture}{\begin{adjustbox}{max width=\linewidth}}
\AfterEndEnvironment{tikzpicture}{\end{adjustbox}}
% légendes pgfplots lisibles par-dessus les courbes
\pgfplotsset{every axis/.append style={legend style={fill=white,fill opacity=0.92,text opacity=1,draw=black!15,inner sep=3pt}}}
% étiquettes d'arêtes (syntaxe edge["..."]) avec fond blanc
\tikzset{every edge quotes/.append style={fill=white,inner sep=1.2pt}}
%% FIGFIX-END
\begin{document}

\pagegarde{Énergie électrique et puissance électrique}{PCT}{3ème}{Module 2 : Électricité}{N14}{2 séances (1 h chacune)}{Cette leçon introduit les grandeurs électriques de base, leurs unités et instruments de mesure, puis établit les formules de la puissance électrique pour un résistor. Elle se termine par des applications concrètes et les consignes de sécurité indispensables.
\vspace{4pt}
\begin{multicols}{2}\small
\begin{itemize}
\item À la fin de cette leçon, je sais définir la tension, le courant, la résistance et la puissance avec leurs unités SI.
\item À la fin de cette leçon, je sais brancher correctement un voltmètre, un ampèremètre et un wattmètre dans un circuit.
\item À la fin de cette leçon, je sais exprimer la puissance électrique d'un résistor sous les formes P = U·I, P = R·I² et P = U²/R.
\item À la fin de cette leçon, je sais calculer l'énergie consommée en kilowattheure à partir de la puissance et du temps d'utilisation.
\item À la fin de cette leçon, je sais lire un compteur électrique et estimer le coût d'une consommation donnée.
\item À la fin de cette leçon, je sais identifier les principaux risques électriques et les dispositifs de protection (disjoncteur, fusible, prise de terre).
\end{itemize}\end{multicols}}

\begin{sommaire}
\begin{multicols}{2}
\begin{enumerate}
\item Grandeurs électriques et instruments de mesure
\item Puissance électrique d'un résistor
\item Mesure de la puissance et de l'énergie consommée
\item Sécurité électrique et consignes de prévention
\item Applications concrètes et résolution de problèmes
\item Activité d'intégration
\item Méthodes et exercices résolus
\item Exercices
\item Corrigés des exercices
\item Fiche bilan
\end{enumerate}
\end{multicols}
\end{sommaire}

\begin{objectifs}
\begin{itemize}
\item À la fin de cette leçon, je sais définir la tension, le courant, la résistance et la puissance avec leurs unités SI.
\item À la fin de cette leçon, je sais brancher correctement un voltmètre, un ampèremètre et un wattmètre dans un circuit.
\item À la fin de cette leçon, je sais exprimer la puissance électrique d'un résistor sous les formes P = U·I, P = R·I² et P = U²/R.
\item À la fin de cette leçon, je sais calculer l'énergie consommée en kilowattheure à partir de la puissance et du temps d'utilisation.
\item À la fin de cette leçon, je sais lire un compteur électrique et estimer le coût d'une consommation donnée.
\item À la fin de cette leçon, je sais identifier les principaux risques électriques et les dispositifs de protection (disjoncteur, fusible, prise de terre).
\item À la fin de cette leçon, je sais appliquer les consignes de sécurité avant toute intervention sur une installation électrique.
\item À la fin de cette leçon, je sais résoudre un problème concret : choisir une ampoule, dimensionner un câble ou proposer des économies d'énergie.
\end{itemize}
\end{objectifs}

\begin{prerequis}
\begin{itemize}
\item Notion de charge électrique et de courant (classe de 4e).
\item Loi d'Ohm U = R·I (début d'année 3e).
\item Unités de base du SI : volt, ampère, ohm, watt, joule.
\item Manipulation simple d'un multimètre (mode tension et courant).
\item Notion d'énergie et de travail (physique 4e).
\end{itemize}
\end{prerequis}

\newpage

\coursec{Énergie électrique et puissance électrique}

Dans le quartier de Nkolbisson, la famille Mbarga vient de recevoir sa facture ENEO : elle a \textbf{doublé} en un seul mois ! Le père fronce les sourcils : « C'est ce vieux réfrigérateur qui consomme trop ! » La mère n'est pas d'accord : « Non, ce sont les nouvelles lampes halogènes du salon ! » Les enfants, curieux, décident de mener l'enquête comme de vrais scientifiques. Leur plan : mesurer la \cle{tension} aux bornes de chaque appareil, mesurer l'\cle{intensité} du courant qui le traverse, puis calculer sa \cle{puissance} réelle et l'\cle{énergie} consommée. Mais comment mesure-t-on ces grandeurs ? Avec quels appareils ? Et comment brancher correctement sans danger ?

\textbf{Problématique :} quelles sont les grandeurs électriques fondamentales, comment les mesure-t-on, comment calcule-t-on la puissance et l'énergie consommée, et quelles sont les règles de sécurité à respecter ?

\courssub{Grandeurs électriques et instruments de mesure}

\courssub{La tension électrique}

\textbf{Intuition.} Imagine deux réservoirs d'eau reliés par un tuyau : l'eau ne coule que s'il existe une \emph{différence de niveau} entre les deux réservoirs. En électricité, c'est pareil : le courant ne circule entre deux points d'un circuit que s'il existe une \emph{différence d'état électrique} entre ces deux points. C'est le générateur (pile, prise ENEO, groupe électrogène) qui crée et maintient cette différence, comme une pompe maintient la différence de niveau d'eau.

\begin{definition}[Tension électrique]
La \cle{tension électrique} entre deux points A et B d'un circuit est la \textbf{différence de potentiel électrique} entre ces deux points. On la note $U_{AB}$.
\begin{itemize}
\item Son unité est le \cle{volt}, de symbole \SI{}{V}.
\item On mesure une tension avec un \cle{voltmètre}, branché \textbf{en dérivation} (en parallèle) entre les deux points.
\end{itemize}
Unités dérivées : le kilovolt ($\SI{1}{kV} = \SI{1000}{V}$) et le millivolt ($\SI{1}{mV} = \SI{0,001}{V}$).
\end{definition}

\begin{exemplebox}[Quelques tensions de la vie courante au Cameroun]
\begin{itemize}
\item Pile plate ou pile bâton (lampe de poche) : $\SI{1,5}{V}$ ; pile carrée (radio) : $\SI{4,5}{V}$ ou $\SI{9}{V}$.
\item Batterie de moto-taxi (démarrage) : $\SI{12}{V}$.
\item Prise de courant domestique ENEO : $\SI{230}{V}$ (valeur efficace, courant alternatif 50 Hz).
\item Lignes de transport haute tension (pylônes) : $\SI{90}{kV}$ à $\SI{225}{kV}$.
\end{itemize}
\end{exemplebox}

\courssub{L'intensité du courant électrique}

\textbf{Intuition.} Reprenons le tuyau d'eau : on peut s'intéresser au \emph{débit}, c'est-à-dire à la quantité d'eau qui passe chaque seconde. En électricité, le courant est un déplacement de porteurs de charges (les électrons dans un fil métallique). L'intensité mesure le \emph{débit} de ces charges.

\begin{definition}[Intensité du courant électrique]
L'\cle{intensité du courant électrique}, notée $I$, est la \textbf{quantité d'électricité qui traverse une section du conducteur par unité de temps}.
\begin{itemize}
\item Son unité est l'\cle{ampère}, de symbole \SI{}{A}.
\item On mesure une intensité avec un \cle{ampèremètre}, branché \textbf{en série} dans le circuit.
\end{itemize}
Unité dérivée : le milliampère ($\SI{1}{mA} = \SI{0,001}{A}$).
\end{definition}

\begin{aretenir}[Sens conventionnel du courant]
Par convention, le courant sort de la borne \textbf{positive} ($+$) du générateur, traverse le circuit extérieur et rentre par la borne \textbf{négative} ($-$). Les électrons, eux, circulent en sens inverse (borne $-$ vers borne $+$).
\end{aretenir}

\courssub{La résistance électrique et la loi d'Ohm}

\textbf{Intuition.} Un tuyau étroit ou obstrué freine le passage de l'eau. De même, certains dipôles s'opposent plus ou moins au passage du courant : c'est la \cle{résistance}.

\begin{definition}[Résistance électrique]
La \cle{résistance électrique} d'un dipôle, notée $R$, traduit sa capacité à \textbf{s'opposer au passage du courant}. Son unité est l'\cle{ohm}, de symbole $\Omega$ (lettre grecque oméga). On la mesure avec un \textbf{ohmmètre}, \textbf{hors circuit} (dipôle déconnecté et sans tension).
\end{definition}

\begin{experience}[Observation : la lampe et le résistor ohmique]
\textbf{Matériel :} un générateur de tension réglable (0--12 V), un résistor de résistance $R = \SI{6}{\Omega}$, un ampèremètre, un voltmètre, des fils de connexion.

\textbf{Protocole :}
\begin{enumerate}
\item Monter le circuit : générateur, interrupteur, ampèremètre en série, résistor. Voltmètre en dérivation aux bornes du résistor.
\item Fermer l'interrupteur. Faire varier la tension $U$ aux bornes du résistor (par exemple 3 V, 6 V, 9 V, 12 V) et relever l'intensité $I$ correspondante.
\item Ouvrir l'interrupteur entre chaque mesure pour ne pas chauffer le résistor.
\end{enumerate}

\textbf{Observations :}
\begin{center}
\begin{tabular}{|c|c|c|c|c|}
\hline
\rowcolor{popblueL} $U$ (V) & 3 & 6 & 9 & 12 \\
\hline
$I$ (A) & 0,5 & 1 & 1,5 & 2 \\
\hline
$\dfrac{U}{I}$ ($\Omega$) & 6 & 6 & 6 & 6 \\
\hline
\end{tabular}
\end{center}

\textbf{Interprétation :} le quotient $\dfrac{U}{I}$ reste constant et égal à la résistance $R = \SI{6}{\Omega}$. La tension et l'intensité sont proportionnelles. Le dipôle est \textbf{ohmique} (résistor).

\textbf{Conclusion :} pour un résistor, $U = R \times I$.
\end{experience}

\begin{propriete}[Loi d'Ohm (rappel exigible au BEPC)]
Pour un \textbf{résistor} (dipôle ohmique) traversé par un courant d'intensité $I$, la tension $U$ à ses bornes est donnée par :
\[ U = R \cdot I \]
avec $U$ en volts (\SI{}{V}), $R$ en ohms ($\Omega$) et $I$ en ampères (\SI{}{A}).
Formes équivalentes : $I = \dfrac{U}{R}$ et $R = \dfrac{U}{I}$.
\end{propriete}

\begin{exoresolu}[Calcul d'une intensité avec la loi d'Ohm]
On mesure une tension $U = \SI{12}{V}$ aux bornes d'une résistance $R = \SI{6}{\Omega}$. Calculer l'intensité $I$ du courant qui la traverse.
\tcblower
\textbf{Données :} $U = \SI{12}{V}$ ; $R = \SI{6}{\Omega}$.

\textbf{Formule :} d'après la loi d'Ohm, $I = \dfrac{U}{R}$.

\textbf{Application numérique :}
\begin{align*}
I &= \dfrac{U}{R} = \dfrac{12}{6} = \SI{2}{A}
\end{align*}
\textbf{Conclusion :} l'intensité du courant qui traverse la résistance est $I = \SI{2}{A}$.
\end{exoresolu}

\courssub{Les appareils de mesure : voltmètre, ampèremètre et wattmètre}

Chaque grandeur a son instrument, et chaque instrument a \textbf{sa manière d'être branché}. C'est un point capital, souvent demandé au BEPC.

\begin{methode}[Brancher correctement un voltmètre et un ampèremètre]
\begin{enumerate}
\item \textbf{Avant tout branchement}, ouvrir l'interrupteur (circuit hors tension) et choisir le calibre le plus grand de l'appareil.
\item \textbf{Voltmètre :} le brancher \textbf{en dérivation} (en parallèle) aux bornes du dipôle dont on veut mesurer la tension. Relier sa borne V (ou $+$) du côté de la borne $+$ du générateur et sa borne COM du côté de la borne $-$.
\item \textbf{Ampèremètre :} \textbf{ouvrir le circuit} à l'endroit où l'on veut mesurer le courant, puis insérer l'ampèremètre \textbf{en série}, borne A (ou $+$) du côté de la borne $+$ du générateur.
\item Fermer l'interrupteur, lire la valeur, puis descendre de calibre si nécessaire pour plus de précision.
\end{enumerate}
\end{methode}

\begin{attention}[Deux erreurs qui grillent les appareils]
\begin{itemize}
\item \textbf{Jamais un voltmètre en série :} sa résistance interne est très grande ; branché en série, il bloque presque tout le courant et le circuit ne fonctionne plus.
\item \textbf{Jamais un ampèremètre en dérivation :} sa résistance interne est très faible ; branché directement aux bornes d'un dipôle ou d'un générateur, il crée un \textbf{court-circuit} : le courant devient énorme et l'appareil (ou le fusible) est détruit.
\end{itemize}
Retiens : \emph{le voltmètre regarde de côté (dérivation), l'ampèremètre se met dans le passage (série)}.
\end{attention}

\begin{definition}[Le wattmètre]
Le \cle{wattmètre} mesure la \textbf{puissance électrique} $P$ reçue par un dipôle, exprimée en \cle{watts} (\SI{}{W}). Il possède \textbf{quatre bornes} : deux bornes « courant » (notées A ou $\sim$) branchées \textbf{en série} (comme un ampèremètre) et deux bornes « tension » (notées V) branchées \textbf{en dérivation} (comme un voltmètre). On dit qu'il a un \textbf{branchement combiné}.
\end{definition}

\begin{popfigure}
\begin{center}
\begin{tikzpicture}[line width=1.1pt, scale=1.05]
% --- Circuit principal ---
% Générateur
\draw (0,0) -- (0,1.6);
\draw[popink] (-0.5,1.6) -- (0.5,1.6);
\draw[popink] (-0.28,2.0) -- (0.28,2.0);
\node[left] at (-0.55,1.6) {\small $+$};
\node[left] at (-0.35,2.05) {\small $-$};
\node[left, align=right] at (-0.7,1.8) {\small Générateur};
\draw (0,2.0) -- (0,3);
% Fil haut vers ampèremètre
\draw (0,3) -- (1.6,3);
% Ampèremètre (série)
\draw[poporange, fill=poporangeL] (2.1,3) circle (0.5);
\node at (2.1,3) {\textbf{A}};
\draw (2.6,3) -- (4.2,3);
% Wattmètre : borne courant en série
\draw[poppurple, fill=poppurpleL] (4.9,3) circle (0.7);
\node at (4.9,3) {\textbf{W}};
\node[above, poppurple] at (4.9,3.85) {\small Wattmètre};
\draw (5.6,3) -- (7.6,3);
% Descente vers résistor
\draw (7.6,3) -- (7.6,2.4);
% Résistor
\draw[popblue] (7.6,2.4) rectangle (8.4,1.2);
\node[right] at (8.5,1.8) {\small Résistor $R$};
\draw (7.6,1.2) -- (7.6,0);
% Fil retour
\draw (7.6,0) -- (0,0);
% Voltmètre en dérivation aux bornes du résistor
\draw (7.6,2.4) -- (9.6,2.4) -- (9.6,2.1);
\draw[popgreen, fill=popgreenL] (9.6,1.8) circle (0.5);
\node at (9.6,1.8) {\textbf{V}};
\draw (9.6,1.3) -- (9.6,1.2) -- (7.6,1.2);
\node[right, popgreen] at (10.2,1.8) {\small Voltmètre};
% Borne tension du wattmètre en dérivation
\draw[poppurple, dashed] (4.9,3.7) -- (4.9,4.4) -- (7.6,4.4) -- (7.6,3.15);
\node[poppurple] at (6.2,4.65) {\small bornes tension (dérivation)};
% Sens du courant
\draw[-{Stealth}, popink, line width=1.3pt] (0.5,3.35) -- (1.3,3.35);
\node[above, popink] at (0.9,3.4) {\small $I$};
\end{tikzpicture}
\end{center}
\legende{Circuit de mesure : l'ampèremètre (A, orange) est branché \textbf{en série} ; le voltmètre (V, vert) est branché \textbf{en dérivation} aux bornes du résistor ; le wattmètre (W, violet) a un branchement \textbf{combiné} : ses bornes courant en série, ses bornes tension (pointillés) en dérivation.}
\end{popfigure}

\begin{aretenir}[Symboles normalisés des appareils de mesure]
Sur un schéma de circuit, un appareil de mesure se représente par un \textbf{cercle} contenant la lettre de la grandeur ou de l'unité :
\begin{itemize}
\item \textbf{V} dans un cercle : voltmètre (mesure $U$ en volts) ;
\item \textbf{A} dans un cercle : ampèremètre (mesure $I$ en ampères) ;
\item \textbf{W} dans un cercle : wattmètre (mesure $P$ en watts) ;
\item $\Omega$ dans un cercle : ohmmètre (mesure $R$ en ohms, hors circuit).
\end{itemize}
\end{aretenir}

\courssub{Le compteur kilowattheure : l'appareil de la facture}

Revenons à la famille Mbarga. Le voltmètre, l'ampèremètre et le wattmètre donnent des mesures \emph{instantanées}. Mais comment ENEO sait-il combien la famille doit payer ? Grâce au \cle{compteur électrique}, fixé à l'entrée de la maison.

\begin{definition}[Compteur kilowattheure]
Le \cle{compteur kilowattheure} mesure l'\textbf{énergie électrique} consommée par l'ensemble des appareils d'une installation. L'unité d'énergie utilisée pour la facturation est le \cle{kilowattheure}, de symbole \SI{}{kWh} :
\[ \SI{1}{kWh} = \SI{1000}{Wh} \]
C'est l'énergie consommée par un appareil de puissance $\SI{1000}{W}$ fonctionnant pendant une heure. La facture ENEO est calculée en multipliant le nombre de kilowattheures consommés par le prix du kilowattheure (tarif en vigueur).
\end{definition}

\begin{exemplebox}[Lecture d'un compteur]
Le compteur de la famille Mbarga indiquait \num{12450} kWh le 1\textsuperscript{er} du mois et \num{12750} kWh le 30. L'énergie consommée dans le mois est :
\[ E = 12750 - 12450 = \SI{300}{kWh} \]
Si le kilowattheure coûte 100 FCFA (tarif indicatif), la facture d'énergie s'élève à $300 \times 100 = \num{30000}$ FCFA (hors taxes et redevances fixes).
\end{exemplebox}

\begin{savaistu}[Et chez toi ?]
De nombreux compteurs modernes au Cameroun sont des compteurs \emph{prépayés} (type Cash Power) : on achète du crédit en kilowattheures (via Mobile Money ou agence) et le compteur coupe l'alimentation quand le crédit est épuisé. Surveiller son compteur, c'est déjà faire de la physique : chaque appareil laissé allumé inutilement fait tourner le compteur... et gonfler la facture ! C'est exactement le mystère que les enfants Mbarga vont résoudre dans la suite de la leçon, en calculant la puissance de chaque appareil.
\end{savaistu}

\begin{aretenir}[Bilan de la section : Grandeurs, unités, appareils, branchements]
\begin{center}
\begin{tabularx}{\linewidth}{|l|X|l|X|}
\hline
\rowcolor{popblueL} \textbf{Grandeur} & \textbf{Symbole et unité} & \textbf{Appareil} & \textbf{Branchement} \\
\hline
Tension & $U$ en volts (\SI{}{V}) & Voltmètre & En dérivation \\
\hline
Intensité & $I$ en ampères (\SI{}{A}) & Ampèremètre & En série \\
\hline
Résistance & $R$ en ohms ($\Omega$) & Ohmmètre & Hors circuit \\
\hline
Puissance & $P$ en watts (\SI{}{W}) & Wattmètre & Combiné (série + dérivation) \\
\hline
Énergie & $E$ en kilowattheures (\SI{}{kWh}) & Compteur & À l'entrée de l'installation \\
\hline
\end{tabularx}
\end{center}
Loi d'Ohm pour un résistor : $U = R \cdot I$.
\end{aretenir}

\courssub{Puissance électrique d'un résistor}

Les enfants Mbarga ont mesuré $U$ et $I$ pour chaque appareil. Maintenant, ils veulent savoir \textbf{lequel consomme le plus}. C'est la notion de \cle{puissance électrique}.

\begin{definition}[Puissance électrique]
La \cle{puissance électrique} $P$ reçue par un dipôle est l'\textbf{énergie électrique transférée par unité de temps}. Elle s'exprime en \cle{watts} (\SI{}{W}).
\[ P = U \cdot I \]
avec $P$ en watts (\SI{}{W}), $U$ en volts (\SI{}{V}) et $I$ en ampères (\SI{}{A}).
\end{definition}

\begin{propriete}[Expressions de la puissance pour un résistor (exigibles au BEPC)]
Pour un résistor de résistance $R$ traversé par un courant $I$ sous une tension $U$, on a :
\begin{align*}
P &= U \cdot I \quad \text{(définition)} \\
P &= R \cdot I^2 \quad \text{(en remplaçant $U = R I$)} \\
P &= \frac{U^2}{R} \quad \text{(en remplaçant $I = \frac{U}{R}$)}
\end{align*}
\textbf{Attention :} ces deux dernières formules ne sont valables \textbf{que pour un résistor} (dipôle ohmique). La formule $P = U I$ est universelle pour tout dipôle en régime continu.
\end{propriete}

\begin{exemplebox}[Puissance et effet Joule]
Quand un courant traverse un résistor, l'énergie électrique est transformée en \textbf{énergie thermique} (chaleur) : c'est l'\cle{effet Joule}. C'est le principe du fer à repasser, du chauffe-eau, du four, de la bouilloire... Mais c'est aussi une perte de'énergie dans les fils de transport (échauffement des câbles ENEO).
\end{exemplebox}

\begin{exoresolu}[Calcul de puissance : le fer à repasser]
Un fer à repasser est traversé par un courant d'intensité $I = \SI{5}{A}$ sous une tension $U = \SI{230}{V}$. Sa résistance mesurée à froid est $R = \SI{46}{\Omega}$.
\begin{enumerate}
\item Calculer sa puissance absorbée $P$.
\item Vérifier le résultat avec la formule $P = R I^2$.
\end{enumerate}
\tcblower
\textbf{Données :} $U = \SI{230}{V}$ ; $I = \SI{5}{A}$ ; $R = \SI{46}{\Omega}$.

\textbf{1. Formule $P = U I$ :}
\begin{align*}
P &= U \times I = 230 \times 5 = \SI{1150}{W}
\end{align*}

\textbf{2. Vérification avec $P = R I^2$ :}
\begin{align*}
P &= R \times I^2 = 46 \times 5^2 = 46 \times 25 = \SI{1150}{W}
\end{align*}
\textbf{Conclusion :} La puissance du fer est $P = \SI{1150}{W}$ (soit $\SI{1,15}{kW}$). Les deux formules donnent le même résultat car le fer se comporte comme un résistor (à température de fonctionnement).
\end{exoresolu}

\courssub{Mesure de la puissance et de l'énergie consommée}

\courssub{Mesure directe de la puissance : le wattmètre}
Le wattmètre donne directement la puissance $P$ en watts. C'est l'instrument le plus simple pour une mesure instantanée. Mais on n'a pas toujours un wattmètre sous la main !

\courssub{Mesure indirecte de la puissance : méthode voltmètre-ampèremètre}
On mesure $U$ aux bornes du dipôle (voltmètre en dérivation) et $I$ qui le traverse (ampèremètre en série), puis on calcule $P = U \times I$. Cette méthode fonctionne pour \textbf{tout type de dipôle} (résistor, moteur, lampe à DEL, etc.).

\begin{methode}[Mesurer la puissance d'un appareil sans wattmètre]
\begin{enumerate}
\item Monter le circuit : générateur, interrupteur, ampèremètre (A) en série, dipôle (D).
\item Brancher le voltmètre (V) en dérivation aux bornes du dipôle D.
\item Fermer l'interrupteur, relever $U$ et $I$.
\item Calculer $P = U \times I$.
\end{enumerate}
\end{methode}

\courssub{Énergie électrique : relation fondamentale}
La puissance est un débit d'énergie. L'énergie $E$ consommée pendant une durée $t$ est :
\[ E = P \times t \]
\begin{itemize}
\item Si $P$ est en watts (\SI{}{W}) et $t$ en secondes (\SI{}{s}), $E$ est en \textbf{joules} (\SI{}{J}).
\item Si $P$ est en kilowatts (\SI{}{kW}) et $t$ en heures (\SI{}{h}), $E$ est en \textbf{kilowattheures} (\SI{}{kWh}).
\end{itemize}

\begin{propriete}[Conversion énergie : Joule $\leftrightarrow$ kilowattheure]
\[ \SI{1}{kWh} = \SI{1000}{W} \times \SI{3600}{s} = \SI{3,6 \times 10^6}{J} = \SI{3,6}{MJ} \]
\end{propriete}

\begin{exoresolu}[Énergie et facture : la lampe du salon]
Une lampe halogène de puissance $P = \SI{500}{W}$ fonctionne pendant $t = \SI{4}{h}$ chaque soir pendant 30 jours.
\begin{enumerate}
\item Calculer l'énergie consommée en kWh sur le mois.
\item Calculer le coût si le kWh est facturé 100 FCFA.
\end{enumerate}
\tcblower
\textbf{Données :} $P = \SI{500}{W} = \SI{0,5}{kW}$ ; $t_{\text{jour}} = \SI{4}{h}$ ; $N = 30$ jours ; Prix = 100 FCFA/kWh.

\textbf{1. Durée totale :} $t = 4 \times 30 = \SI{120}{h}$.
\textbf{Énergie :} $E = P \times t = 0,5 \times 120 = \SI{60}{kWh}$.

\textbf{2. Coût :} $\text{Coût} = 60 \times 100 = \num{6000}$ FCFA.

\textbf{Conclusion :} Cette seule lampe coûte 6000 FCFA par mois. Si on la remplace par une LED de 50 W (10 fois moins), le coût tombe à 600 FCFA ! L'enquête des enfants Mbarga porte ses fruits.
\end{exoresolu}

\courssub{Sécurité électrique et consignes de prévention}

L'électricité tue. Au Cameroun, les accidents domestiques (électrocution, incendies) sont fréquents : fils dénudés, rallonges surchargées, branchements sauvages (« \emph{by-pass} » sur le compteur), utilisation d'appareils défectueux près de l'eau.

\begin{aretenir}[Les 3 dangers principaux de l'électricité domestique]
\begin{enumerate}
\item \textbf{L'électrocution :} passage du courant dans le corps humain (seuil de dangerosité : $\SI{30}{mA}$ en courant alternatif 50 Hz). Causes : fil nu, appareil non mis à la terre, mains mouillées.
\item \textbf{L'incendie par échauffement (effet Joule) :} surcharge sur une rallonge ou une multiprise (trop d'appareils branchés dessus), section de fil trop faible pour l'intensité demandée, contact lâche dans une prise.
\item \textbf{L'arc électrique / court-circuit :} contact direct entre phase et neutre (fil coupé, isolation fondue). Intensité énorme, chaleur extrême, projection de métal fondu.
\end{enumerate}
\end{aretenir}

\begin{propriete}[Les protections obligatoires dans une installation (norme NF C 15-100 / ENEO)]
\begin{itemize}
\item \textbf{Le disjoncteur différentiel (interrupteur différentiel) :} coupe le courant si une fuite vers la terre est détectée ($> \SI{30}{mA}$). \textbf{Protection des personnes.} Indispensable sur chaque circuit (prises, éclairage, gros appareils).
\item \textbf{Le disjoncteur divisionnaire (ou fusible) :} coupe le courant si l'intensité dépasse la valeur nominale du fil (ex: \SI{16}{A} pour prises, \SI{10}{A} pour éclairage, \SI{32}{A} pour four/chaudière). \textbf{Protection des fils (incendie).}
\item \textbf{La mise à la terre (prise de terre) :} relie les parties métalliques des appareils (carcasse machine à laver, réfrigérateur) à la terre. En cas de fuite, le courant part vers la terre et déclenche le différentiel.
\end{itemize}
\end{propriete}

\begin{attention}[Erreurs fréquentes et pratiques dangereuses au Cameroun]
\begin{itemize}
\item \textbf{Le « by-pass » (raccordement frauduleux avant le compteur) :} \emph{Illégal, dangereux (pas de protection différentielle), puni par la loi.} C'est la cause majeure d'électrocutions mortelles.
\item \textbf{Rallonge / multiprise en « guirlande » :} brancher une multiprise sur une autre multiprise. Risque de surchauffe du fil de la première multiprise (section souvent \SI{1,5}{mm^2} pour \SI{16}{A} max).
\item \textbf{Fil de section insuffisante :} utiliser du fil d'éclairage (\SI{1,5}{mm^2}) pour un fer à repasser ou un climatiseur. Le fil fond, l'isolant brûle.
\item \textbf{Manipuler l'électricité pieds nus ou mains mouillées :} la résistance du corps chute drastiquement (de $\sim \SI{100}{k\Omega}$ sec à $\sim \SI{1}{k\Omega}$ mouillé), le courant devient mortel.
\item \textbf{Remplacer un fusible par un fil de cuivre ou un trombone :} \textbf{Interdit !} Plus de protection contre la surcharge = incendie garanti.
\end{itemize}
\end{attention}

\begin{methode}[Réagir face à un accident électrique]
\begin{enumerate}
\item \textbf{NE PAS TOUCHER LA VICTIME} tant qu'elle est en contact avec le courant (risque de second accident).
\item \textbf{COUPER LE COURANT} immédiatement (disjoncteur général, débrancher l'appareil si prise accessible).
\item \textbf{APPELER LES SECOURS} (SAMU \num{112} ou \num{119}, Pompiers \num{118}) en donnant l'adresse précise.
\item Si formation secourisme : commencer la RCP (Réanimation Cardio-Pulmonaire) si la victime ne respire plus.
\end{enumerate}
\end{methode}

\courssub{Applications concrètes et résolution de problèmes}

Voici des exercices types BEPC pour t'entraîner. Résous-les avant de regarder les corrigés.

\begin{exercice}[Entraînement BEPC - Puissance et Énergie]
\begin{enumerate}
\item \textbf{Choix multiple :} Un résistor de résistance $R = \SI{10}{\Omega}$ est traversé par un courant $I = \SI{0,5}{A}$. Sa puissance dissipée est :
\begin{itemize}
\item A. $\SI{2,5}{W}$ \item B. $\SI{5}{W}$ \item C. $\SI{25}{W}$ \item D. $\SI{50}{W}$
\end{itemize}
\item \textbf{Calcul direct :} Un moteur de pompe à eau (forage) fonctionne sous $U = \SI{230}{V}$ et absorbe une intensité $I = \SI{4,5}{A}$.
\begin{enumerate}
\item Calculer sa puissance absorbée $P$ en watts.
\item Quelle énergie en kWh consomme-t-il s'il fonctionne 2 heures par jour pendant 30 jours ?
\item Coût mensuel si kWh = 100 FCFA ?
\end{enumerate}
\item \textbf{Inverse :} Une bouilloire de puissance $P = \SI{2000}{W}$ est branchée sur le secteur $\SI{230}{V}$.
\begin{enumerate}
\item Calculer l'intensité $I$ qu'elle demande au réseau.
\item Quelle section de fil minimum faut-il prévoir (rappel : \SI{1,5}{mm^2} $\to$ \SI{16}{A} ; \SI{2,5}{mm^2} $\to$ \SI{25}{A}) ?
\end{enumerate}
\item \textbf{Sécurité :} On veut brancher sur une même multiprise (max \SI{16}{A} / \SI{3680}{W}) : un réfrigérateur (\SI{150}{W}), un micro-ondes (\SI{1200}{W}) et une bouilloire (\SI{2000}{W}). Est-ce raisonnable ? Justifier par un calcul.
\end{enumerate}
\end{exercice}

\begin{corrige}[Exercice Entraînement BEPC]
\begin{enumerate}
\item \textbf{Réponse B.} $P = R I^2 = 10 \times (0,5)^2 = 10 \times 0,25 = \SI{2,5}{W}$ ? Non, $10 \times 0,25 = 2,5$. Attendez. $P = R I^2 = 10 \times 0,25 = \SI{2,5}{W}$. C'est la réponse A.
\textbf{Correction :} $P = R I^2 = 10 \times 0,5^2 = 10 \times 0,25 = \SI{2,5}{W}$. La bonne réponse est \textbf{A}. (Piège : ne pas confondre $U=RI=5V$ puis $P=UI=2,5W$).
\item \textbf{Pompe à eau :}
\begin{enumerate}
\item $P = U I = 230 \times 4,5 = \SI{1035}{W} \approx \SI{1,04}{kW}$.
\item $t = 2 \times 30 = \SI{60}{h}$. $E = P \times t = 1,035 \times 60 = \SI{62,1}{kWh}$.
\item Coût $= 62,1 \times 100 = \num{6210}$ FCFA.
\end{enumerate}
\item \textbf{Bouilloire :}
\begin{enumerate}
\item $I = \dfrac{P}{U} = \dfrac{2000}{230} \approx \SI{8,7}{A}$.
\item $\SI{8,7}{A} < \SI{16}{A}$. Un fil de \SI{1,5}{mm^2} (prise standard) suffit pour la bouilloire seule. Mais attention à la multiprise si d'autres appareils sont branchés dessus.
\end{enumerate}
\item \textbf{Multiprise :}
Puissance totale $P_{\text{tot}} = 150 + 1200 + 2000 = \SI{3350}{W}$.
Intensité totale $I_{\text{tot}} = \dfrac{3350}{230} \approx \SI{14,6}{A}$.
Théoriquement $14,6 A < 16 A$ (limite multiprise). \textbf{MAIS} c'est limite : au démarrage du frigo (pic de courant) ou si la tension baisse (courant monte), le disjoncteur de la multiprise va sauter, voire faire fondre les contacts. \textbf{Déconseillé.} Mettre la bouilloire sur une prise murale dédiée.
\end{enumerate}
\end{corrige}

\begin{exercice}[Problème synthèse : L'enquête des Mbarga]
Les enfants Mbarga mesurent les grandeurs suivantes sur trois appareils suspects :
\begin{center}
\begin{tabular}{|c|c|c|c|}
\hline
\rowcolor{popblueL} \textbf{Appareil} & \textbf{Tension $U$ (V)} & \textbf{Intensité $I$ (A)} & \textbf{Durée d'usage/jour} \\
\hline
Vieux Réfrigérateur & 230 & 0,8 & 24 h (en marche 1/3 du temps) \\
\hline
Lampes Halogènes (5 x 500 W) & 230 & -- & 5 h \\
\hline
Fer à repasser & 230 & 5,2 & 1 h \\
\hline
\end{tabular}
\end{center}
\begin{enumerate}
\item Compléter le tableau : calculer la puissance $P$ de chaque appareil et l'intensité des lampes (sachant qu'elles sont des résistors).
\item Calculer l'énergie journalière (en kWh) de chaque appareil. (Pour le frigo : temps de marche réel = 24 h / 3 = 8 h).
\item Calculer l'énergie mensuelle (30 j) et le coût (100 FCFA/kWh).
\item Quel est l'appareil le plus coûteux ? Quelle recommandation faites-vous à la famille ?
\end{enumerate}
\end{exercice}

\begin{corrige}[Problème synthèse : L'enquête des Mbarga]
\begin{enumerate}
\item \textbf{Calculs de puissance et intensité manquante :}
\begin{itemize}
\item \textbf{Frigorifique :} $P = U I = 230 \times 0,8 = \SI{184}{W}$.
\item \textbf{Lampes (5 x 500 W) :} $P_{\text{tot}} = 5 \times 500 = \SI{2500}{W}$. $I = \dfrac{P}{U} = \dfrac{2500}{230} \approx \SI{10,9}{A}$.
\item \textbf{Fer :} $P = U I = 230 \times 5,2 = \SI{1196}{W} \approx \SI{1,2}{kW}$.
\end{itemize}
\item \textbf{Énergie journalière :}
\begin{itemize}
\item Frigo : $E = 0,184 \times 8 = \SI{1,47}{kWh}$.
\item Lampes : $E = 2,5 \times 5 = \SI{12,5}{kWh}$.
\item Fer : $E = 1,196 \times 1 \approx \SI{1,2}{kWh}$.
\end{itemize}
\item \textbf{Énergie mensuelle et coût :}
\begin{itemize}
\item Frigo : $1,47 \times 30 = \SI{44,1}{kWh}$ $\to$ \num{4410} FCFA.
\item Lampes : $12,5 \times 30 = \SI{375}{kWh}$ $\to$ \num{37500} FCFA.
\item Fer : $1,2 \times 30 = \SI{36}{kWh}$ $\to$ \num{3600} FCFA.
\end{itemize}
\item \textbf{Conclusion :} Ce sont les \textbf{lampes halogènes} qui coûtent le plus cher (\num{37500} FCFA/mois, soit 85\% de la facture de ces 3 appareils !).
\textbf{Recommandation :} Remplacer les 5 halogènes (2500 W) par 5 lampes LED de 50 W chacune (puissance totale 250 W, 10 fois moins). Économie mensuelle : $\approx \num{33750}$ FCFA. L'investissement (ampoules LED) est rentabilisé en moins d'un mois.
\end{enumerate}
\end{corrige}

\begin{aretenir}[Synthèse de la leçon : Ce qu'il faut retenir pour le BEPC]
\begin{itemize}
\item \textbf{Grandeurs :} $U$ (V, voltmètre //), $I$ (A, ampèremètre --), $R$ ($\Omega$, ohmmètre hors circuit), $P$ (W, wattmètre combiné), $E$ (kWh, compteur).
\item \textbf{Loi d'Ohm (résistor) :} $U = R I$.
\item \textbf{Puissance (universel) :} $P = U I$.
\item \textbf{Puissance (résistor seulement) :} $P = R I^2 = \dfrac{U^2}{R}$.
\item \textbf{Énergie :} $E = P \times t$ (Joules si $P$ en W, $t$ en s ; kWh si $P$ en kW, $t$ en h).
\item \textbf{Conversion :} $\SI{1}{kWh} = \SI{3,6}{MJ}$.
\item \textbf{Sécurité :} Différentiel 30 mA (protection personnes), Disjoncteur/Fusible (protection fils), Mise à la terre, Pas de surcharge multiprises, Pas de by-pass, Mains sèches.
\end{itemize}
\end{aretenir}

\coursec{Puissance électrique d'un résistor}

Dans la section précédente, tu as appris à mesurer les grandeurs électriques : la tension $U$ avec un voltmètre, l'intensité $I$ avec un ampèremètre. Mais à quoi servent ces mesures dans la vie quotidienne ? Quand ta mère branche le fer à repasser, le chauffe-eau ou la bouilloire, elle sait que certains appareils « consomment beaucoup » et font grimper la facture d'ENEO, tandis qu'une petite ampoule LED éclaire longtemps pour presque rien. Ce qui distingue ces appareils, c'est leur \cle{puissance électrique}. C'est cette grandeur que nous allons définir, calculer et mesurer dans cette section.

\courssub{Notion de puissance électrique : l'intuition d'abord}

\textbf{Une situation de la vie courante.}\par
Imagine deux appareils branchés sur la même prise de \SI{220}{V} : une ampoule LED de \SI{10}{W} et un radiateur électrique de \SI{2000}{W}. Au bout d'une heure de fonctionnement, le radiateur a chauffé toute la pièce et le compteur a beaucoup tourné ; l'ampoule, elle, a à peine fait bouger le compteur. Pourtant, les deux appareils reçoivent la même tension ! La différence vient de la \emph{vitesse} à laquelle chacun transforme l'énergie électrique en une autre forme d'énergie (lumière, chaleur).

\begin{experience}[Observer l'effet de la puissance]
\textbf{Matériel :} un générateur de tension réglable, deux lampes de puissances nominales différentes (par exemple \SI{3}{W} et \SI{10}{W}), des fils de connexion, un ampèremètre et un voltmètre.\par
\textbf{Protocole :}
\begin{enumerate}
\item Brancher la lampe de \SI{3}{W} seule aux bornes du générateur, l'ampèremètre en série et le voltmètre en dérivation aux bornes de la lampe. Mesurer $U$ et $I$.
\item Remplacer par la lampe de \SI{10}{W}, sous la même tension. Mesurer à nouveau $U$ et $I$.
\item Comparer les éclats des deux lampes.
\end{enumerate}
\textbf{Observations :} sous la même tension, la lampe de plus forte puissance est traversée par une intensité plus grande et brille davantage.\par
\textbf{Interprétation :} la puissance dépend à la fois de la tension $U$ et de l'intensité $I$. Plus le produit $U \times I$ est grand, plus l'appareil transforme rapidement l'énergie électrique.
\end{experience}

\begin{definition}[Puissance électrique]
La \cle{puissance électrique} $P$ d'un appareil est l'énergie électrique qu'il transfère (reçoit ou transforme) \textbf{par unité de temps}. C'est un taux de transfert d'énergie.\par
Son unité dans le Système international est le \cle{watt} (symbole : W), en hommage à l'ingénieur écossais James Watt.\par
Une puissance de \SI{1}{W} correspond à un transfert d'énergie de \SI{1}{J} par seconde : $1~\text{W} = 1~\text{J/s}$.\par
Multiples usuels : le kilowatt ($1~\text{kW} = 1000~\text{W}$) et le mégawatt ($1~\text{MW} = 10^{6}~\text{W}$, utilisé pour les centrales comme celle de Song-Loulou).
\end{definition}

\courssub{L'expression fondamentale : $P = U \cdot I$}

L'expérience précédente nous a montré que la puissance dépend du couple $(U, I)$. La physique donne le résultat exact :

\begin{propriete}[Puissance électrique d'un dipôle]
La puissance électrique $P$ reçue par un dipôle soumis à une tension $U$ et traversé par un courant d'intensité $I$ est :
\[ P = U \cdot I \]
avec $P$ en watts (W), $U$ en volts (V) et $I$ en ampères (A).
\end{propriete}

\textbf{Pourquoi ce produit ?} La tension $U$ représente l'énergie fournie \emph{par unité de charge électrique} (en joules par coulomb), et l'intensité $I$ représente la quantité de charge qui circule \emph{par seconde} (en coulombs par seconde). En les multipliant, on obtient bien des joules par seconde, c'est-à-dire des watts :
\[ \frac{\text{J}}{\text{C}} \times \frac{\text{C}}{\text{s}} = \frac{\text{J}}{\text{s}} = \text{W}. \]

\begin{exemplebox}[Quelques puissances nominales d'appareils courants]
La \cle{puissance nominale} est la puissance pour laquelle l'appareil est conçu, indiquée par le fabricant sur la plaque signalétique.
\begin{center}
\begin{tabularx}{\linewidth}{|l|X|}\hline
\rowcolor{popblueL} \textbf{Appareil} & \textbf{Puissance nominale} \\ \hline
Ampoule LED & \SI{10}{W} \\ \hline
Téléviseur & \SI{100}{W} \\ \hline
Fer à repasser & \SI{1000}{W} \\ \hline
Radiateur électrique & \SI{2000}{W} \\ \hline
Groupe électrogène (petit modèle) & environ \SI{3000}{W} \\ \hline
\end{tabularx}
\end{center}
\end{exemplebox}

\courssub{Cas du résistor : les trois formes de la puissance}

Un \cle{résistor} (conducteur ohmique) est un dipôle qui obéit à la loi d'Ohm : $U = R \cdot I$. Toute l'énergie électrique qu'il reçoit est transformée en chaleur : c'est l'\cle{effet Joule}, mis à profit dans les fers à repasser, les chauffe-eau et les fours.

En combinant $P = U \cdot I$ avec la loi d'Ohm, on obtient deux autres expressions très utiles.

\textbf{Première forme : en fonction de $R$ et $I$.}\par
Remplaçons $U$ par $R \cdot I$ dans $P = U \cdot I$ :
\begin{align*}
P &= U \cdot I \\
  &= (R \cdot I) \cdot I \\
  &= R \cdot I^2
\end{align*}

\textbf{Deuxième forme : en fonction de $U$ et $R$.}\par
De la loi d'Ohm, on tire $I = \dfrac{U}{R}$. Remplaçons dans $P = U \cdot I$ :
\begin{align*}
P &= U \cdot I \\
  &= U \cdot \frac{U}{R} \\
  &= \frac{U^2}{R}
\end{align*}

\begin{propriete}[Puissance dissipée par un résistor]
Pour un résistor de résistance $R$, soumis à une tension $U$ et traversé par un courant d'intensité $I$, la puissance électrique convertie en chaleur (effet Joule) s'écrit sous trois formes équivalentes :
\[ P = U \cdot I = R \cdot I^2 = \frac{U^2}{R} \]
avec $P$ en W, $U$ en V, $I$ en A et $R$ en $\Omega$.\par
Ces formules ne sont valables que pour un dipôle ohmique (résistor). La dérivation à partir de la loi d'Ohm est exigible au BEPC : sache la refaire.
\end{propriete}

\begin{methode}[Calculer la puissance connaissant deux grandeurs parmi $U$, $I$, $R$]
\begin{enumerate}
\item \textbf{Repérer} les deux grandeurs connues dans l'énoncé et noter leurs valeurs en unités du Système international (V, A, $\Omega$).
\item \textbf{Choisir} la formule qui ne fait intervenir que ces deux grandeurs :
\begin{itemize}
\item $U$ et $I$ connus : $P = U \cdot I$ ;
\item $R$ et $I$ connus : $P = R \cdot I^2$ ;
\item $U$ et $R$ connus : $P = \dfrac{U^2}{R}$.
\end{itemize}
\item \textbf{Calculer} en remplaçant les lettres par les valeurs, sans oublier les parenthèses pour les carrés.
\item \textbf{Conclure} par une phrase avec l'unité : la puissance s'exprime en watts (W).
\item \textbf{Vérifier} (facultatif mais recommandé) : retrouver la troisième grandeur avec la loi d'Ohm et recalculer $P$ par une autre forme ; on doit trouver le même résultat.
\end{enumerate}
\end{methode}

\begin{exoresolu}[Lampe de \SI{60}{W} branchée sur le secteur \SI{220}{V}]
Énoncé. Une lampe à incandescence porte les indications « \SI{220}{V} ; \SI{60}{W} ». Elle est branchée sur une prise ENEO délivrant une tension de \SI{220}{V}.\par
1. Calculer l'intensité $I$ du courant qui la traverse.\par
2. En déduire la résistance $R$ du filament de la lampe en fonctionnement.
\tcblower
Solution détaillée.\par
\textbf{1. Calcul de l'intensité.} On connaît $P$ et $U$. On utilise $P = U \cdot I$, donc :
\[ I = \frac{P}{U} = \frac{60}{220} \approx \SI{0,27}{A}. \]
\textbf{2. Calcul de la résistance.} On applique la loi d'Ohm :
\[ R = \frac{U}{I} = \frac{220}{0,27} \approx \SI{815}{\Omega}. \]
\textbf{Vérification} avec la forme $P = \dfrac{U^2}{R}$ :
\[ P = \frac{220^2}{815} = \frac{48\,400}{815} \approx \SI{59,4}{W} \approx \SI{60}{W}. \]
Le petit écart vient de l'arrondi de $I$. Le résultat est cohérent.\par
\textbf{Conclusion :} la lampe est traversée par un courant d'environ \SI{0,27}{A} et son filament a une résistance d'environ \SI{815}{\Omega} en fonctionnement.
\end{exoresolu}

\courssub{Interprétation graphique : la puissance croît comme le carré de l'intensité}

Pour un résistor de résistance $R$ fixée, la relation $P = R \cdot I^2$ montre que la puissance est \textbf{proportionnelle au carré} de l'intensité. Conséquence spectaculaire : si l'intensité double, la puissance est multipliée par quatre ; si elle triple, la puissance est multipliée par neuf. C'est pourquoi une surintensité, même brève, peut rapidement échauffer et détériorer un appareil ou un conducteur.

\begin{popfigure}
\begin{center}
\begin{tikzpicture}
\begin{axis}[
    width=0.85\linewidth,
    height=6.5cm,
    xlabel={Intensité $I$ (A)},
    ylabel={Puissance $P$ (W)},
    xmin=0, xmax=3.2,
    ymin=0, ymax=100,
    grid=both,
    grid style={popline!40},
    xtick={0,0.5,1,1.5,2,2.5,3},
    ytick={0,20,40,60,80,100},
    axis line style={popdark},
    label style={popdark},
]
\addplot[popcurve,domain=0:3.1,samples=200]{10*x^2};
\addplot[only marks,mark=*,mark size=1.8pt,color=poporange] coordinates {(1,10) (2,40) (3,90)};
\node[color=popdark,anchor=west] at (axis cs:1.6,75) {$P = R \cdot I^2$ avec $R = 10~\Omega$};
\end{axis}
\end{tikzpicture}
\end{center}
\legende{Courbe de la puissance $P$ dissipée par un résistor de \SI{10}{\Omega} en fonction de l'intensité $I$ : c'est une branche de parabole. Les points marqués correspondent à $I = \SI{1}{A}$ ($P = \SI{10}{W}$), $I = \SI{2}{A}$ ($P = \SI{40}{W}$) et $I = \SI{3}{A}$ ($P = \SI{90}{W}$) : quand $I$ est multipliée par 3, $P$ est multipliée par 9.}
\end{popfigure}

\courssub{De la puissance à l'énergie : $E = P \cdot t$}

La puissance indique la \emph{vitesse} du transfert d'énergie ; l'\cle{énergie} est la \emph{quantité totale} transférée pendant une durée donnée. Le lien entre les deux est simple :

\begin{propriete}[Relation entre énergie et puissance]
L'énergie électrique $E$ consommée par un appareil de puissance $P$ fonctionnant pendant une durée $t$ est :
\[ E = P \cdot t \]
Si $P$ est en watts (W) et $t$ en secondes (s), alors $E$ est en \cle{joules} (J).
\end{propriete}

Le joule étant une très petite unité à l'échelle domestique, les fournisseurs d'électricité comme ENEO utilisent le \cle{kilowattheure} (kWh) : c'est l'énergie consommée par un appareil de \SI{1}{kW} fonctionnant pendant \SI{1}{h}.

\begin{aretenir}[Conversion joule et kilowattheure]
\[ 1~\text{kWh} = 1~\text{kW} \times 1~\text{h} = 1000~\text{W} \times 3600~\text{s} = 3,6 \times 10^{6}~\text{J} \]
Pour utiliser directement les kWh : $E$ (kWh) $= P$ (kW) $\times\, t$ (h).\par
Exemple : un radiateur de \SI{2000}{W} $=$ \SI{2}{kW} fonctionnant \SI{3}{h} consomme $E = 2 \times 3 = \SI{6}{kWh}$, soit $6 \times 3,6 \times 10^{6} = 2,16 \times 10^{7}~\text{J}$.
\end{aretenir}

\begin{attention}[Ne pas confondre puissance et énergie]
C'est l'erreur la plus fréquente au BEPC !
\begin{itemize}
\item La \textbf{puissance} $P$ se mesure en \textbf{watts} (W) : c'est une vitesse de transfert, une caractéristique de l'appareil à un instant donné.
\item L'\textbf{énergie} $E$ se mesure en \textbf{joules} (J) ou en \textbf{kilowattheures} (kWh) : c'est une quantité accumulée pendant une durée.
\end{itemize}
Dire « cette lampe consomme \SI{60}{W} par heure » est faux : elle a une puissance de \SI{60}{W}, et elle consomme une énergie de $60 \times 3600 = \SI{216000}{J}$ (soit \SI{0,06}{kWh}) en une heure. Autre piège : dans $E = P \cdot t$, vérifie toujours la cohérence des unités (W avec s pour obtenir des J ; kW avec h pour obtenir des kWh).
\end{attention}

\begin{savaistu}[Le compteur ENEO mesure l'énergie, pas la puissance]
Le compteur électrique installé dans les maisons camerounaises n'affiche pas la puissance instantanée : il totalise l'\textbf{énergie} consommée en kilowattheures (kWh). C'est cette énergie cumulée qui est facturée. Sur les anciens compteurs à disque, le disque tourne d'autant plus vite que la puissance totale des appareils branchés est grande : quand tu allumes le fer à repasser et la bouilloire en même temps, le disque s'emballe ! Les compteurs modernes à prépaiement décomptent les kWh du crédit acheté. Comprendre la différence entre puissance et énergie, c'est comprendre sa facture d'électricité.
\end{savaistu}

\begin{exercice}[À toi de jouer]
Un grille-pain de résistance $R = \SI{44}{\Omega}$ est branché sous une tension $U = \SI{220}{V}$.\par
1. Calculer l'intensité $I$ du courant qui le traverse.\par
2. Calculer sa puissance électrique $P$ de deux façons différentes.\par
3. Calculer, en joules puis en kWh, l'énergie consommée s'il fonctionne \SI{5}{minutes} chaque matin pendant 30 jours.
\end{exercice}

\begin{corrige}[Exercice : le grille-pain]
\textbf{1.} Loi d'Ohm : $I = \dfrac{U}{R} = \dfrac{220}{44} = \SI{5}{A}$.\par
\textbf{2.} Première façon : $P = U \cdot I = 220 \times 5 = \SI{1100}{W}$. Deuxième façon : $P = R \cdot I^2 = 44 \times 5^2 = 44 \times 25 = \SI{1100}{W}$. Les deux résultats concordent.\par
\textbf{3.} Durée totale : $t = 30 \times 5 = \SI{150}{min} = \SI{9000}{s} = \SI{2,5}{h}$.\par
En joules : $E = P \cdot t = 1100 \times 9000 = 9,9 \times 10^{6}~\text{J}$.\par
En kWh : $E = 1,1~\text{kW} \times 2,5~\text{h} = \SI{2,75}{kWh}$. Vérification : $2,75 \times 3,6 \times 10^{6} = 9,9 \times 10^{6}~\text{J}$. Cohérent.
\end{corrige}

\begin{aretenir}[L'essentiel de la section]
\begin{itemize}
\item La puissance électrique est l'énergie transférée par unité de temps ; elle s'exprime en watts (W).
\item Pour tout dipôle : $P = U \cdot I$.
\item Pour un résistor, grâce à la loi d'Ohm : $P = U \cdot I = R \cdot I^2 = \dfrac{U^2}{R}$.
\item À résistance constante, $P$ varie comme le carré de $I$ : la courbe $P(I)$ est une parabole.
\item L'énergie consommée est $E = P \cdot t$, en joules (W et s) ou en kWh (kW et h), avec $1~\text{kWh} = 3,6 \times 10^{6}~\text{J}$.
\end{itemize}
\end{aretenir}

Maintenant que tu sais calculer une puissance et une énergie, la section suivante te montrera comment on les \emph{mesure} concrètement, avec le wattmètre et le compteur électrique.

\coursec{Mesure de la puissance et de l'énergie consommée}

Dans la section précédente, tu as appris à calculer la puissance électrique d'un résistor : $P = U \cdot I = R \cdot I^2$. Mais une question pratique se pose immédiatement dans chaque maison camerounaise : combien d'électricité la famille a-t-elle réellement consommée ce mois-ci, et combien va-t-elle payer à ENEO ? Pour répondre, il ne suffit pas de connaître la puissance des appareils : il faut savoir \emph{mesurer} cette puissance et \emph{compter} l'énergie consommée dans le temps. C'est exactement le rôle du wattmètre et du compteur électrique. C'est ce que nous allons découvrir maintenant.

\courssub{Le wattmètre : mesurer directement la puissance}

\textbf{Situation concrète.} Ton père achète un fer à repasser au marché. Sur la plaque signalétique, il est écrit « 1000 W ». Mais ce fer est-il vraiment conforme ? Un appareil défectueux ou contrefait peut consommer bien plus que prévu. Comment vérifier sa puissance réelle en fonctionnement ?

\begin{definition}[Wattmètre]
Le \cle{wattmètre} est un appareil de mesure qui indique directement la \cle{puissance active} $P$ (en watts) consommée par un dipôle en fonctionnement. Il combine en interne la mesure de la tension $U$ (comme un voltmètre, branché en dérivation) et celle de l'intensité $I$ (comme un ampèremètre, branché en série), puis effectue le produit $P = U \cdot I$.
\end{definition}

\begin{experience}[Mesurer la puissance d'une lampe avec un wattmètre]
\textbf{Matériel :} un générateur de tension alternative, une lampe (par exemple \SI{60}{W}), un wattmètre, des fils de connexion.

\textbf{Protocole :}
\begin{enumerate}
\item Brancher la lampe aux bornes du générateur.
\item Insérer le wattmètre : ses bornes « courant » en série avec la lampe, ses bornes « tension » en dérivation aux bornes de la lampe.
\item Allumer et lire l'indication du wattmètre.
\end{enumerate}

\textbf{Observation :} le wattmètre affiche par exemple \SI{58}{W}.

\textbf{Interprétation :} la valeur mesurée est proche de la valeur nominale (\SI{60}{W}). L'écart s'explique par les tolérances de fabrication et par la tension réelle du réseau, qui n'est pas toujours exactement \SI{220}{V}. Le wattmètre mesure la puissance \emph{réellement consommée à l'instant $t$}, et non une valeur théorique.
\end{experience}

\begin{attention}[Puissance nominale et puissance réelle]
La puissance inscrite sur un appareil (puissance \cle{nominale}) n'est garantie que si l'appareil est alimenté sous sa tension nominale. Si la tension du réseau chute (phénomène fréquent aux heures de pointe), la puissance réelle diminue : c'est pourquoi les lampes éclairent parfois plus faiblement le soir.
\end{attention}

\courssub{Le compteur électrique : compter l'énergie}

\textbf{Situation concrète.} Devant presque chaque maison connectée au réseau, on trouve un boîtier scellé : le compteur ENEO. L'agent releveur vient chaque mois noter des chiffres. Que représentent-ils ?

\begin{definition}[Énergie électrique et compteur]
L'\cle{énergie électrique} $E$ est la quantité d'électricité consommée par les appareils pendant une durée donnée. Elle se mesure avec un \cle{compteur électrique}, dont l'\cle{index} s'exprime en \cle{kilowattheures} (kWh).
\begin{itemize}
\item Dans un compteur \emph{électromécanique}, un disque métallique tourne d'autant plus vite que la puissance consommée est grande ; sa rotation fait avancer les chiffres de l'index.
\item Dans un compteur \emph{électronique}, un circuit mesure en permanence $U$ et $I$, calcule l'énergie et l'affiche sur un écran (souvent avec une diode qui clignote : plus elle clignote vite, plus la consommation est forte).
\end{itemize}
\end{definition}

\begin{popfigure}
\begin{tikzpicture}[scale=1]
% Boîtier du compteur
\draw[fill=popcream, thick, rounded corners=4pt] (0,0) rectangle (7,4.6);
\draw[fill=popblueL, thick, rounded corners=2pt] (0.4,3.6) rectangle (6.6,4.3);
\node[popdark] at (3.5,3.95) {\small \textbf{COMPTEUR ÉLECTRIQUE --- ENEO}};
% Fenêtre d'index
\draw[fill=white, thick] (0.6,2.4) rectangle (6.4,3.3);
\foreach \x/\c in {1.05/0, 1.85/2, 2.65/3, 3.45/4, 4.25/5, 5.05/7}{
  \draw[fill=popdark] (\x-0.32,2.55) rectangle (\x+0.32,3.15);
  \node[white] at (\x,2.85) {\textbf{\c}};
}
\draw[fill=poporange] (5.55,2.55) rectangle (6.15,3.15);
\node[white] at (5.85,2.85) {\textbf{8}};
\node[popdark] at (3.5,2.25) {\small kWh};
% Disque
\draw[fill=popgoldL, thick] (2.2,1.1) circle (0.75);
\draw[fill=popdark] (2.2,1.1) circle (0.08);
\draw[popdark, thick] (2.2,1.1) -- (2.75,1.45);
\draw[-{Stealth}, poporange, thick] (3.25,0.75) arc (-30:50:0.6);
% Écran électronique
\draw[fill=popgreenL, thick] (4.1,0.7) rectangle (6.3,1.6);
\node[popdark] at (5.2,1.3) {\small 023457,8 kWh};
\node[popdark] at (5.2,0.95) {\tiny affichage électronique};
% Annotations
\draw[-{Stealth}, popblue, thick] (1.05,2.5) -- (0.7,1.9) node[below, popdark, align=center, font=\scriptsize] {index : chiffres\\à relever};
\draw[-{Stealth}, popblue, thick] (2.2,0.35) -- (2.2,-0.15) node[below, popdark, align=center, font=\scriptsize] {disque tournant\\(compteur électromécanique)};
\draw[-{Stealth}, popblue, thick] (5.85,3.2) -- (6.5,3.5) node[right, popdark, align=left, font=\scriptsize] {chiffre rouge :\\dixièmes de kWh};
\end{tikzpicture}
\legende{Schéma annoté d'un compteur kilowattheure : l'index (ici \num{23457,8} kWh) augmente au fur et à mesure de la consommation. Sur les compteurs électromécaniques, le disque tourne ; sur les compteurs électroniques, l'affichage est numérique.}
\end{popfigure}

\begin{methode}[Lire un compteur kilowattheure et calculer une consommation]
\begin{enumerate}
\item \textbf{Relever l'index} au début de la période : noter tous les chiffres de gauche à droite (le chiffre rouge indique les dixièmes, on peut l'ignorer pour une facture). Exemple : $I_1 = \num{23457}$ kWh.
\item \textbf{Relever l'index} à la fin de la période (un mois plus tard). Exemple : $I_2 = \num{23612}$ kWh.
\item \textbf{Calculer la différence} : $E = I_2 - I_1 = \num{23612} - \num{23457} = \SI{155}{kWh}$.
\item \textbf{Conclure} : la famille a consommé \SI{155}{kWh} pendant le mois.
\end{enumerate}
\end{methode}

\courssub{Relation entre énergie, puissance et durée}

\textbf{Intuition.} Une ampoule de \SI{100}{W} allumée une heure consomme peu ; la même ampoule allumée toute une nuit consomme beaucoup plus. L'énergie dépend donc de deux choses : la \emph{puissance} de l'appareil et la \emph{durée} d'utilisation. C'est un produit, comme la distance parcourue est le produit de la vitesse par le temps.

\begin{propriete}[Énergie électrique consommée]
L'énergie électrique $E$ consommée par un appareil de puissance $P$ fonctionnant pendant une durée $t$ est :
\[ E = P \times t \]
Deux systèmes d'unités cohérents :
\begin{itemize}
\item $P$ en watts (W), $t$ en secondes (s) $\Rightarrow$ $E$ en \cle{joules} (J), unité officielle du système international ;
\item $P$ en kilowatts (kW), $t$ en heures (h) $\Rightarrow$ $E$ en \cle{kilowattheures} (kWh), unité pratique des factures.
\end{itemize}
Conversion exigible : $1~\text{kWh} = \SI{1000}{W} \times \SI{3600}{s} = \num{3600000}~\text{J}$, soit $1~\text{kWh} = 3,6 \cdot 10^{6}~\text{J}$.
\end{propriete}

\begin{attention}[Le piège classique : kW, kWh et minutes]
\begin{itemize}
\item Le \textbf{kW} est une unité de \emph{puissance} (un débit d'énergie) ; le \textbf{kWh} est une unité d'\emph{énergie} (une quantité). On ne dit jamais « consommer des kW » !
\item Avant d'appliquer $E = P \times t$, convertis toujours : puissance en kW et durée en heures. Si la durée est donnée en minutes, divise par 60 : $30~\text{min} = 0,5~\text{h}$. Multiplier des kW par des minutes sans conversion est l'erreur la plus fréquente au BEPC.
\item Si $P$ est en watts et $t$ en heures, le résultat est en Wh : divise alors par 1000 pour obtenir des kWh.
\end{itemize}
\end{attention}

\begin{exemplebox}[Estimer le coût d'une consommation]
Pour obtenir le montant à payer, on multiplie l'énergie consommée par le prix du kilowattheure :
\[ \text{Coût (FCFA)} = E~(\text{kWh}) \times \text{tarif}~(\text{FCFA/kWh}) \]
Exemple : pour $E = \SI{155}{kWh}$ à un tarif moyen de 100 FCFA/kWh, le coût est $155 \times 100 = \num{15500}$ FCFA.
\end{exemplebox}

\begin{exoresolu}[Ampoule allumée chaque soir]
\textbf{Énoncé.} Une ampoule LED de puissance $P = \SI{10}{W}$ éclaire la véranda d'une maison à Yaoundé pendant 5 h chaque jour, pendant 30 jours. 1) Calculer l'énergie consommée en kWh. 2) Calculer le coût, au tarif de 100 FCFA/kWh. 3) Exprimer cette énergie en joules.
\tcblower
\textbf{Solution détaillée.}

1) Convertissons d'abord la puissance en kilowatts : $P = \SI{10}{W} = \SI{0,010}{kW}$. La durée totale de fonctionnement est :
\[ t = 5~\text{h/jour} \times 30~\text{jours} = \SI{150}{h} \]
Appliquons la relation :
\[ E = P \times t = 0,010 \times 150 = \SI{1,5}{kWh} \]

2) Coût de la consommation :
\[ \text{Coût} = E \times \text{tarif} = 1,5 \times 100 = 150~\text{FCFA} \]

3) Conversion en joules :
\[ E = 1,5 \times 3,6 \cdot 10^{6} = 5,4 \cdot 10^{6}~\text{J} \]

\textbf{Conclusion.} Une petite ampoule LED coûte très peu (environ 150 FCFA par mois) : c'est pourquoi remplacer les vieilles ampoules à incandescence de \SI{100}{W} par des LED divise la facture d'éclairage par 10 !
\end{exoresolu}

\courssub{Analyser une courbe de consommation journalière}

\textbf{Observation.} ENEO enregistre la puissance totale appelée par les abonnés à chaque instant de la journée. Représentée graphiquement, cette consommation n'est pas constante : elle présente des \emph{creux} et des \emph{pics}.

\begin{popfigure}
\begin{tikzpicture}
\begin{axis}[
    width=\linewidth, height=6.5cm,
    xlabel={Heure de la journée (h)},
    ylabel={Puissance consommée (kW)},
    xmin=0, xmax=24, ymin=0, ymax=5.4,
    xtick={0,2,4,6,8,10,12,14,16,18,20,22,24},
    ytick={0,1,2,3,4,5},
    grid=both, grid style={popline!40},
    axis line style={popdark},
    label style={popdark},
    tick label style={popdark, font=\scriptsize},
]
\addplot[popcurve, smooth, thick] coordinates {
(0,1.1) (2,0.9) (4,0.8) (6,1.4) (8,2.2) (10,2.4) (12,2.6) (14,2.3) (16,2.2) (18,3.4) (19,4.3) (20,4.6) (21,4.1) (22,2.9) (24,1.4)
};
\addplot[poporange, dashed, thick] coordinates {(19,0) (19,4.3)};
\addplot[poporange, dashed, thick] coordinates {(20,0) (20,4.6)};
\node[popdark, font=\scriptsize, align=center, fill=poporangeL, rounded corners=2pt] at (axis cs:14.5,4.6) {pic du soir (19 h -- 21 h)};
\node[popdark, font=\scriptsize, align=center, fill=popblueL, rounded corners=2pt] at (axis cs:4,1.9) {creux de la nuit};
\end{axis}
\end{tikzpicture}
\legende{Courbe de consommation électrique type d'un quartier résidentiel sur 24 h : faible la nuit, moyenne la journée, et pic marqué entre 19 h et 21 h, quand les familles allument lampes, téléviseurs, ventilateurs et cuisinières en même temps.}
\end{popfigure}

\textbf{Interprétation (démarche d'investigation).}
\begin{itemize}
\item \emph{Observation :} la courbe montre un maximum net vers 20 h et un minimum vers 4 h du matin.
\item \emph{Hypothèse :} le pic correspond au moment où le plus grand nombre d'appareils fonctionnent simultanément.
\item \emph{Vérification :} entre 19 h et 21 h, les habitants rentrent, allument l'éclairage, la télévision, les ventilateurs et préparent le repas : les puissances s'ajoutent, donc $P_{\text{totale}} = P_1 + P_2 + \cdots$ est maximale.
\item \emph{Conclusion :} la consommation d'un réseau dépend de la \emph{simultanéité} d'usage. C'est pendant ces pics que le réseau est le plus sollicité, que les tensions chutent et que surviennent les délestages.
\end{itemize}

\begin{aretenir}[L'essentiel de la section]
\begin{itemize}
\item Le \cle{wattmètre} mesure directement la puissance active $P$ d'un appareil en fonctionnement.
\item Le \cle{compteur} électrique mesure l'énergie $E$ consommée ; son index s'exprime en kWh. Consommation = différence de deux index.
\item $E = P \times t$, avec $1~\text{kWh} = 3,6 \cdot 10^{6}~\text{J}$.
\item Coût = $E~(\text{kWh}) \times$ tarif (FCFA/kWh).
\item La consommation d'un quartier présente un \cle{pic le soir}, quand tous les appareils fonctionnent en même temps.
\end{itemize}
\end{aretenir}

\courssub{Pour aller plus loin : puissances active, réactive et apparente}

\begin{savaistu}[Trois puissances, et les tarifs progressifs au Cameroun]
En courant alternatif, on distingue en réalité trois puissances (tu les étudieras en détail au lycée et pour le Baccalauréat) :
\begin{itemize}
\item la puissance \cle{active} $P$ (en W), seule réellement transformée en chaleur, lumière ou mouvement : c'est elle que paie le consommateur ;
\item la puissance \cle{réactive} $Q$ (en var), échangée par les bobines et condensateurs des moteurs et transformateurs sans être consommée ;
\item la puissance \cle{apparente} $S$ (en VA), liée aux deux précédentes par $S^2 = P^2 + Q^2$ : c'est elle qui dimensionne les câbles et les transformateurs d'ENEO.
\end{itemize}
Au Cameroun, le prix du kilowattheure n'est pas fixe : il est \emph{progressif}. La consommation mensuelle est facturée par tranches : les premiers kilowattheures (usage domestique essentiel) sont à tarif réduit, mais au-delà d'environ \SI{200}{kWh} par mois, le prix du kWh augmente fortement. Conséquence pratique : gaspiller l'électricité coûte de plus en plus cher ! Éteindre les appareils inutiles et utiliser des lampes LED permet de rester dans les tranches bon marché.
\end{savaistu}

\begin{exercice}[Facture de fin de mois]
L'index du compteur d'une famille de Douala indique \num{41230} kWh le 1\textsuperscript{er} du mois et \num{41475} kWh le 30. 1) Calculer l'énergie consommée pendant le mois. 2) Le tarif moyen étant de 100 FCFA/kWh, calculer le montant de la facture. 3) Cette famille dépasse-t-elle le seuil des \SI{200}{kWh} ? Que risque-t-elle ?
\end{exercice}

\begin{corrige}[Exercice : Facture de fin de mois]
1) $E = I_2 - I_1 = \num{41475} - \num{41230} = \SI{245}{kWh}$.

2) Coût $= 245 \times 100 = \num{24500}$ FCFA.

3) Oui : $245 > 200$, la famille dépasse le seuil de \SI{200}{kWh/mois}. Une partie de sa consommation sera donc facturée dans la tranche supérieure, à un tarif plus élevé : sa facture réelle sera en pratique supérieure au calcul simplifié de la question 2. Elle a intérêt à réduire sa consommation (LED, extinction des appareils en veille) pour rester sous le seuil.
\end{corrige}

\coursec{Sécurité électrique et consignes de prévention}

Dans la section précédente, tu as appris à mesurer la puissance et l'énergie électriques consommées par tes appareils, et tu as vu que le compteur d'ENEO enregistre chaque kilowattheure utilisé à la maison. Mais l'électricité qui fait fonctionner ta télévision, ton fer à repasser ou ton groupe électrogène n'est pas sans danger : chaque année, des accidents domestiques graves — brûlures, électrocutions, incendies — sont causés par une mauvaise installation ou un mauvais geste. Comprendre les risques et connaître les dispositifs de protection, c'est savoir protéger sa famille. C'est l'objet de cette section.

\courssub{Les principaux risques électriques}

\textbf{Une situation pour commencer.}\par
Dans une maison de quartier, on branche sur une seule rallonge un fer à repasser, une bouilloire et un réfrigérateur. Au bout de quelques minutes, le câble de la rallonge devient chaud, puis une odeur de plastique brûlé se répand. Que s'est-il passé ? Pourquoi le câble a-t-il chauffé ? C'est en répondant à ces questions que nous allons dégager les quatre grands risques électriques.

\begin{experience}[Observation : un câble qui chauffe par effet Joule]
\textbf{Matériel :} un générateur de tension continue (pile ou alimentation réglée sur faible tension, ex. \SI{12}{V}), un fil très fin (brin de laine d'acier ou fil de fer de \SI{0,1}{mm}), un interrupteur, des fils de connexion.\par
\textbf{Protocole :} on réalise un circuit simple : générateur -- interrupteur -- fil fin. On ferme l'interrupteur pendant quelques secondes seulement, sous la surveillance de l'enseignant, puis on l'ouvre immédiatement.\par
\textbf{Observations :} le fil fin devient très chaud, rougeoie, et peut même se rompre (fondre) si le courant est assez intense.\par
\textbf{Interprétation :} le passage du courant dans un conducteur produit de la chaleur : c'est l'\cle{effet Joule}. La puissance dissipée en chaleur vaut $P = R \cdot I^{2}$. Plus le courant $I$ est intense et plus la résistance $R$ du fil est grande (fil fin, long), plus l'échauffement est important.\par
\textbf{Conclusion :} un câble trop fin ou parcouru par un courant trop intense (surcharge) peut s'échauffer dangereusement, faire fondre son isolant et enflammer les matériaux voisins.
\end{experience}

On retient quatre risques principaux :

\begin{itemize}
\item L'\cle{électrocution} : le courant traverse le corps humain (contact direct avec un fil dénudé, appareil défectueux, mains mouillées). Dès quelques dizaines de milliampères, le courant perturbe le fonctionnement du cœur.
\item Le \cle{court-circuit} : contact direct entre la phase et le neutre (ou la terre) \textbf{sans passer par un appareil}. La résistance du circuit devient quasi nulle, l'intensité devient énorme ($I = U/R$), provoquant un échauffement brutal et violent des conducteurs.
\item La \cle{surchauffe} (ou surcharge) : elle survient quand trop d'appareils sont branchés sur une même ligne ou quand les câbles sont de section insuffisante ; l'effet Joule échauffe \emph{durablement} les conducteurs.
\item L'\cle{incendie} : conséquence possible d'un court-circuit ou d'une surchauffe, quand la chaleur enflamme l'isolant des câbles, les rideaux ou les meubles proches.
\end{itemize}

\begin{definition}[Court-circuit]
Un \cle{court-circuit} est un contact direct, accidentel ou non, entre la phase et le neutre (ou entre la phase et la terre) \textbf{sans passer par un récepteur} (sans charge). La résistance du circuit devient alors quasi nulle, donc l'intensité devient très grande : d'après $I = \dfrac{U}{R}$, si $R$ tend vers zéro, $I$ devient énorme. Ce courant brutal provoque un échauffement violent et quasi instantané des conducteurs.
\end{definition}

\begin{exemplebox}[Pourquoi la rallonge a-t-elle chauffé ?]
Le fer à repasser (\SI{1000}{W}), la bouilloire (\SI{1500}{W}) et le réfrigérateur (\SI{150}{W}) branchés ensemble demandent une puissance totale $P = \SI{2650}{W}$. Sous une tension $U = \SI{220}{V}$, l'intensité appelée vaut $I = \dfrac{P}{U} = \dfrac{2650}{220} \approx \SI{12,0}{A}$. Si le câble de la rallonge est prévu pour \SI{10}{A} seulement (section trop faible), il est en \textbf{surcharge} : il chauffe par effet Joule ($P = R \cdot I^{2}$), son isolant fond, et deux fils dénudés peuvent se toucher : c'est le court-circuit, puis l'incendie.
\end{exemplebox}

\courssub{Le danger du courant pour le corps humain}

Ce n'est pas la tension qui tue directement, c'est l'\textbf{intensité du courant qui traverse le corps}. Le tableau suivant donne les effets du courant alternatif (fréquence \SI{50}{Hz}) en fonction de son intensité effective traversant le corps (chemin main-pied).

\begin{center}
\begin{tabularx}{\linewidth}{|l|X|}
\hline
\rowcolor{popblueL} \textbf{Intensité traversant le corps} & \textbf{Effet sur l'organisme} \\
\hline
$I < 0{,}5\ \text{mA}$ & Aucune sensation (seuil de perception). \\
\hline
$0{,}5\ \text{mA} < I < 10\ \text{mA}$ & Picotements, sensation de fourmillement (seuil de lâcher possible). \\
\hline
$10\ \text{mA} < I < 30\ \text{mA}$ & Contraction musculaire soutenue, \textbf{impossibilité de lâcher le conducteur} (tétanisation). \\
\hline
$30\ \text{mA} < I < 50\ \text{mA}$ & Tétanisation généralisée, gêne respiratoire grave, risque de fibrillation ventriculaire. \\
\hline
$I > 50\ \text{mA}$ & \textbf{Risque élevé de fibrillation cardiaque} : le cœur bat de façon désordonnée, arrêt de la circulation sanguine, mort possible en quelques minutes sans réanimation. \\
\hline
\end{tabularx}
\end{center}

\begin{exemplebox}[Un courant de fuite de 30 mA]
Un chauffe-eau défectueux laisse « fuir » un courant de \SI{30}{mA} vers sa carcasse métallique. Une personne qui touche cette carcasse (pieds au sol) est traversée par ce courant. Or \SI{30}{mA} se situe dans la zone de tétanisation et de gêne respiratoire, proche du seuil de fibrillation (\SI{50}{mA}) : la personne est en danger de mort imminent. Si l'installation possède un disjoncteur différentiel \SI{30}{mA}, celui-ci coupe le courant en moins de \SI{40}{ms} (pour un différentiel domestique), avant que le cœur ne soit atteint. Ce dispositif sauve des vies.
\end{exemplebox}

\begin{attention}[L'eau diminue drastiquement la résistance du corps]
Le corps humain sec présente une résistance élevée (quelques \si{k\ohm} à \SI{100}{k\ohm}), mais la peau mouillée ou des pieds dans l'eau font chuter cette résistance à \SI{1}{k\ohm} ou moins. Pour une même tension de \SI{220}{V}, l'intensité traversant le corps ($I = U/R$) devient alors bien plus grande (plusieurs centaines de mA), mortelle. C'est pourquoi il est \textbf{formellement interdit} de manipuler un appareil électrique (sèche-cheveux, fer, prise, interrupteur) avec les mains mouillées, les pieds dans l'eau ou dans une salle de bain non aux normes.
\end{attention}

\courssub{Les dispositifs de protection de l'installation}

Pour se protéger, toute installation domestique moderne comporte un \textbf{tableau de répartition} (ou tableau électrique) où sont placés les organes de protection.

\begin{popfigure}
\begin{tikzpicture}[
    scale=0.95,
    every node/.style={font=\small},
    protblock/.style={draw, thick, fill=popgreenL, rounded corners=3pt, minimum width=2.2cm, minimum height=1.4cm, align=center},
    diffblock/.style={draw, thick, fill=poporangeL, rounded corners=3pt, minimum width=2.2cm, minimum height=1.6cm, align=center},
    wire/.style={thick},
    earthwire/.style={thick, popgreen!70!black},
    node distance=0.8cm and 0.5cm
]
% Arrivée ENEO
\node[left] (eneo_l) at (0,3.5) {Arrivée ENEO};
\node[left] (eneo_n) at (0,2.5) {Neutre};
\draw[wire] (eneo_l) -- (1.5,3.5);
\draw[wire] (eneo_n) -- (1.5,2.5);

% Disjoncteur Différentiel Général
\node[diffblock, align=center] (diff) at (3.5,3.0){Disjoncteur\\différentiel\\30 mA\\Type A / AC};
\draw[wire] (1.5,3.5) -- (diff.north west);
\draw[wire] (1.5,2.5) -- (diff.south west);

% Barres de répartition (Phase / Neutre)
\draw[wire, popblue] (diff.north east) -- ++(0.8,0) coordinate (bus_ph);
\draw[wire] (diff.south east) -- ++(0.8,0) coordinate (bus_n);
% Barre de terre
\draw[earthwire] (diff.south) -- ++(0,-0.5) -| ++(2.5,-1.5) coordinate (earth_bar);

% Disjoncteurs Divisionnaires
\foreach \i/\name/\y in {1/Prises/3.5, 2/Éclairage/2.5, 3/Chauffe-eau/1.5} {
    \node[protblock, anchor=west] (dj\i) at ($(bus_ph)+(1.2, \y-3.0)$) {Disj. divisionnaire\\\name};
    \draw[wire] (bus_ph) -- (dj\i.north west);
    \draw[wire] (bus_n) -- (dj\i.south west);
    \draw[wire] (dj\i.north east) -- ++(0.7,0);
    \draw[wire] (dj\i.south east) -- ++(0.7,0);
    % Terre vers chaque circuit
    \draw[earthwire] (dj\i.south) |- (earth_bar);
}

% Prise de terre
\draw[earthwire] (earth_bar) -- ++(0,-0.8) coordinate (earth_end);
\draw[earthwire] (earth_end) -- ++(-0.4,0) -- ++(0.8,-0.3) -- ++(-0.8,-0.3) -- ++(0.8,-0.3);
\node[popgreen!70!black, right=0.2cm of earth_end, anchor=west] {Prise de terre (piquet)};

% Légende interne simplifiée
\node[align=left, font=\tiny, text=popdark, anchor=north west] at (12, 4.5) {
    \textbf{Légende :}\\
    \textcolor{popblue}{---} Phase / Neutre (circuits)\\
    \textcolor{popgreen!70!black}{---} Conducteur de protection (Vert-Jaune)\\
    \textcolor{poporange}{■} Protection personnes (Diff. 30 mA)\\
    \textcolor{popgreen}{■} Protection circuits (Div. 10/16/20 A)
};
\end{tikzpicture}
\legende{Schéma simplifié d'un tableau de répartition domestique : 1) Le disjoncteur différentiel \SI{30}{mA} (orange) protège les personnes en détectant les fuites de courant. 2) Les disjoncteurs divisionnaires (verts) protègent chaque circuit (prises, éclairage, chauffe-eau) contre les surcharges et courts-circuits. 3) Le conducteur vert-jaune relie les carcasses métalliques des appareils à la prise de terre (piquet).}
\end{popfigure}

\textbf{Le disjoncteur différentiel (protection des personnes).}\par
En fonctionnement normal, le courant qui part par la phase revient intégralement par le neutre : les deux intensités sont égales. Si un appareil est défectueux (isolement abîmé), une partie du courant « fuit » vers la terre (par la carcasse métallique, puis le conducteur de protection, ou à travers le corps d'une personne) : l'intensité dans la phase devient supérieure à celle dans le neutre. Le disjoncteur différentiel détecte cette différence et coupe le circuit.

\begin{propriete}[Fonctionnement du disjoncteur différentiel]
Le \cle{disjoncteur différentiel} compare en permanence l'intensité du courant qui entre (phase) et celle qui ressort (neutre). Si la différence (courant de fuite) dépasse sa sensibilité nominale, typiquement
\[ I_{\text{fuite}} = |I_{\text{phase}} - I_{\text{neutre}}| > \SI{30}{mA}, \]
il coupe automatiquement l'alimentation en moins de \SI{40}{ms} (pour les modèles domestiques), protégeant ainsi les personnes contre l'électrocution.
\end{propriete}

\textbf{Le disjoncteur divisionnaire et le fusible (protection des circuits).}\par
Le \cle{disjoncteur divisionnaire} (ou disjoncteur modulaire) protège un circuit précis (prises, éclairage, four, chauffe-eau) contre les \textbf{surcharges} (trop d'appareils branchés) et les \textbf{courts-circuits} (contact direct phase-neutre) : si l'intensité dépasse sa valeur nominale (ex. \SI{16}{A} pour un circuit de prises, \SI{20}{A} pour un four), il « saute » (se déclenche) et coupe le circuit. On le réarme après avoir corrigé le défaut.
Le \cle{fusible} (ou coupe-circuit à fusible) joue le même rôle mais de façon \textbf{destructible} : c'est un fil ou une lame calibrée qui \textbf{fond} quand l'intensité devient trop grande pendant un certain temps. Une fois fondu, il doit être remplacé par un fusible \textbf{identique} (même calibre, même type : gG, aM, etc.).

\begin{attention}[Ne jamais remplacer un fusible par un fil de cuivre ou un clou]
Quand un fusible grille, c'est le signe d'un défaut réel (surcharge ou court-circuit). Remplacer le fusible par un fil de cuivre ordinaire — ou par un clou, comme on le voit parfois — supprime toute protection : le fil de cuivre ne fondra pas, et en cas de défaut ce sont les câbles de la maison qui chaufferont par effet Joule ($P = R \cdot I^{2}$) jusqu'à l'incendie. On remplace toujours un fusible par un fusible \textbf{de même calibre et de même type}, après avoir recherché et réparé la cause du défaut.
\end{attention}

\courssub{La mise à la terre et le conducteur de protection}

\textbf{Observation.}\par
Regarde une prise de courant moderne (type E, standard au Cameroun et en France) : en plus des deux trous (phase et neutre), elle possède une broche métallique saillante. Cette broche est reliée au \textbf{conducteur de protection}, un fil bicolore \textbf{vert-jaune}, lui-même connecté à une tige métallique (piquet) enfoncée dans le sol : la \cle{prise de terre}.

\begin{definition}[Mise à la terre]
La \cle{mise à la terre} consiste à relier les parties métalliques accessibles des appareils (carcasse de machine à laver, réfrigérateur, chauffe-eau, tableau électrique) au sol par le conducteur de protection vert-jaune. Si un défaut d'isolement met la carcasse sous tension (phase touchant la tôle), le courant de fuite s'écoule vers la terre par ce conducteur (faible résistance) au lieu de traverser le corps de l'utilisateur (haute résistance) ; le disjoncteur différentiel détecte cette fuite et coupe le circuit.
\end{definition}

La terre et le différentiel forment donc un \textbf{couple inséparable} : la terre offre un chemin de faible impédance au courant de fuite, et le différentiel détecte cette fuite et coupe. Sans prise de terre, le différentiel ne peut pas détecter le défaut avant que quelqu'un ne touche la carcasse (le courant de fuite ne circule pas).

\begin{methode}[Vérifier la continuité de la prise de terre avec un multimètre]
\begin{enumerate}
\item \textbf{Couper le courant} au disjoncteur général (ou au disjoncteur de branchement) : on ne travaille \textbf{jamais} sur une installation sous tension.
\item Régler le multimètre sur la fonction \textbf{ohmmètre} (position $\Omega$, calibre le plus bas, ex. \SI{200}{\ohm}) ou sur le testeur de continuité (symbole diode/buzzer).
\item Placer une pointe de touche sur la broche de terre de la prise murale (ou la borne de terre d'une rallonge) et l'autre pointe sur une borne de terre de référence (barrette de terre du tableau, ou la carcasse métallique d'un appareil voisin dont on est sûr qu'il est bien mis à la terre).
\item \textbf{Lire la mesure} : une résistance très faible (proche de \SI{0}{\ohm}, généralement $< \SI{1}{\ohm}$, avec signal sonore continu) indique que la liaison de terre est continue et correcte.
\item Si l'ohmmètre affiche une résistance infinie (OL, pas de signal), la liaison de terre est coupée (fil sectionné, borne desserrée, piquet arraché) : la prise ne protège plus, il faut faire réparer l'installation par un électricien qualifié.
\end{enumerate}
Remarque : on peut aussi utiliser un \textbf{testeur de prise} (petit boîtier à trois voyants), qui, branché sur la prise, signale immédiatement l'absence de terre, une inversion phase-neutre ou l'absence de neutre.
\end{methode}

\begin{savaistu}[La norme NF C 15-100 et le contexte camerounais]
En France, la norme \textbf{NF C 15-100} impose les règles des installations électriques domestiques : un disjoncteur différentiel de sensibilité \SI{30}{mA} (Type AC ou A) est obligatoire en tête de chaque circuit de prises de courant, et les salles d'eau (salle de bain, buanderie) sont soumises à des règles encore plus strictes (volumes de protection, liaison équipotentielle). Au Cameroun, les installations modernes neuves suivent des règles équivalentes (normes IEC 60364, reprises par les normes locales). \textbf{Exige toujours} qu'un électricien qualifié installe un différentiel \SI{30}{mA} en tête d'installation et une prise de terre (piquet) dans une maison neuve ou rénovée. C'est une question de vie ou de mort, pas de confort ni d'option budgétaire.
\end{savaistu}

\courssub{Consignes de sécurité et signalisation normalisée}

\begin{aretenir}[Les gestes qui sauvent — Règles d'or]
\begin{itemize}
\item \textbf{Couper le courant} au disjoncteur général (ou divisionnaire concerné) avant toute intervention sur l'installation (changer une prise, un interrupteur, réparer un fil, changer une ampoule).
\item Ne \textbf{jamais} manipuler un appareil électrique (fiche, interrupteur, appareil) avec les mains mouillées, les pieds nus sur un sol humide, ou dans un local humide (salle de bain, buanderie) sauf si l'installation est aux normes (volumes, différentiel 30 mA).
\item Ne jamais mettre les doigts ou un objet métallique (couteau, tournevis, fil de fer) dans une prise de courant.
\item Respecter les \textbf{sections de câbles} adaptées à l'intensité du circuit : ex. \SI{1,5}{mm^{2}} pour l'éclairage (\SI{10}{A} ou \SI{16}{A}), \SI{2,5}{mm^{2}} pour les prises (\SI{16}{A} ou \SI{20}{A}), \SI{6}{mm^{2}} ou plus pour le chauffe-eau / four (\SI{32}{A}). Ne pas brancher plusieurs appareils puissants sur une même rallonge ou multiprise bas de gamme.
\item Remplacer un fusible fondu par un fusible de \textbf{même calibre et même type}, jamais par un fil de cuivre, un trombone ou un clou.
\item Ne pas surcharger les multiprises (regarder la puissance max indiquée, ex. \SI{2500}{W} ou \SI{3500}{W}) ; débrancher les appareils sensibles (TV, ordi, frigo) en cas d'orage violent ou d'inondation.
\item En cas d'accident électrique : \textbf{1. Couper le courant} (disjoncteur général) ; \textbf{2. Ne jamais toucher directement la victime} tant qu'elle est en contact avec le circuit (risque de second accident) ; \textbf{3. Appeler les secours} (Pompiers \textbf{118} ou SAMU) ; \textbf{4. Pratiquer les gestes de premiers secours} (RCP) si la victime ne respire plus, une fois le courant coupé.
\end{itemize}
\end{aretenir}

\textbf{La signalisation normalisée (ISO 7010).}\par
Sur les appareils, les tableaux électriques, les transformateurs ENEO et les lieux de travail, des pictogrammes normalisés indiquent les dangers et les règles à respecter. Il faut savoir les reconnaître :

\begin{itemize}
\item \textbf{Signalisation de DANGER} (triangle jaune bordé de noir, symbole noir) : un éclair dans un triangle signale un \textbf{danger électrique} — risque de choc électrique, haute tension. Exemple : sur un transformateur ENEO, sur un tableau électrique, sur un groupe électrogène.
\item \textbf{Signalisation d'INTERDICTION} (cercle rouge barré sur fond blanc, symbole noir) : « Ne pas toucher », « Accès interdit aux personnes non autorisées », « Ne pas éteindre à l'eau » (sur un feu d'origine électrique).
\item \textbf{Signalisation d'OBLIGATION} (disque bleu, symbole blanc) : « Port de gants isolants obligatoire », « Couper l'alimentation avant intervention », « Port de chaussures de sécurité isolantes obligatoire ».
\item \textbf{Signalisation de SECOURS / SAUVETAGE} (carré ou rectangle vert, symbole blanc) : Issue de secours, poste de premiers secours, extincteur (souvent rouge avec pictogramme blanc).
\end{itemize}

\begin{attention}[Un feu d'origine électrique ne s'éteint JAMAIS à l'eau]
L'eau conduit le courant : jeter de l'eau sur un appareil en feu encore branché, c'est risquer l'électrocution du sauveteur et propager le feu. La bonne conduite : \textbf{1. Couper le courant} au disjoncteur général ; \textbf{2. Utiliser un extincteur adapté} : poudre polyvalente (ABC) ou \ce{CO2} (dioxyde de carbone, neige carbonique) — \textbf{pas d'eau, pas de mousse} ; \textbf{3. Si pas d'extincteur}, étouffer les flammes avec une couverture épaisse (couverture de survie, tapis non synthétique) ou du sable sec.
\end{attention}

\begin{exoresolu}[Diagnostiquer une installation dangereuse]
Dans une maison, on constate trois anomalies :\\
\textbf{(a)} Le fusible du circuit des prises a été remplacé par un fil de cuivre.\\
\textbf{(b)} La machine à laver donne des « picotements » (petits chocs) quand on la touche pendant qu'elle fonctionne.\\
\textbf{(c)} Trois appareils puissants (fer, bouilloire, micro-ondes) sont branchés sur une même rallonge à enrouleur.

Pour chaque point, identifier le risque principal et proposer la mesure de sécurité correcte.
\tcblower
\textbf{(a) Fusible remplacé par un fil de cuivre.}\\
\textit{Risque :} En cas de surcharge ou de court-circuit, plus rien ne coupe le circuit (le cuivre ne fond pas). Les câbles de l'installation chauffent par effet Joule ($P = R \cdot I^{2}$) jusqu'à la fusion de l'isolant et l'incendie.\\
\textit{Mesure :} Couper le courant. Remettre un fusible de calibre adapté (ex. \SI{16}{A} gG pour circuit prises \SI{2,5}{mm^{2}}) ou, mieux, installer un disjoncteur divisionnaire. Rechercher la cause de la fonte du fusible initial (surcharge ? court-circuit ?).

\textbf{(b) Picotements sur la machine à laver.}\\
\textit{Risque :} Un défaut d'isolement interne met la carcasse métallique sous tension (fuite de courant). Le courant traverse le corps de l'utilisateur vers la terre. Le seuil de fibrillation (\SI{\approx 50}{mA}) peut être atteint. Danger de mort par électrocution.\\
\textit{Mesure :} \textbf{1. Débrancher l'appareil immédiatement.} \textbf{2. Vérifier la continuité de la prise de terre} (méthode ci-dessus) : la broche de terre de la prise est-elle bien reliée au piquet ? \textbf{3. Vérifier la présence et le bon fonctionnement du disjoncteur différentiel \SI{30}{mA}} (bouton test « T »). \textbf{4. Faire réparer la machine} (changement de la résistance ou du moteur défectueux) par un professionnel avant de la rebrancher.

\textbf{(c) Rallonge surchargée (et enrouleur non déroulé).}\\
\textit{Risque :} Surchauffe du câble de la rallonge par effet Joule (section insuffisante pour l'intensité appelée). Si l'enrouleur n'est pas totalement déroulé, l'effet d'auto-induction (bobine) aggrave l'échauffement. Fonte de l'isolant, court-circuit, incendie.\\
\textit{Mesure :} Répartir les appareils sur plusieurs prises murales distinctes (circuits différents). Utiliser des rallonges de section adaptée (\SI{2,5}{mm^{2}} pour fortes puissances) et \textbf{toujours dérouler complètement} les enrouleurs. Ne pas chaîner les multiprises.
\end{exoresolu}

\begin{exercice}[À toi de jouer]
1. Explique pourquoi un disjoncteur différentiel \SI{30}{mA} protège efficacement une personne contre l'électrocution, alors qu'un fusible (ou disjoncteur divisionnaire) de \SI{16}{A} ne le peut pas.\par
2. Une installation possède une prise de terre correcte (piquet, fil vert-jaune relié aux carcasses) mais \textbf{pas} de disjoncteur différentiel. Un défaut met la carcasse d'un lave-linge sous tension. La protection est-elle efficace ? Justifie ta réponse en décrivant ce qui se passe si quelqu'un touche la machine.\par
3. Cite trois pictogrammes normalisés que l'on doit trouver obligatoirement sur ou près d'un transformateur de distribution ENEO (poste de transformation) et précise pour chacun sa forme, sa couleur et sa signification.
\end{exercice}

\begin{corrige}[Exercice]
1. Un fusible ou disjoncteur divisionnaire de \SI{16}{A} ne se déclenche que si l'intensité dépasse \SI{16}{A} (soit plus de 300 fois le seuil de fibrillation \SI{\approx 50}{mA}). Un courant mortel de \SI{100}{mA} ou \SI{500}{mA} traversant une personne ne le fait pas déclencher : il protège les \textbf{fils}, pas les gens. Le différentiel \SI{30}{mA}, lui, coupe dès que la fuite dépasse \SI{30}{mA}, donc \textbf{bien avant} que le courant traversant le corps n'atteigne des valeurs mortelles. Il protège les \textbf{personnes}.

2. \textbf{Non, la protection n'est pas complète.}\\
Sans différentiel, le courant de fuite du lave-linge s'écoule bien vers la terre par le fil vert-jaune (faible résistance), mais \textbf{rien ne coupe le circuit automatiquement}. Le défaut persiste. La carcasse reste sous tension (proche de \SI{220}{V} par rapport à la terre). Si quelqu'un touche la carcasse, il se met en parallèle avec la liaison de terre : comme la résistance du corps (\SI{\approx 1}{k\ohm} mouillé) est bien plus grande que celle de la liaison de terre (\SI{< 1}{\ohm}), la majeure partie du courant continue de passer par la terre, mais une partie non négligeable (plusieurs dizaines de mA) traverse la personne. Le danger d'électrocution reste réel. Le différentiel est indispensable pour \textbf{couper l'alimentation} dès l'apparition de la fuite.

3. Exemples obligatoires près d'un transformateur ENEO :\\
\begin{itemize}
\item \textbf{Triangle jaune bordé noir + éclair noir} : \textit{Danger} — « Haute tension / Danger de mort / Risque d'électrocution ».
\item \textbf{Cercle rouge barré + symbole « personne + éclair » ou « clé »} : \textit{Interdiction} — « Accès interdit aux personnes non autorisées / Ne pas manipuler / Travaux interdits sans consignation ».
\item \textbf{Disque bleu + symbole « casque » ou « gants » ou « couper courant »} : \textit{Obligation} — « Port du casque / gants isolants obligatoire » ou « Couper l'alimentation avant toute intervention ».
\end{itemize}
(Autre pictogramme fréquent : carré rouge + symbole extincteur = localisation extincteur).
\end{corrige}

\begin{aretenir}[L'essentiel de la section — Sécurité électrique]
\begin{itemize}
\item \textbf{4 risques majeurs :} Électrocution (courant dans le corps), Court-circuit (contact direct phase-neutre, $I$ énorme), Surchauffe / Surcharge ($I$ trop fort pour la section du câble, effet Joule durable), Incendie (conséquence des deux précédents).
\item \textbf{Protection des personnes :} \textbf{Disjoncteur différentiel \SI{30}{mA}} (coupe si fuite $> \SI{30}{mA}$). Seuil de danger : \SI{\approx 50}{mA} (fibrillation).
\item \textbf{Protection des circuits :} \textbf{Disjoncteurs divisionnaires} (réarmables) ou \textbf{Fusibles} (à changer) — calibrés selon la section des câbles (\SI{16}{A} pour \SI{2,5}{mm^{2}}, etc.). Protègent contre surcharges et courts-circuits.
\item \textbf{Mise à la terre (conducteur vert-jaune + piquet) :} Évacue les courants de fuite vers le sol. Indispensable pour que le différentiel détecte le défaut.
\item \textbf{Couple inséparable :} Terre + Différentiel = Protection efficace des personnes.
\item \textbf{Consignes vitales :} Couper le courant avant intervention \textbf{!} Pas d'électricité avec mains mouillées / pieds dans l'eau \textbf{!} Pas de fil de cuivre à la place d'un fusible \textbf{!} Pas d'eau sur feu électrique \textbf{!} Respecter la signalisation (Triangle=Danger, Rond barré=Interdit, Disque bleu=Obligation).
\end{itemize}
\end{aretenir}

\coursec{Applications concrètes et résolution de problèmes}

Dans la section précédente, nous avons vu comment la sécurité électrique repose sur des règles précises : fusibles et disjoncteurs adaptés, prise de terre, respect des intensités admissibles. Mais savoir se protéger ne suffit pas : il faut aussi savoir \cle{choisir} son matériel, \cle{calculer} sa consommation et \cle{payer} sa facture ENEO en toute connaissance de cause. C'est l'objet de cette dernière section : utiliser les formules $P = U \cdot I$ et $E = P \cdot t$ comme de véritables outils de la vie quotidienne, du choix d'une ampoule au dimensionnement d'un câble, jusqu'à la lecture de la facture mensuelle d'un ménage camerounais.

\courssub{Choisir une ampoule : puissance, flux lumineux et efficacité}

\textbf{Situation de départ.} À la boutique du quartier, tu hésites entre trois ampoules : une incandescente de \SI{100}{W}, une halogène de \SI{75}{W} et une LED de \SI{10}{W}. Le vendeur affirme qu'elles éclairent « pareil ». Comment vérifier ?

\textbf{L'intuition.} La puissance en watts indique ce que l'ampoule \emph{consomme}, pas ce qu'elle \emph{éclaire}. Une vieille ampoule à incandescence transforme plus de 90 \% de l'énergie en chaleur : c'est presque un petit chauffage ! Pour comparer l'éclairage, il faut une autre grandeur : le \cle{flux lumineux}, mesuré en \cle{lumens} (lm).

\begin{definition}[Flux lumineux et efficacité lumineuse]
Le \cle{flux lumineux} d'une lampe, noté $\Phi$, est la quantité de lumière qu'elle émet ; il se mesure en \cle{lumens} (lm). L'\cle{efficacité lumineuse} est le rapport du flux lumineux à la puissance électrique consommée :
\[
\eta = \dfrac{\Phi}{P} \quad \text{en lumens par watt (lm/W)}.
\]
Plus l'efficacité est grande, plus la lampe produit de lumière pour la même dépense d'électricité.
\end{definition}

\begin{center}
\begin{tabularx}{\linewidth}{|l|X|X|X|}
\hline
\rowcolor{popblueL} \textbf{Type de lampe} & \textbf{Puissance} & \textbf{Flux lumineux} & \textbf{Efficacité} \\
\hline
Incandescence & \SI{100}{W} & environ \SI{1300}{lm} & environ \SI{13}{lm/W} \\
\hline
Halogène & \SI{75}{W} & environ \SI{1300}{lm} & environ \SI{17}{lm/W} \\
\hline
LED & \SI{10}{W} & environ \SI{1000}{lm} & environ \SI{100}{lm/W} \\
\hline
\end{tabularx}
\end{center}

\textbf{Conclusion de l'investigation.} À flux lumineux comparable, la LED consomme 7 à 10 fois moins que l'halogène ou l'incandescente. Le bon réflexe à l'achat : regarder d'abord les \emph{lumens} (pour la lumière voulue), puis les \emph{watts} (pour la dépense).

\begin{savaistu}[La révolution LED]
Une LED de \SI{10}{W} éclaire autant qu'une halogène de \SI{75}{W} : cela représente une économie d'environ 85 \% sur la consommation d'éclairage ! De plus, une LED dure 15 000 à 25 000 heures, contre 1 000 heures pour une incandescente. Dans un pays où chaque kilowattheure compte, remplacer les ampoules des maisons, des écoles et des boutiques par des LED est l'un des gestes d'économie d'énergie les plus rentables.
\end{savaistu}

\begin{attention}[Piège fréquent]
Ne confonds jamais la \textbf{puissance} (W), qui mesure la consommation, avec le \textbf{flux lumineux} (lm), qui mesure la lumière produite. Dire « cette ampoule de 100 W éclaire plus » n'est vrai qu'entre ampoules de \emph{même technologie}. Une LED de 10 W éclaire autant qu'une incandescente de 75 à 100 W.
\end{attention}

\courssub{Dimensionner un câble électrique}

\textbf{Situation de départ.} Ton oncle installe un chauffe-eau dans la salle de bain. Il hésite entre un câble fin (moins cher) et un câble plus gros. Que risque-t-on avec un câble trop fin ?

\textbf{L'intuition.} Quand un courant traverse un conducteur, celui-ci s'échauffe (effet Joule : $P = R \cdot I^2$). Plus le câble est fin, plus sa résistance est grande, plus il chauffe. Un câble sous-dimensionné peut fondre son isolant et provoquer un incendie. Il existe donc, pour chaque \cle{section} de câble (surface de cuivre, en mm\textsuperscript{2}), une \cle{intensité admissible} maximale.

\begin{methode}[Choisir la section d'un câble]
\begin{enumerate}
\item \textbf{Calculer l'intensité} du circuit : $I = \dfrac{P}{U}$ (avec $U = \SI{230}{V}$ au Cameroun).
\item \textbf{Consulter le tableau} des intensités admissibles :
\begin{center}
\begin{tabular}{|c|c|c|}
\hline
\rowcolor{popblueL} \textbf{Section} & \textbf{Intensité max} & \textbf{Usage typique} \\
\hline
\SI{1,5}{mm^2} & $\leq \SI{16}{A}$ & éclairage, petit chauffage \\
\hline
\SI{2,5}{mm^2} & $\leq \SI{25}{A}$ & prises de courant, chauffe-eau \\
\hline
\SI{6}{mm^2} & $\leq \SI{40}{A}$ & cuisinière, gros appareils \\
\hline
\end{tabular}
\end{center}
\item \textbf{Choisir la section} dont l'intensité admissible est \emph{supérieure} à l'intensité calculée, avec une marge de sécurité.
\item \textbf{Vérifier} que le disjoncteur ou fusible de protection correspond à la section choisie.
\end{enumerate}
\end{methode}

\begin{exemplebox}[Chauffage de 2000 W sous 230 V]
Un radiateur de puissance $P = \SI{2000}{W}$ est branché sur le réseau $U = \SI{230}{V}$.
\[
I = \dfrac{P}{U} = \dfrac{2000}{230} \approx \SI{8,7}{A}.
\]
Comme $8{,}7 < 16$ A, un câble de section \SI{1,5}{mm^2} est suffisant. En pratique, on choisit souvent \SI{2,5}{mm^2} pour un chauffage fixe, afin d'avoir une marge de sécurité confortable.
\end{exemplebox}

\begin{attention}[Ni trop fin, ni trop gros]
\textbf{Sous-dimensionner} un câble entraîne un échauffement dangereux : isolant qui fond, court-circuit, incendie. C'est un risque réel quand on branche un gros appareil sur une rallonge fine. À l'inverse, \textbf{sur-dimensionner} inutilement augmente le coût de l'installation (le cuivre est cher) sans gain de sécurité. Le bon choix est le juste dimensionnement, avec une marge raisonnable.
\end{attention}

\courssub{Calculer la facture mensuelle d'un ménage}

\begin{propriete}[Énergie facturée]
L'énergie électrique consommée par une installation est la somme des énergies consommées par chaque appareil :
\[
E = \sum_i P_i \times t_i
\]
où $P_i$ est la puissance de chaque appareil en kilowatts (kW) et $t_i$ sa durée de fonctionnement en heures (h). L'énergie s'exprime alors en \cle{kilowattheures} (kWh), l'unité qui figure sur la facture ENEO. Le coût s'obtient en multipliant $E$ par le prix du kilowattheure.
\end{propriete}

\begin{methode}[Établir un bilan énergétique simple d'un logement]
\begin{enumerate}
\item \textbf{Recenser} les appareils : noter la puissance de chacun (plaque signalétique).
\item \textbf{Estimer} la durée quotidienne d'utilisation de chaque appareil.
\item \textbf{Calculer} l'énergie quotidienne de chaque appareil : $E_i = P_i \times t_i$ (en kWh si $P$ en kW et $t$ en h).
\item \textbf{Additionner} toutes les énergies, puis multiplier par 30 pour le mois.
\item \textbf{Multiplier} par le prix du kWh pour obtenir le coût mensuel.
\item \textbf{Identifier} les postes les plus coûteux et proposer des économies.
\end{enumerate}
\end{methode}

\begin{exoresolu}[Facture mensuelle de la famille Ngo Bell]
\textbf{Énoncé.} La famille Ngo Bell utilise chaque jour : 5 lampes LED de \SI{10}{W} pendant 5 h ; un téléviseur de \SI{100}{W} pendant 4 h ; un réfrigérateur qui consomme en moyenne \SI{1,2}{kWh} par jour ; un fer à repasser de \SI{1200}{W} pendant 30 min. Le kWh coûte 100 FCFA. Calculer l'énergie mensuelle (30 jours) et le montant de la facture.
\tcblower
\textbf{Solution détaillée.} On applique $E_i = P_i \times t_i$ pour chaque appareil, en convertissant les puissances en kW.
\begin{align*}
E_{\text{lampes}} &= 5 \times 0{,}010 \times 5 = \SI{0,25}{kWh/jour} \\
E_{\text{TV}} &= 0{,}100 \times 4 = \SI{0,40}{kWh/jour} \\
E_{\text{frigo}} &= \SI{1,2}{kWh/jour} \\
E_{\text{fer}} &= 1{,}2 \times 0{,}5 = \SI{0,60}{kWh/jour}
\end{align*}
Énergie quotidienne totale :
\[
E_{\text{jour}} = 0{,}25 + 0{,}40 + 1{,}2 + 0{,}60 = \SI{2,45}{kWh}.
\]
Énergie mensuelle : $E = 2{,}45 \times 30 = \SI{73,5}{kWh}$.

Coût mensuel : $73{,}5 \times 100 = 7\,350$ FCFA.

\textbf{Conclusion.} La famille paie environ 7 350 FCFA par mois. Le réfrigérateur est le poste le plus coûteux (environ 36 kWh, soit près de la moitié de la facture).
\end{exoresolu}

\courssub{Interpréter une courbe puissance-temps : le chauffage à accumulation}

\textbf{Situation de départ.} Certains chauffages « à accumulation » fonctionnent par périodes : ils chauffent l'eau ou des briques réfractaires pendant quelques heures, puis restituent la chaleur. Leur consommation n'est pas continue : comment la représenter et la calculer ?

\begin{popfigure}
\begin{center}
\begin{tikzpicture}
\begin{axis}[
    width=0.95\linewidth, height=6cm,
    xlabel={Temps (h)}, ylabel={Puissance (W)},
    xmin=0, xmax=24, ymin=0, ymax=2500,
    xtick={0,2,5,8,18,21,24},
    ytick={0,500,1000,1500,2000},
    grid=major, grid style={popline!50},
    axis lines=left,
    every axis x label/.style={at={(ticklabel* cs:1.0)},anchor=west},
    every axis y label/.style={at={(ticklabel* cs:1.0)},anchor=south},
]
\addplot[popcurve, very thick, popblue, const plot] coordinates {
    (0,0) (5,2000) (7,0) (18,2000) (21,0) (24,0)
};
\addplot[poporange, dashed, thick] coordinates {(0,2000) (24,2000)};
\node[poporange, anchor=south west] at (axis cs:0.3,2050) {\small $P = \SI{2000}{W}$};
\node[popdark, anchor=north] at (axis cs:6,2000) {\small matin : 5 h--7 h};
\node[popdark, anchor=north] at (axis cs:19.5,2000) {\small soir : 18 h--21 h};
\end{axis}
\end{tikzpicture}
\end{center}
\legende{Courbe puissance--temps d'un chauffage de \SI{2000}{W} fonctionnant 2 h le matin (5 h à 7 h) et 3 h le soir (18 h à 21 h). L'énergie consommée est l'aire sous la courbe : $E = 2000 \times 2 + 2000 \times 3 = \SI{10}{kWh}$ par jour.}
\end{popfigure}

\begin{aretenir}[Lire une courbe puissance--temps]
Sur un graphique $P(t)$, l'\cle{énergie consommée} est l'\textbf{aire comprise entre la courbe et l'axe du temps}. Pour un fonctionnement par paliers (puissance constante pendant une durée $t$), chaque palier contribue pour $P \times t$. Ici : $E = 2000 \times 2 + 2000 \times 3 = 10\,000$ Wh $= \SI{10}{kWh}$ par jour, soit environ \SI{300}{kWh} par mois : le chauffage est souvent le premier poste de consommation d'un logement !
\end{aretenir}

\courssub{Proposer des économies d'énergie}

À partir du bilan énergétique, on peut agir sur trois leviers :

\begin{itemize}
\item \textbf{Remplacer les appareils énergivores} : passer des halogènes aux LED (économie d'environ 85 \% sur l'éclairage), choisir un réfrigérateur de classe énergétique économe.
\item \textbf{Supprimer les consommations cachées} : un téléviseur ou un décodeur en veille consomme 5 à 15 W en permanence. Sur un mois, 10 W de veille représentent $0{,}010 \times 24 \times 30 = \SI{7,2}{kWh}$, soit plus de 700 FCFA jetés ! Éteindre complètement les appareils ou utiliser une multiprise à interrupteur.
\item \textbf{Programmer les usages} : un programmateur fait fonctionner le chauffe-eau uniquement aux heures utiles, au lieu de chauffer 24 h/24. Un chauffe-eau de \SI{1500}{W} ramené de 24 h à 3 h par jour économise $1{,}5 \times 21 \times 30 = \SI{945}{kWh}$ par mois.
\end{itemize}

\begin{exoresolu}[Exercice de synthèse : l'atelier de menuiserie]
\textbf{Énoncé.} Dans son atelier à Bafoussam, M. Kamga utilise une scie électrique de \SI{2200}{W} sous \SI{230}{V}, 3 h par jour, 26 jours par mois. Le kWh coûte 100 FCFA.
\begin{enumerate}
\item Calculer l'intensité du courant dans la scie en fonctionnement.
\item Choisir la section de câble adaptée (\SI{1,5}{mm^2} $\leq$ 16 A ; \SI{2,5}{mm^2} $\leq$ 25 A).
\item Calculer l'énergie mensuelle consommée et son coût.
\item M. Kamga veut brancher en plus un rabot de \SI{1500}{W} sur la même rallonge. Le câble de la rallonge supporte 16 A au total. Est-ce prudent ?
\end{enumerate}
\tcblower
\textbf{Solution détaillée.}
\begin{enumerate}
\item Intensité : $I = \dfrac{P}{U} = \dfrac{2200}{230} \approx \SI{9,6}{A}$.
\item $9{,}6 < 16$ A : le câble de \SI{1,5}{mm^2} convient ; on choisit \SI{2,5}{mm^2} pour plus de marge dans un atelier (poussière, chaleur).
\item Énergie mensuelle : $E = P \times t = 2{,}2 \times 3 \times 26 = \SI{171,6}{kWh}$. Coût : $171{,}6 \times 100 = 17\,160$ FCFA.
\item Courant total si les deux appareils fonctionnent ensemble :
\[
I_{\text{total}} = \dfrac{2200 + 1500}{230} = \dfrac{3700}{230} \approx \SI{16,1}{A}.
\]
C'est \textbf{supérieur} à 16 A : la rallonge va chauffer dangereusement. Il faut soit ne pas utiliser les deux machines en même temps, soit installer une ligne dédiée de section \SI{2,5}{mm^2} protégée par un disjoncteur adapté. La sécurité d'abord !
\end{enumerate}
\end{exoresolu}

\begin{exercice}[À toi de jouer]
Une famille remplace 6 ampoules halogènes de \SI{75}{W} (utilisées 4 h par jour) par 6 LED de \SI{10}{W} de flux lumineux équivalent.
\begin{enumerate}
\item Calculer l'énergie économisée chaque jour, puis chaque mois (30 jours).
\item À 100 FCFA le kWh, calculer l'économie mensuelle en FCFA.
\item Les 6 LED coûtent 12 000 FCFA. En combien de mois l'investissement est-il remboursé ?
\end{enumerate}
\end{exercice}

\begin{corrige}[Exercice « À toi de jouer »]
\begin{enumerate}
\item Puissance économisée par ampoule : $75 - 10 = \SI{65}{W}$. Pour 6 ampoules : $6 \times 65 = \SI{390}{W} = \SI{0,39}{kW}$. Énergie quotidienne économisée : $E = 0{,}39 \times 4 = \SI{1,56}{kWh}$. Par mois : $1{,}56 \times 30 = \SI{46,8}{kWh}$.
\item Économie mensuelle : $46{,}8 \times 100 = 4\,680$ FCFA.
\item Durée de remboursement : $\dfrac{12\,000}{4\,680} \approx 2{,}6$ mois, soit moins de 3 mois. Ensuite, c'est du gain net !
\end{enumerate}
\end{corrige}

\begin{aretenir}[L'essentiel de la section]
\begin{itemize}
\item Pour choisir une ampoule : comparer les \cle{lumens} (lumière) et l'\cle{efficacité} (lm/W), pas seulement les watts.
\item Pour choisir un câble : calculer $I = P/U$, puis prendre une section dont l'intensité admissible dépasse $I$ (\SI{1,5}{mm^2} $\leq$ 16 A ; \SI{2,5}{mm^2} $\leq$ 25 A).
\item L'énergie facturée est $E = \sum P_i \times t_i$ en kWh ; le coût est $E$ multiplié par le prix du kWh.
\item Sur une courbe $P(t)$, l'énergie est l'aire sous la courbe.
\item Les grandes économies : LED, extinction des veilles, programmation des appareils de chauffage.
\end{itemize}
\end{aretenir}

\newpage

\coursec{Activité d'intégration}

\coursec{Leçon 14 : Énergie électrique et puissance électrique}

\courssub{Objectifs de l'activité}

À l'issue de cette activité d'intégration, tu seras capable de :

\begin{itemize}
\item lire un relevé de compteur électrique et calculer une consommation d'énergie en kilowattheures (\si{kWh}) ;
\item utiliser les relations $P = U \times I$, $P = R \times I^2$ et $E = P \times t$ dans une situation concrète ;
\item identifier les appareils les plus énergivores d'une habitation ;
\item vérifier la conformité des protections électriques (disjoncteur différentiel, calibre des disjoncteurs divisionnaires) ;
\item proposer des actions d'économie d'énergie chiffrées et argumentées ;
\item présenter tes résultats à l'oral à l'aide d'une affiche.
\end{itemize}

Cette activité mobilise l'ensemble des savoirs et savoir-faire de la leçon : grandeurs électriques, instruments de mesure, puissance, énergie et sécurité.

\courssub{Situation : Audit énergétique de la maison de la famille Ngo Bell à Yaoundé}

\textbf{Contexte.} La famille Ngo Bell habite un quartier de Yaoundé. Elle est alimentée en électricité par le réseau ENEO sous une tension de $U = \SI{220}{V}$. À la fin du mois, Monsieur Ngo Bell trouve sa facture d'électricité trop élevée. Il demande à son fils Alain, élève en classe de 3ème, de l'aider à comprendre d'où vient cette consommation et comment la réduire.

Alain décide de réaliser un \cle{audit énergétique} de la maison avec l'aide de ses camarades. Il dispose des documents suivants.

\textbf{Document 1 : relevés du compteur électrique sur 7 jours.}

\begin{center}
\begin{tabularx}{\linewidth}{|l|X|X|}
\hline
\rowcolor{popblueL} \textbf{Relevé} & \textbf{Date et heure} & \textbf{Index du compteur} \\
\hline
Relevé 1 & Lundi 06 h 00 & \num{12450,0} \si{kWh} \\
\hline
Relevé 2 & Lundi suivant 06 h 00 & \num{12534,5} \si{kWh} \\
\hline
\end{tabularx}
\end{center}

\textbf{Document 2 : fiches techniques des appareils de la maison.}

\begin{center}
\begin{tabularx}{\linewidth}{|l|X|X|X|}
\hline
\rowcolor{popblueL} \textbf{Appareil} & \textbf{Puissance nominale} & \textbf{Tension} & \textbf{Durée d'utilisation par jour} \\
\hline
Réfrigérateur & \SI{150}{W} & \SI{220}{V} & 8 h (fonctionnement effectif) \\
\hline
Téléviseur & \SI{100}{W} & \SI{220}{V} & 5 h \\
\hline
4 lampes halogènes & $4 \times \SI{50}{W}$ & \SI{220}{V} & 6 h \\
\hline
Fer à repasser & \SI{1200}{W} & \SI{220}{V} & 0,5 h \\
\hline
Bouilloire électrique & \SI{2000}{W} & \SI{220}{V} & 0,3 h \\
\hline
Ventilateur & \SI{75}{W} & \SI{220}{V} & 8 h \\
\hline
\end{tabularx}
\end{center}

\textbf{Document 3 : observations du tableau électrique.}

En examinant le tableau électrique, Alain note les informations suivantes :

\begin{itemize}
\item le disjoncteur général est un disjoncteur différentiel de calibre \SI{30}{A} avec une sensibilité de \SI{500}{mA} ;
\item le circuit d'éclairage est protégé par un disjoncteur divisionnaire de \SI{10}{A} ;
\item le circuit des prises de courant est protégé par un disjoncteur divisionnaire de \SI{16}{A} ;
\item la bouilloire et le fer à repasser sont branchés sur une même rallonge multiprise reliée à une seule prise murale ;
\item plusieurs fils de la rallonge sont dénudés et la prise de terre n'est pas raccordée.
\end{itemize}

\begin{popfigure}
\begin{tikzpicture}[
    box/.style={draw, thick, rounded corners, minimum width=2.5cm, minimum height=1cm, align=center, fill=popblueL},
    wire/.style={draw, thick},
    label/.style={font=\small, text=popdark},
    danger/.style={draw=poporange, thick, fill=poporangeL, rounded corners, minimum width=2cm, minimum height=0.8cm, align=center}
]
% Grid
\node[box, fill=popgreenL, align=center] (eneo) at (0,3){Réseau ENEO\\220 V};
% Main breaker
\node[box, align=center] (diff) at (0,1.5){Disjoncteur général\\Différentiel 30 A / 500 mA};
% Busbar
\draw[wire] (eneo.south) -- (diff.north);
\draw[wire] (diff.south) -- ++(0,-0.5) coordinate (bus);
% Lighting circuit
\draw[wire] (bus) -| ++(-3,-1.5) node[box, fill=poppinkL, align=center] (dl){Disjoncteur\\Éclairage 10 A};
\draw[wire] (dl.south) -- ++(0,-1) node[label, anchor=north, align=center]{Lampes\\Halogènes};
% Socket circuit
\draw[wire] (bus) -| ++(3,-1.5) node[box, fill=poppinkL, align=center] (dp){Disjoncteur\\Prises 16 A};
\draw[wire] (dp.south) -- ++(0,-1) coordinate (socket);
% Power strip
\node[danger, anchor=north, align=center] (strip) at (socket){Rallonge multiprise\\Fils dénudés, pas de terre};
\draw[wire] (socket) -- (strip.north);
% Appliances on strip
\node[label, anchor=north west] at (strip.south west) {Bouilloire 2000 W};
\node[label, anchor=north east] at (strip.south east) {Fer 1200 W};
\draw[wire, poporange] (strip.south west) -- ++(-1.5,-1) node[label, anchor=north] {$\approx$ 9,1 A};
\draw[wire, poporange] (strip.south east) -- ++(1.5,-1) node[label, anchor=north] {$\approx$ 5,5 A};
% Ground symbol missing
\node[label, text=poporange, font=\small\bfseries] at (4.5, -0.5) {$\times$ Prise de terre absente};
\end{tikzpicture}
\legende{Schéma simplifié du tableau électrique et de la rallonge surchargée (Document 3).}
\end{popfigure}

\textbf{Document 4 : tarif ENEO (valeur simplifiée pour l'exercice).}

Le kilowattheure est facturé en moyenne \SI{75}{FCFA} le \si{kWh}.

\textbf{Problème à résoudre :} d'où vient la consommation élevée de la famille ? Les installations sont-elles sûres ? Quelles économies concrètes peut-on proposer ?

\courssub{Consignes}

Travail en groupes de 4 élèves. Durée : 2 séances de 1 h. Restitution orale de 5 minutes par groupe avec une affiche.

\begin{enumerate}
\item \textbf{Consommation hebdomadaire.} À l'aide du document 1, calcule l'énergie électrique consommée par la famille pendant les 7 jours, en \si{kWh}. Déduis-en la consommation moyenne par jour.
\item \textbf{Consommation théorique des appareils.} À l'aide du document 2, calcule l'énergie consommée \emph{par jour} par chaque appareil (en \si{kWh}), puis l'énergie totale théorique par jour. Compare avec le résultat de la question 1 et commente.
\item \textbf{Classement.} Range les appareils du plus énergivore au moins énergivore sur une journée. Quel appareil aurait-on tendance à accuser à tort ? Pourquoi ?
\item \textbf{Courants et protections.} Calcule l'intensité du courant traversant la bouilloire seule, puis le fer seul, puis les deux branchés ensemble sur la rallonge. Le disjoncteur de \SI{16}{A} du circuit de prises déclenche-t-il ? La rallonge est-elle utilisée dans de bonnes conditions ?
\item \textbf{Sécurité des personnes.} La sensibilité du différentiel est-elle conforme à la norme de protection des personnes (\SI{30}{mA}) ? Cite deux autres défauts de sécurité relevés dans le document 3 et explique le danger correspondant.
\item \textbf{Recommandations chiffrées.} Rédige trois recommandations concrètes et chiffrées. Exemple attendu : remplacer les 4 lampes halogènes de \SI{50}{W} par 4 lampes LED de \SI{7}{W} ; calcule l'économie d'énergie réalisée en un mois de 30 jours, puis l'économie en francs CFA.
\item \textbf{Restitution.} Prépare une affiche présentant : le bilan de consommation, les défauts de sécurité et les trois recommandations chiffrées.
\end{enumerate}

\courssub{Ressources à mobiliser}

\begin{aretenir}[Boîte à outils de l'activité]
\begin{itemize}
\item \textbf{Puissance électrique} d'un appareil : $P = U \times I$, avec $P$ en watts (\si{W}), $U$ en volts (\si{V}) et $I$ en ampères (\si{A}).
\item Pour un résistor (appareil purement résistif : résistance, radiateur, bouilloire, fer, lampe à filament) : $P = R \times I^2 = \dfrac{U^2}{R}$.
\item \textbf{Énergie électrique} consommée : $E = P \times t$. Si $P$ est en kilowatts (\si{kW}) et $t$ en heures (h), alors $E$ est en kilowattheures (\si{kWh}). Rappel : $\SI{1}{kW} = \SI{1000}{W}$.
\item \textbf{Consommation lue au compteur} : $E = \text{index final} - \text{index initial}$.
\item \textbf{Intensité} débitée par un appareil : $I = \dfrac{P}{U}$.
\item \textbf{Sécurité} : le disjoncteur différentiel de \SI{30}{mA} protège les personnes contre les fuites de courant (électrisation) ; le disjoncteur divisionnaire protège les fils du circuit contre les surintensités (échauffement, incendie) ; la prise de terre évacue les courants de défaut vers le sol.
\end{itemize}
\end{aretenir}

\begin{attention}[Pièges fréquents]
\begin{itemize}
\item Ne pas oublier de convertir les watts en kilowatts avant de calculer l'énergie en \si{kWh} : $E(\si{kWh}) = P(\si{kW}) \times t(\si{h})$.
\item Ne pas confondre la puissance (en \si{W}, instantanée) et l'énergie (en \si{kWh}, cumulée dans le temps) : un appareil puissant utilisé peu de temps peut consommer moins qu'un appareil modeste utilisé longtemps.
\item Pour les 4 lampes, multiplier la puissance d'une lampe par 4 avant tout calcul.
\item Le disjoncteur divisionnaire (\SI{16}{A}) protège le \emph{circuit} (les fils dans les murs). La rallonge multiprise a son propre calibre (souvent \SI{10}{A} ou \SI{16}{A} max) : c'est le maillon faible en cas de surcharge concentrée.
\end{itemize}
\end{attention}

\courssub{Démarche et corrigé commenté}

\begin{exoresolu}[Corrigé commenté de l'audit énergétique]
Énoncé : réaliser l'audit complet de la maison Ngo Bell (consommation, classement, sécurité, recommandations).

\tcblower

\textbf{Étape 1 : consommation hebdomadaire mesurée au compteur.}

On applique $E = \text{index final} - \text{index initial}$ :
\begin{align*}
E_{\text{semaine}} &= \num{12534,5} - \num{12450,0} = \SI{84,5}{kWh} \\
E_{\text{jour}} &= \frac{\num{84,5}}{7} \approx \SI{12,1}{kWh}
\end{align*}
La famille consomme environ \SI{12,1}{kWh} par jour en moyenne.

\textbf{Étape 2 : consommation théorique des appareils (par jour).}

On utilise $E = P \times t$ avec $P$ en \si{kW} ($1~\text{kW} = 1000~\text{W}$) :

\begin{center}
\begin{tabularx}{\linewidth}{|l|X|X|}
\hline
\rowcolor{popblueL} \textbf{Appareil} & \textbf{Calcul ($P$ en kW $\times$ $t$ en h)} & \textbf{Énergie / jour} \\
\hline
Réfrigérateur & $0{,}150 \times 8$ & \SI{1,20}{kWh} \\
\hline
Téléviseur & $0{,}100 \times 5$ & \SI{0,50}{kWh} \\
\hline
4 halogènes & $4 \times 0{,}050 \times 6$ & \SI{1,20}{kWh} \\
\hline
Fer à repasser & $1{,}200 \times 0{,}5$ & \SI{0,60}{kWh} \\
\hline
Bouilloire & $2{,}000 \times 0{,}3$ & \SI{0,60}{kWh} \\
\hline
Ventilateur & $0{,}075 \times 8$ & \SI{0,60}{kWh} \\
\hline
\textbf{Total théorique} & & \textbf{\SI{4,70}{kWh}} \\
\hline
\end{tabularx}
\end{center}

\textbf{Comparaison et commentaire :} la consommation théorique (\SI{4,7}{kWh}/jour) est nettement inférieure à la consommation mesurée (\SI{12,1}{kWh}/jour). L'écart d'environ \SI{7,4}{kWh}/jour peut s'expliquer par :
\begin{itemize}
\item des appareils non recensés (congélateur, chauffe-eau électrique, recharge de téléphones, éclairage extérieur, ordinateur) ;
\item des appareils vieillissants (réfrigérateur mal entretenu, joints usés) qui consomment plus que leur puissance nominale ;
\item des durées d'utilisation réelles plus longues que les estimations (télévision allumée en veille, lumières oubliées).
\end{itemize}
C'est une piste d'investigation supplémentaire pour la famille : il faut tout recenser.

\textbf{Étape 3 : classement des appareils par énergie quotidienne.}

Du plus énergivore au moins énergivore :
\begin{enumerate}
\item Réfrigérateur et 4 halogènes (ex æquo, \SI{1,20}{kWh/jour}) ;
\item Fer à repasser, Bouilloire, Ventilateur (ex æquo, \SI{0,60}{kWh/jour}) ;
\item Téléviseur (\SI{0,50}{kWh/jour}).
\end{enumerate}

\textbf{Remarque importante (piège corrigé) :} on accuse souvent la bouilloire (\SI{2000}{W}) parce qu'elle est très puissante. Or elle ne fonctionne que \SI{0,3}{h} par jour : elle consomme deux fois moins que les lampes halogènes, bien moins puissantes (\SI{50}{W} chacune) mais allumées 6 h. C'est bien le produit $P \times t$ qui compte pour l'énergie, pas la puissance seule.

\textbf{Étape 4 : courants et protections.}

Avec $I = \dfrac{P}{U}$ sous \SI{220}{V} :
\begin{align*}
I_{\text{bouilloire}} &= \frac{2000}{220} \approx \SI{9,1}{A} \\
I_{\text{fer}} &= \frac{1200}{220} \approx \SI{5,5}{A} \\
I_{\text{total sur rallonge}} &= 9{,}1 + 5{,}5 = \SI{14,6}{A}
\end{align*}

\begin{itemize}
\item \textbf{Disjoncteur divisionnaire 16 A :} l'intensité totale (\SI{14,6}{A}) reste inférieure au calibre (\SI{16}{A}). Le disjoncteur \textbf{ne déclenche pas} pour surintensité. Le circuit mural est protégé.
\item \textbf{Rallonge multiprise :} la plupart des rallonges domestiques sont calibrées pour \SI{10}{A} maximum (parfois \SI{16}{A} pour les modèles renforcés, mais rarement). Ici, \SI{14,6}{A} traverse la rallonge : elle est \textbf{surchargée}. Elle chauffe anormalement, l'isolant fond, risque d'incendie \textbf{très élevé}.
\item \textbf{Conclusion :} il faut brancher la bouilloire et le fer sur \textbf{deux prises murales distinctes} (sur le même circuit ou circuits différents), et \textbf{jamais} en même temps sur une même rallonge.
\end{itemize}

\textbf{Étape 5 : sécurité des personnes.}

\begin{itemize}
\item \textbf{Différentiel 500 mA : NON CONFORME.} La protection des personnes (contact direct/indirect) exige une sensibilité $\le \SI{30}{mA}$ (norme NF C 15-100). Un différentiel \SI{500}{mA} ne protège que contre les risques d'incendie (fuites de courant importantes), pas contre l'électrisation mortelle. Au-delà de \SI{30}{mA} traversant le corps humain, il y a danger de fibrillation ventriculaire ; un différentiel de \SI{500}{mA} déclenche trop tard (ou pas du tout pour un défaut de faible intensité).
\item \textbf{Fils dénudés sur la rallonge :} exposition au contact direct avec des parties sous tension (\SI{220}{V}). Risque immédiat d'électrocution (surtout avec mains mouillées en cuisine) et de court-circuit (arc électrique, incendie).
\item \textbf{Absence de prise de terre raccordée :} si la carcasse métallique d'un appareil (fer, réfrigérateur, bouilloire) est mise accidentellement sous tension (fil dénudé touchant le châssis), le courant de défaut ne peut pas s'écouler vers la terre. La carcasse reste sous tension : toute personne la touchant est électrisée. Le différentiel (même 30 mA) ne voit pas le défaut si le courant ne fuit pas vers la terre.
\end{itemize}

\textbf{Étape 6 : recommandations chiffrées.}

\emph{Recommandation 1 : remplacer les 4 lampes halogènes par des lampes LED.}
\begin{align*}
P_{\text{avant}} &= 4 \times 50 = \SI{200}{W} = \SI{0,200}{kW} \\
P_{\text{après}} &= 4 \times 7 = \SI{28}{W} = \SI{0,028}{kW} \\
\Delta P &= 0{,}200 - 0{,}028 = \SI{0,172}{kW} \\
E_{\text{économisée/mois}} &= 0{,}172 \times 6 \times 30 = \SI{30,96}{kWh} \approx \SI{31}{kWh} \\
\text{Économie FCFA} &= 31 \times 75 = \SI{2325}{FCFA}
\end{align*}

\emph{Recommandation 2 : réduire l'éclairage inutile (éteindre en quittant une pièce).}
Gain estimé : 1 h de moins par jour pour les 4 lampes.
\begin{itemize}
\item Avec les LED (après recommandation 1) : $0{,}028 \times 1 \times 30 = \SI{0,84}{kWh} \approx \SI{63}{FCFA}$.
\item Avec les halogènes (avant) : $0{,}200 \times 1 \times 30 = \SI{6}{kWh} \approx \SI{450}{FCFA}$.
\end{itemize}
Le geste simple rapporte plus tant que les halogènes sont en place.

\emph{Recommandation 3 : remettre l'installation en conformité sécurité (priorité absolue).}
\begin{itemize}
\item Remplacer le différentiel 500 mA par un différentiel \SI{30}{mA} (coût : \SIrange{5000}{10000}{FCFA} environ).
\item Raccorder la prise de terre de la maison (piquet de terre + liaison équipotentielle).
\item Remplacer la rallonge dénudée par une rallonge neuve calibre \SI{16}{A} (ou mieux : faire tirer une prise dédiée pour la bouilloire).
\item Répartir les gros appareils sur des prises distinctes.
\end{itemize}
Coût unique, mais protection durable des personnes et des biens (valeur inestimable).

\emph{Recommandation 4 (bonus) : entretenir le réfrigérateur.}
Dégivrer régulièrement (givre = isolant = surconsommation), vérifier les joints de porte, laisser un espace de \SI{5}{cm} derrière pour la ventilation. Un réfrigérateur givré peut consommer jusqu'à 30 \% de plus.
\begin{align*}
\text{Gaspillage estimé} &= 0{,}30 \times \SI{1,20}{kWh/jour} \times 30 = \SI{10,8}{kWh/mois} \approx \SI{11}{kWh} \\
\text{Économie} &= 11 \times 75 = \SI{825}{FCFA}
\end{align*}

\textbf{Étape 7 : affiche et restitution orale.}
L'affiche comporte trois rubriques claires :
\begin{enumerate}
\item \textbf{Combien consommons-nous ?} Tableau comparatif (compteur vs théorie), graphique en barres du classement.
\item \textbf{Sommes-nous en sécurité ?} Schéma du tableau électrique annoté avec les 3 défauts majeurs (Différentiel 500 mA, Fils dénudés, Pas de terre).
\item \textbf{Que faisons-nous ?} Les 3-4 recommandations chiffrées (kWh et FCFA).
\end{enumerate}

\end{exoresolu}

\courssub{Bilan de l'activité}

Cet audit énergétique a permis de mettre en évidence plusieurs idées essentielles de la leçon :

\begin{itemize}
\item \textbf{L'énergie se mesure et se calcule.} Le compteur électrique est l'instrument de référence : la différence de deux index donne l'énergie consommée, que l'on retrouve (au moins partiellement) par le calcul $E = P \times t$ appliqué à chaque appareil.
\item \textbf{Puissance et énergie sont deux grandeurs différentes.} L'appareil le plus puissant (la bouilloire, \SI{2000}{W}) n'est pas le plus énergivore : la durée d'utilisation est déterminante. C'est une erreur fréquente que l'activité corrige durablement.
\item \textbf{La relation $P = U \times I$ protège.} Le calcul des intensités permet de vérifier qu'un circuit, une rallonge ou un disjoncteur n'est pas surchargé : c'est une démarche de prévention concrète contre l'incendie.
\item \textbf{La sécurité électrique repose sur des règles précises, non négociables :} différentiel de \SI{30}{mA} (protection vie), prise de terre raccordée (protection défaut), conducteurs en bon état (pas de contact direct), pas de surconcentration d'appareils sur une multiprise (protection incendie).
\item \textbf{Les économies d'énergie se chiffrent.} Remplacer des halogènes par des LED, éteindre les appareils inutiles, entretenir le réfrigérateur : chaque action se traduit en \si{kWh} puis en francs CFA, ce qui rend la physique utile à la vie quotidienne de la famille.
\end{itemize}

\begin{savaistu}[Pour aller plus loin : la facture ENEO réelle]
Au Cameroun, la facture ENEO est calculée par \textbf{tranches de consommation} (tarification progressive) : plus tu consommes, plus le prix du \si{kWh} augmente (ex. : tranche 1 \textasciitilde 50 FCFA, tranche 2 \textasciitilde 75 FCFA, tranche 3 \textasciitilde 100 FCFA...). Réduire sa consommation permet donc de rester dans les tranches les moins chères : l'économie réelle est souvent \textbf{supérieure} à celle calculée avec le tarif moyen de \SI{75}{FCFA/kWh} utilisé ici. De plus, chaque \si{kWh} économisé réduit la demande sur le réseau national et limite les délestages !
\end{savaistu}

\begin{exercice}[Pour t'entraîner à la maison]
Releve l'index du compteur de ta famille à une semaine d'intervalle (même jour, même heure), calcule la consommation hebdomadaire, puis identifie les trois appareils les plus énergivores de la maison en appliquant $E = P \times t$ (puissance sur la plaque signalétique $\times$ durée d'utilisation estimée). Propose une recommandation chiffrée à tes parents.
\end{exercice}

\newpage

\coursec{Méthodes et exercices résolus}

\coursec{Leçon 14 : Énergie électrique et puissance électrique}

\courssub{Grandeurs électriques et instruments de mesure}

\begin{definition}[Grandeurs électriques de base]
Dans un circuit électrique, on manipule cinq grandeurs fondamentales :
\begin{itemize}
\item la \cle{tension} $U$ (voltage), en \cle{volt} (\SI{}{V}) : elle mesure la « pression » électrique fournie par le générateur ;
\item l'\cle{intensité} $I$ du courant, en \cle{ampère} (\SI{}{A}) : elle mesure la quantité d'électricité qui traverse une section du fil par seconde ;
\item la \cle{résistance} $R$, en \cle{ohm} (\SI{}{\ohm}) : elle caractérise l'opposition d'un dipôle au passage du courant ;
\item la \cle{puissance} $P$, en \cle{watt} (\SI{}{W}) : elle mesure l'énergie échangée par seconde ;
\item l'\cle{énergie} $E$, en \cle{joule} (\SI{}{J}) ou en \cle{kilowattheure} (\SI{}{kWh}) : elle mesure la quantité totale d'énergie transférée.
\end{itemize}
\end{definition}

\begin{propriete}[Instruments de mesure et leurs branchements]
\begin{center}
\begin{tabularx}{\linewidth}{|l|X|X|l|}\hline
\rowcolor{popblueL}
\textbf{Grandeur} & \textbf{Instrument} & \textbf{Branchement} & \textbf{Unité} \\ \hline
Tension $U$ & Voltmètre & En \cle{dérivation} (aux bornes du dipôle) & Volt (V) \\ \hline
Intensité $I$ & Ampèremètre & En \cle{série} (sur le trajet du courant) & Ampère (A) \\ \hline
Puissance $P$ & Wattmètre & 4 bornes : 2 en série (courant), 2 en dérivation (tension) & Watt (W) \\ \hline
Énergie $E$ & Compteur électrique (ENEO) & En série sur l'arrivée générale & Kilowattheure (kWh) \\ \hline
\end{tabularx}
\end{center}
\textbf{Rappel sécurité} : un ampèremètre a une résistance quasi nulle ; branché en dérivation, il crée un \cle{court-circuit} dangereux. Un voltmètre a une résistance très grande ; branché en série, il bloque le courant.
\end{propriete}

\begin{methode}[Choisir et brancher le bon instrument]
Pour répondre à une question sur la mesure d'une grandeur électrique, applique la démarche suivante :
\begin{enumerate}
\item \textbf{Repère la grandeur à mesurer} : tension $U$ (volt, V), intensité $I$ (ampère, A), puissance $P$ (watt, W) ou énergie $E$ (kilowattheure, kWh).
\item \textbf{Choisis l'instrument} : voltmètre pour $U$, ampèremètre pour $I$, wattmètre pour $P$, compteur électrique pour $E$.
\item \textbf{Place l'instrument correctement} : l'ampèremètre se branche en \cle{série} dans le circuit ; le voltmètre se branche en \cle{dérivation} (aux bornes du dipôle).
\item \textbf{Respecte les bornes} : la borne COM est reliée au pôle négatif du générateur ; le courant entre par la borne marquée A (ou V).
\item \textbf{Choisis le calibre} : commence toujours par le calibre le plus élevé, puis diminue pour gagner en précision.
\end{enumerate}
\textbf{Réflexe} : « voltmètre = en V comme en déri\textbf{V}ation ; ampèremètre = en série comme un \textbf{A}gent sur la route du courant ».
\end{methode}

\begin{experience}[Mesure de la tension et de l'intensité dans un circuit simple]
\textbf{Matériel} : générateur réglable (0--12 V), lampe 6 V, rhéostat, ampèremètre, voltmètre, fils de connexion.\\
\textbf{Protocole} :
\begin{enumerate}
\item Monte le circuit : générateur -- ampèremètre -- rhéostat -- lampe -- retour générateur.
\item Branche le voltmètre en dérivation aux bornes de la lampe.
\item Varie le rhéostat pour faire varier la luminosité de la lampe.
\item Releve 5 couples de valeurs $(U ; I)$.
\end{enumerate}
\textbf{Observations} : quand $U$ augmente, $I$ augmente et la lampe brille plus fort.\\
\textbf{Interprétation} : la puissance lumineuse et thermique de la lampe dépend de $U$ et $I$. On vérifiera plus loin que $P = U \times I$.
\legende{Schéma de montage : ampèremètre (A) en série, voltmètre (V) en dérivation aux bornes de la lampe (L).}
\begin{popfigure}
\begin{tikzpicture}[scale=0.9, every node/.style={font=\small}]
% Generator
\draw[thick] (0,0) -- (0,2);
\draw[thick] (0,2) -- (1,2);
\draw (0.5,2) node[above] {G};
\draw[thick] (1,2) -- (2,2);
% Ammeter
\draw[thick] (2,2) -- (2.5,2);
\draw (2.75,2) circle (0.4);
\draw (2.75,2) node {A};
\draw[thick] (3.15,2) -- (4,2);
% Rheostat
\draw[thick] (4,2) -- (4.5,2);
\draw (4.75,2) node[draw, rectangle, minimum width=1cm, minimum height=0.6cm] {Rh};
\draw[thick] (5.25,2) -- (6,2);
% Lamp
\draw[thick] (6,2) -- (6.5,2);
\draw (6.75,2) circle (0.4);
\draw (6.75,2) node {L};
\draw[thick] (7.15,2) -- (8,2);
% Return wire
\draw[thick] (8,2) -- (8,0);
\draw[thick] (8,0) -- (0,0);
% Voltmeter
\draw[thick] (6,2) -- (6,3);
\draw (6,3.5) circle (0.4);
\draw (6,3.5) node {V};
\draw[thick] (6,4) -- (7.15,4);
\draw[thick] (7.15,4) -- (7.15,2);
% Labels
\draw[popblue, <->] (2.75,1.4) -- (3.15,1.4) node[midway, below] {Série};
\draw[poporange, <->] (6,3.5) -- (7.15,3.5) node[midway, above] {Dérivation};
\end{tikzpicture}
\legende{Montage de mesure : l'ampèremètre (A) est en série, le voltmètre (V) est en dérivation aux bornes de la lampe (L).}
\end{popfigure}
\end{experience}

\begin{exoresolu}[Branchement des appareils de mesure]
Un élève de 3ème veut mesurer la tension aux bornes d'une lampe et l'intensité du courant qui la traverse dans un circuit alimenté par un générateur de \SI{6}{V}.\\
1) Nomme les deux instruments nécessaires.\\
2) Indique comment chacun doit être branché.\\
3) L'élève branche le voltmètre en série : que se passe-t-il ?
\tcblower
\textbf{1) Instruments nécessaires}\\
Pour mesurer la tension $U$ aux bornes de la lampe, on utilise un \cle{voltmètre}. Pour mesurer l'intensité $I$ qui traverse la lampe, on utilise un \cle{ampèremètre}.\\
\textbf{2) Branchements}\\
Le voltmètre se branche \textbf{en dérivation} aux bornes de la lampe : ses deux fils sont reliés directement aux deux bornes de la lampe, sans couper le circuit.\\
L'ampèremètre se branche \textbf{en série} : on coupe le circuit et on l'insère sur le trajet du courant, par exemple entre le générateur et la lampe.\\
\textbf{3) Voltmètre branché en série}\\
Un voltmètre possède une très grande résistance interne. Branché en série, il empêche pratiquement le passage du courant : l'intensité devient quasi nulle et \textbf{la lampe ne brille plus}. Le voltmètre indique alors environ la tension du générateur (\SI{6}{V}), mais la mesure voulue n'est pas réalisée.\\
\textbf{Conclusion} : le voltmètre se branche en dérivation, l'ampèremètre en série ; un mauvais branchement fausse la mesure et peut empêcher le circuit de fonctionner.
\end{exoresolu}

\courssub{Puissance électrique d'un résistor}

\begin{propriete}[Expression de la puissance électrique reçue par un résistor]
Pour un dipôle ohmique (résistor) traversé par un courant d'intensité $I$ et aux bornes duquel la tension est $U$, la \cle{puissance électrique} $P$ reçue s'exprime de trois manières équivalentes :
\[
P = U \cdot I \qquad P = R \cdot I^{2} \qquad P = \frac{U^{2}}{R}
\]
\begin{itemize}
\item $P$ en watts (W), $U$ en volts (V), $I$ en ampères (A), $R$ en ohms ($\Omega$).
\item La démonstration de $P = R \cdot I^{2}$ à partir de $P = U \cdot I$ et de la loi d'Ohm $U = R \cdot I$ est \textbf{exigible au BEPC}.
\end{itemize}
\end{propriete}

\begin{methode}[Utiliser P = U.I = R.I²]
\begin{enumerate}
\item \textbf{Fais l'inventaire des données} : note $U$, $I$, $R$ avec leurs unités. Convertis si nécessaire (mA en A : diviser par 1000 ; kW en W : multiplier par 1000).
\item \textbf{Choisis la formule adaptée} :
\begin{itemize}
\item si tu connais $U$ et $I$ : $P = U \cdot I$ ;
\item si tu connais $R$ et $I$ : $P = R \cdot I^{2}$ ;
\item si tu connais $U$ et $R$ : $P = \dfrac{U^{2}}{R}$.
\end{itemize}
\item \textbf{Calcule} en respectant les unités : $U$ en volts, $I$ en ampères, $R$ en ohms, alors $P$ s'obtient en \textbf{watts (W)}.
\item \textbf{Vérifie l'ordre de grandeur} : une lampe de poche : quelques watts ; un fer à repasser : environ \SI{1000}{W} ; un climatiseur : \SI{1500}{W} à \SI{3000}{W}.
\item \textbf{Rédige la conclusion} avec la valeur et l'unité.
\end{enumerate}
\textbf{Astuce} : si on te demande $I$ à partir de $P$ et $U$, utilise $I = \dfrac{P}{U}$ ; si on te demande $R$ à partir de $P$ et $I$, utilise $R = \dfrac{P}{I^{2}}$.
\end{methode}

\begin{exoresolu}[Puissance d'un fer à repasser]
Un fer à repasser porte les indications « \SI{220}{V} ; \SI{1000}{W} ». Il est branché sur le réseau ENEO de tension $U = \SI{220}{V}$.\\
1) Calcule l'intensité $I$ du courant qui le traverse.\\
2) Calcule la résistance $R$ de sa résistance chauffante.\\
3) Retrouve la puissance en utilisant la formule $P = R \cdot I^{2}$.
\tcblower
\textbf{1) Calcul de l'intensité}\\
On connaît $P = \SI{1000}{W}$ et $U = \SI{220}{V}$. De $P = U \cdot I$, on tire :
\[ I = \frac{P}{U} = \frac{1000}{220} \approx \SI{4,55}{A}. \]
\textbf{2) Calcul de la résistance}\\
D'après la loi d'Ohm, $U = R \cdot I$, donc :
\[ R = \frac{U}{I} = \frac{220}{4,55} \approx \SI{48,4}{\ohm}. \]
\textbf{3) Vérification avec $P = R \cdot I^{2}$}\\
\[ P = R \cdot I^{2} = 48,4 \times (4,55)^{2} = 48,4 \times 20,7025 \approx \SI{1002}{W} \approx \SI{1000}{W}. \]
On retrouve bien la puissance indiquée sur la plaque signalétique du fer.\\
\textbf{Conclusion} : le fer à repasser est traversé par un courant d'environ \SI{4,55}{A}, sa résistance vaut environ \SI{48,4}{\ohm}, et les trois formules $P = U \cdot I$, $P = R \cdot I^{2}$ et $P = \dfrac{U^{2}}{R}$ donnent le même résultat.
\end{exoresolu}

\courssub{Mesure de la puissance et de l'énergie consommée}

\begin{definition}[Énergie électrique]
L'\cle{énergie électrique} $E$ consommée (ou fournie) par un appareil de puissance $P$ pendant une durée $t$ est :
\[ E = P \cdot t \]
L'unité légale pour la facturation (compteur ENEO) est le \cle{kilowattheure} (\SI{}{kWh}).
\begin{center}
\fcolorbox{popdark}{popcream}{\parbox{0.9\linewidth}{\centering
\textbf{Conversions utiles} :\\
$1 \text{ kW} = 1000 \text{ W}$ \quad $1 \text{ h} = 60 \text{ min} = 3600 \text{ s}$\\
$1 \text{ kWh} = 1000 \text{ W} \times 3600 \text{ s} = \SI{3,6 \times 10^{6}}{J} = \SI{3,6}{MJ}$
}}
\end{center}
\end{definition}

\begin{methode}[Utiliser E = P.t pour l'énergie en kWh]
\begin{enumerate}
\item \textbf{Identifie} la puissance $P$ de l'appareil (en W ou kW) et la durée $t$ de fonctionnement (en h).
\item \textbf{Convertis} : pour obtenir l'énergie en kilowattheures (kWh), exprime $P$ en \textbf{kilowatts} (diviser par 1000) et $t$ en \textbf{heures}.
\item \textbf{Applique la formule} : $E = P \cdot t$, avec $E$ en kWh.
\item \textbf{Calcule le coût} si demandé : coût = $E$ (kWh) $\times$ prix du kWh (en FCFA).
\item \textbf{Conclus} par une phrase complète avec l'unité.
\end{enumerate}
\textbf{Réflexe} : le kWh est une unité d'\textbf{énergie}, pas de puissance ! C'est l'unité du compteur ENEO.
\end{methode}

\begin{exoresolu}[Facture ENEO d'une famille]
Une famille de Yaoundé utilise chaque jour : une télévision de \SI{100}{W} pendant 5 h, un ventilateur de \SI{60}{W} pendant 8 h et 5 lampes de \SI{10}{W} pendant 6 h.\\
1) Calcule l'énergie totale consommée par jour en kWh.\\
2) Calcule l'énergie consommée en 30 jours.\\
3) Le kWh coûte 99 FCFA : calcule la dépense mensuelle.
\tcblower
\textbf{1) Énergie quotidienne}\\
On applique $E = P \cdot t$ pour chaque appareil, avec $P$ en kW :
\begin{align*}
E_{\text{TV}} &= 0,100 \times 5 = \SI{0,500}{kWh} \\
E_{\text{ventilateur}} &= 0,060 \times 8 = \SI{0,480}{kWh} \\
E_{\text{lampes}} &= 5 \times 0,010 \times 6 = \SI{0,300}{kWh}
\end{align*}
Énergie totale par jour :
\[ E_{\text{jour}} = 0,500 + 0,480 + 0,300 = \SI{1,28}{kWh}. \]
\textbf{2) Énergie mensuelle}\\
\[ E_{\text{mois}} = 1,28 \times 30 = \SI{38,4}{kWh}. \]
\textbf{3) Dépense mensuelle}\\
\[ \text{Coût} = 38,4 \times 99 = \SI{3801,6}{FCFA} \approx \SI{3802}{FCFA}. \]
\textbf{Conclusion} : cette famille consomme \SI{38,4}{kWh} par mois, ce qui représente une dépense d'environ 3 800 FCFA. Éteindre les appareils inutiles permet de réduire la facture.
\end{exoresolu}

\begin{methode}[Exploiter les indications d'un compteur électrique]
\begin{enumerate}
\item \textbf{Repère les deux relevés} : index du compteur au début et à la fin de la période (en kWh).
\item \textbf{Calcule la consommation} : $E = \text{index final} - \text{index initial}$.
\item \textbf{Vérifie la puissance souscrite} : la somme des puissances des appareils fonctionnant \textbf{en même temps} ne doit pas dépasser la puissance maximale du compteur (ex. \SI{3300}{W}, \SI{6600}{W}), sinon le disjoncteur \textbf{coupe} le courant (disjonction).
\item \textbf{Conclus} en proposant éventuellement une solution (éteindre un appareil, décaler les usages, souscrire une puissance plus élevée).
\end{enumerate}
\end{methode}

\begin{exoresolu}[Disjonction à la maison]
Le compteur d'une maison autorise une puissance maximale de \SI{3300}{W}. Un soir, fonctionnent simultanément : un climatiseur de \SI{1500}{W}, une bouilloire de \SI{1200}{W}, un fer à repasser de \SI{1000}{W} et des lampes totalisant \SI{200}{W}.\\
1) Calcule la puissance totale utilisée.\\
2) Le disjoncteur va-t-il couper le courant ? Justifie.\\
3) Propose une solution sans changer d'abonnement.
\tcblower
\textbf{1) Puissance totale}\\
Les appareils branchés simultanément ajoutent leurs puissances :
\[ P_{\text{totale}} = 1500 + 1200 + 1000 + 200 = \SI{3900}{W}. \]
\textbf{2) Comparaison à la puissance maximale}\\
On compare : $P_{\text{totale}} = \SI{3900}{W} > P_{\text{max}} = \SI{3300}{W}$.\\
La puissance demandée dépasse la puissance souscrite : \textbf{le disjoncteur coupe le courant} pour protéger l'installation. C'est la « disjonction ».\\
\textbf{3) Solution}\\
Il faut réduire la puissance simultanée d'au moins $3900 - 3300 = \SI{600}{W}$. Par exemple, \textbf{ne pas utiliser le fer à repasser} (\SI{1000}{W}) en même temps que les autres appareils : la puissance devient $3900 - 1000 = \SI{2900}{W} < \SI{3300}{W}$.\\
\textbf{Conclusion} : la puissance totale de \SI{3900}{W} dépasse la limite de \SI{3300}{W}, donc le disjoncteur saute ; en repassant à un autre moment, la consommation simultanée redescend à \SI{2900}{W} et l'installation fonctionne normalement.
\end{exoresolu}

\courssub{Sécurité électrique et consignes de prévention}

\begin{aretenir}[Risques électriques principaux et prévention]
\begin{center}
\begin{tabularx}{\linewidth}{|l|X|X|}\hline
\rowcolor{poporangeL}
\textbf{Danger} & \textbf{Origine physique} & \textbf{Prévention} \\ \hline
\cle{Électrisation} / \cle{Électrocution} & Passage du courant dans le corps (résistance corps $\approx \SI{1000}{\ohm}$ sec, $\SI{300}{\ohm}$ mouillé). $I = U/R$. & Ne pas toucher un appareil avec les mains mouillées ; ne pas insérer d'objet dans une prise ; mise à la terre (prise 3 broches) + disjoncteur différentiel (30 mA). \\ \hline
\cle{Court-circuit} & Contact direct entre les deux bornes du générateur ($R \approx 0 \Rightarrow I$ énorme). & Fusibles, disjoncteurs magnétothermiques ; isoler les fils dénudés. \\ \hline
\cle{Surchauffe} / \cle{Incendie} & Effet Joule dans les fils : $P = R_{\text{fil}} \cdot I^{2}$. Si $I$ trop grand (surcharge, rallonge enroulée), $T$ monte, isolant fond. & Ne pas brancher plusieurs appareils puissants sur une même rallonge/multiprise ; dérouler totalement les rallonges ; respecter la section des fils (norme NF C 15-100). \\ \hline
\end{tabularx}
\end{center}
\end{aretenir}

\begin{methode}[Analyser une situation de danger]
Face à une situation concrète (rallonge, prise, appareil mouillé...), raisonne ainsi :
\begin{enumerate}
\item \textbf{Repère le danger} : surcharge d'une multiprise, fil dénudé, appareil près de l'eau, appareil de puissance trop élevée pour le circuit.
\item \textbf{Nomme le risque} : échauffement des fils, incendie, électrisation (passage du courant dans le corps), court-circuit.
\item \textbf{Cite la règle de sécurité} : ne pas brancher plusieurs appareils puissants sur une même rallonge ; ne jamais toucher un appareil électrique avec les mains mouillées ; couper le courant (disjoncteur général) avant toute intervention ; utiliser des prises avec terre.
\item \textbf{Justifie avec la physique} : si $I$ augmente trop dans un fil, l'effet Joule ($P = R \cdot I^{2}$) échauffe le fil et peut provoquer un incendie. Si le corps est en contact avec la tension secteur (\SI{220}{V}), le courant $I = U/R_{\text{corps}}$ peut être mortel ($> \SI{30}{mA}$).
\end{enumerate}
\end{methode}

\begin{exoresolu}[La rallonge surchargée]
Pendant les fêtes, un père de famille branche sur une même rallonge : un réfrigérateur (\SI{200}{W}), un four (\SI{2000}{W}) et une bouilloire (\SI{1200}{W}). La rallonge supporte au maximum \SI{16}{A} sous \SI{220}{V}.\\
1) Calcule l'intensité totale du courant dans la rallonge.\\
2) Explique le danger encouru en utilisant l'effet Joule.\\
3) Donne deux conseils de sécurité.
\tcblower
\textbf{1) Intensité totale}\\
Puissance totale : $P = 200 + 2000 + 1200 = \SI{3400}{W}$.\\
De $P = U \cdot I$, on tire :
\[ I = \frac{P}{U} = \frac{3400}{220} \approx \SI{15,5}{A}. \]
L'intensité est proche de la limite de \SI{16}{A} : la rallonge fonctionne à sa capacité maximale, sans aucune marge de sécurité.\\
\textbf{2) Danger}\\
D'après l'effet Joule, la puissance dégagée en chaleur dans un fil est $P = R \cdot I^{2}$ : elle augmente avec le \textbf{carré} de l'intensité. Si l'intensité dépasse la limite, les fils de la rallonge s'échauffent fortement, l'isolant fond et un \textbf{incendie} peut se déclarer. De plus, si la rallonge est enroulée, la chaleur ne s'évacue pas (effet bobine).\\
\textbf{3) Conseils de sécurité}\\
\begin{itemize}
\item Brancher le four sur une prise murale séparée, jamais sur une rallonge avec d'autres appareils puissants.
\item Ne jamais laisser la rallonge enroulée pendant l'utilisation et vérifier que les fils ne sont pas dénudés ou écrasés par une porte.
\end{itemize}
\textbf{Conclusion} : avec $I \approx \SI{15,5}{A}$, la rallonge est au bord de la surcharge ; il faut répartir les appareils puissants sur des prises différentes pour éviter l'échauffement et l'incendie.
\end{exoresolu}

\courssub{Applications concrètes et résolution de problèmes}

\begin{methode}[Problème de synthèse : démarche en 5 étapes]
\begin{enumerate}
\item \textbf{Lis l'énoncé deux fois} et souligne les données : tension, puissance, durée, prix du kWh, rendement éventuel.
\item \textbf{Écris les formules utiles} avant tout calcul : $P = U \cdot I$, $E = P \cdot t$, coût $= E \times$ prix, $P = R \cdot I^{2}$.
\item \textbf{Convertis toutes les unités} au début : W en kW, minutes en heures, mA en A.
\item \textbf{Enchaîne les calculs} en justifiant chaque étape par la formule utilisée.
\item \textbf{Termine par une phrase de conclusion} qui répond exactement à la question posée, avec l'unité.
\end{enumerate}
\end{methode}

\begin{exoresolu}[Le groupe électrogène du quartier]
Pendant une coupure ENEO, un commerçant fait fonctionner son groupe électrogène. Il alimente un congélateur de puissance $P = \SI{300}{W}$ sous $U = \SI{220}{V}$ pendant $t = 4$ h.\\
1) Calcule l'intensité du courant dans le congélateur.\\
2) Calcule l'énergie consommée en kWh.\\
3) Le groupe consomme 1 litre d'essence pour produire \SI{3}{kWh}. Calcule le volume d'essence utilisé, puis sa dépense si le litre coûte 730 FCFA.
\tcblower
\textbf{1) Intensité du courant}\\
De $P = U \cdot I$, on tire :
\[ I = \frac{P}{U} = \frac{300}{220} \approx \SI{1,36}{A}. \]
\textbf{2) Énergie consommée}\\
On convertit la puissance en kW : $P = \SI{0,300}{kW}$. Alors :
\[ E = P \cdot t = 0,300 \times 4 = \SI{1,2}{kWh}. \]
\textbf{3) Essence utilisée et dépense}\\
Volume d'essence :
\[ V = \frac{E}{3} = \frac{1,2}{3} = \SI{0,4}{L}. \]
Dépense :
\[ \text{Coût} = 0,4 \times 730 = \SI{292}{FCFA}. \]
\textbf{Conclusion} : le congélateur est traversé par un courant d'environ \SI{1,36}{A} ; il consomme \SI{1,2}{kWh} en 4 h, ce qui nécessite \SI{0,4}{L} d'essence, soit une dépense de 292 FCFA pour la soirée de coupure.
\end{exoresolu}

\begin{savaistu}[Le réseau ENEO et la production au Cameroun]
Au Cameroun, l'électricité distribuée par ENEO est principalement d'origine \cle{hydroélectrique} (barrages de Song Loulou, Edea, Memve'ele) et \cle{thermique} (centrales à fioul/gaz de Kribi, Limbe, Dibamba). La tension secteur est de \SI{220}{V} (monophasé) ou \SI{380}{V} (triphasé) à une fréquence de \SI{50}{Hz}. Le compteur Linky (électronique) remplace peu à peu les compteurs à disque : il relève l'index (kWh) à distance et permet de gérer la puissance souscrite (ex. 3,3 kVA, 6,6 kVA, 12 kVA). Le \cle{kilovoltampère (kVA)} est l'unité de la puissance apparente ; pour les charges résistives (fer, bouilloire), \SI{1}{kVA} $\approx$ \SI{1}{kW}. Pour les moteurs (ventilateur, pompe, climatiseur), la puissance active (W) est inférieure à la puissance apparente (VA) à cause du \cle{facteur de puissance} ($\cos \varphi$), notion approfondie en Terminale.
\end{savaistu}

\begin{attention}[Les 5 erreurs qui coûtent des points au BEPC]
\begin{enumerate}
\item \textbf{Confondre puissance et énergie} : la puissance s'exprime en watts (W), l'énergie en kilowattheures (kWh) ou en joules (J). Écrire « l'énergie consommée est de \SI{100}{W} » est une erreur d'unité éliminatoire.
\item \textbf{Oublier les conversions avant $E = P \cdot t$} : pour obtenir des kWh, la puissance doit être en \textbf{kW} et le temps en \textbf{heures}. Avec $P$ en W et $t$ en minutes, le résultat est faux : convertis d'abord (30 min = 0,5 h ; \SI{300}{W} = \SI{0,3}{kW}).
\item \textbf{Mal brancher les instruments dans un schéma} : l'ampèremètre en dérivation provoque un court-circuit ; le voltmètre en série bloque le courant. Retiens : ampèremètre en \textbf{série}, voltmètre en \textbf{dérivation}.
\item \textbf{Appliquer $P = R \cdot I^{2}$ avec $I$ en mA} : toutes les formules exigent $I$ en \textbf{ampères}. $I = \SI{50}{mA} = \SI{0,05}{A}$. Une intensité non convertie donne une puissance un million de fois trop grande !
\item \textbf{Oublier la phrase de conclusion et la justification} : au BEPC, un résultat sans formule citée, sans calcul détaillé ou sans unité perd des points. Rédige toujours : formule $\rightarrow$ application numérique $\rightarrow$ résultat avec unité $\rightarrow$ phrase de conclusion.
\end{enumerate}
\end{attention}

\begin{exercice}[Entraînement BEPC - Niveau 1]
\textbf{Exercice 1 :} Une lampe de chevet porte la mention « \SI{220}{V} - \SI{40}{W} ».
\begin{enumerate}
\item Calcule l'intensité du courant qui traverse cette lampe lorsqu'elle fonctionne normalement.
\item Calcule la résistance de son filament à chaud.
\item Quelle énergie en kWh consomme-t-elle si elle reste allumée 5 heures ?
\end{enumerate}

\textbf{Exercice 2 :} Un chauffe-eau de résistance $R = \SI{22}{\ohm}$ est branché sur le secteur \SI{220}{V}.
\begin{enumerate}
\item Calcule la puissance absorbée par le chauffe-eau.
\item Calcule l'intensité du courant.
\item Si le chauffe-eau fonctionne 2 h par jour pendant 30 jours, quelle est l'énergie consommée en kWh ? (Prix du kWh = 99 FCFA). Combien coûte ce mois de fonctionnement ?
\end{enumerate}
\end{exercice}

\begin{corrige}[Exercice 1]
\begin{enumerate}
\item $I = \frac{P}{U} = \frac{40}{220} \approx \SI{0,182}{A}$.
\item $R = \frac{U}{I} = \frac{220}{0,182} \approx \SI{1209}{\ohm}$ (ou $R = \frac{U^2}{P} = \frac{220^2}{40} = \SI{1210}{\ohm}$).
\item $E = P \cdot t = 0,040 \times 5 = \SI{0,20}{kWh}$.
\end{enumerate}
\end{corrige}

\begin{corrige}[Exercice 2]
\begin{enumerate}
\item $P = \frac{U^2}{R} = \frac{220^2}{22} = \frac{48400}{22} = \SI{2200}{W} = \SI{2,2}{kW}$.
\item $I = \frac{P}{U} = \frac{2200}{220} = \SI{10}{A}$ (ou $I = \frac{U}{R} = \frac{220}{22} = \SI{10}{A}$).
\item $E = 2,2 \times 2 \times 30 = \SI{132}{kWh}$. Coût $= 132 \times 99 = \SI{13068}{FCFA}$.
\end{enumerate}
\end{corrige}

\begin{exercice}[Entraînement BEPC - Niveau 2 : Synthèse]
Un artisan soudeur utilise un poste à souder branché sur une prise triphasée \SI{380}{V} (puissance souscrite \SI{12}{kVA}). Le poste indique une puissance absorbée de \SI{8}{kW} sous \SI{380}{V}. Sur le même circuit, il branche une meuleuse de \SI{2000}{W} et un éclairage de \SI{500}{W}.
\begin{enumerate}
\item Calcule l'intensité appelée par le poste à souder seul.
\item Calcule la puissance totale si les trois appareils fonctionnent ensemble.
\item Le disjoncteur du compteur (réglé sur \SI{12}{kVA} $\approx$ \SI{12}{kW} pour charges résistives) va-t-il disjoncter ? Justifie.
\item Propose une organisation du travail pour ne pas faire disjoncter.
\end{enumerate}
\end{exercice}

\begin{corrige}[Exercice Niveau 2]
\begin{enumerate}
\item Poste à souder : $P = \SI{8000}{W}$, $U = \SI{380}{V}$. $I = \frac{P}{U} = \frac{8000}{380} \approx \SI{21,1}{A}$ (par phase, en régime triphasé équilibré $P = \sqrt{3} U I \cos\varphi$, mais en 3ème on simplifie souvent en monophasé équivalent ou on donne $I = P/U$ pour l'ordre de grandeur. \textbf{Note} : Au BEPC 3ème, on reste généralement en monophasé \SI{220}{V}. Si l'exercice impose triphasé, on précise : $I = \frac{P}{\sqrt{3} U} \approx \SI{12,1}{A}$ par phase. Ici, supposons une simplification monophasée à \SI{380}{V} pour le calcul d'intensité globale ou on précise le contexte). \textbf{Correction adaptée 3ème} : On considère souvent la puissance par phase ou on simplifie. Disons $I \approx \SI{21}{A}$ (courant dans un conducteur si modélisé simple).
\item $P_{\text{totale}} = 8000 + 2000 + 500 = \SI{10500}{W} = \SI{10,5}{kW}$.
\item $P_{\text{totale}} = \SI{10,5}{kW} < \SI{12}{kW}$ (puissance souscrite). \textbf{Non}, le disjoncteur ne devrait pas couper pour raison de puissance active. (Attention : en triphasé, il faut vérifier l'équilibre des phases).
\item Si d'autres appareils s'ajoutent (ex. compresseur \SI{2}{kW}), le total dépasse \SI{12}{kW}. Solution : ne pas utiliser la meuleuse en même temps que le poste à souder à pleine puissance, ou demander une augmentation de puissance souscrite.
\end{enumerate}
\end{corrige}

\newpage

\coursec{Exercices}

\coursec{Énergie électrique et puissance électrique}

\courssub{Grandeurs électriques, instruments de mesure et définitions}

\begin{definition}[Tension, courant, résistance]
Dans un circuit électrique en régime continu :
\begin{itemize}
\item La \cle{tension} $U$ (en volts, \SI{}{V}) est la différence de potentiel aux bornes d'un dipôle. Elle se mesure au \cle{voltmètre} branché \emph{en dérivation} (en parallèle).
\item L'\cle{intensité} du courant $I$ (en ampères, \SI{}{A}) est la quantité de charge traversant une section par seconde. Elle se mesure à l'\cle{ampèremètre} branché \emph{en série}.
\item La \cle{résistance} $R$ (en ohms, \SI{}{\ohm}) caractérise l'opposition au passage du courant. Pour un conducteur ohmique, la \cle{loi d'Ohm} s'applique : $U = R \cdot I$.
\end{itemize}
\end{definition}

\begin{propriete}[Puissance électrique reçue par un dipôle]
La \cle{puissance électrique} $P$ (en watts, \SI{}{W}) reçue par un dipôle soumis à une tension $U$ et traversé par un courant d'intensité $I$ est :
\[
P = U \cdot I
\]
Pour un \cle{résistor} (conducteur ohmique), en utilisant la loi d'Ohm ($U = R \cdot I$), on obtient deux formes équivalentes :
\[
P = R \cdot I^2 \qquad \text{et} \qquad P = \frac{U^2}{R}
\]
\emph{Au BEPC, il faut savoir utiliser les trois formules selon les grandeurs connues.}
\end{propriete}

\begin{aretenir}[Unités et instruments de mesure de la puissance et de l'énergie]
\begin{center}
\begin{tabularx}{\linewidth}{|l|X|X|}\hline
\rowcolor{popblueL} \textbf{Grandeurs} & \textbf{Unité légale (SI)} & \textbf{Instrument de mesure} \\ \hline
Tension $U$ & Volt (\SI{}{V}) & Voltmètre (branché en dérivation) \\ \hline
Intensité $I$ & Ampère (\SI{}{A}) & Ampèremètre (branché en série) \\ \hline
Puissance $P$ & Watt (\SI{}{W}) & Wattmètre (branché en série + dérivation) \\ \hline
Énergie $E$ & Joule (\SI{}{J}) & Compteur électrique (kWh) \\ \hline
\end{tabularx}
\end{center}
L'\cle{énergie électrique} $E$ reçue pendant une durée $t$ est : $E = P \cdot t$ (en joules si $P$ en watts et $t$ en secondes).
L'unité courante pour la facturation (ENEO) est le \cle{kilowattheure} (\SI{}{kWh}) : $1\,\text{kWh} = 3\,600\,000\,\text{J} = 3,6\,\text{MJ}$.
\end{aretenir}

\begin{attention}[Pièges classiques]
\begin{itemize}
\item \textbf{Confondre puissance et énergie :} La puissance est un débit (vitesse de transfert), l'énergie est une quantité. $P$ en W, $E$ en J ou kWh.
\item \textbf{Unités incohérentes :} Dans $E = P \cdot t$, si $P$ est en kW, $t$ doit être en h pour avoir $E$ en kWh. Si $P$ en W et $t$ en s, $E$ en J. \textbf{Convertis toujours avant de calculer.}
\item \textbf{Branchement voltmètre/ampèremètre :} Voltmètre \emph{toujours} en parallèle (dérivation) aux bornes du dipôle ; Ampèremètre \emph{toujours} en série dans la branche.
\item \textbf{Formule $P = R \cdot I^2$ :} Valable \emph{uniquement} pour un résistor (conducteur ohmique). Pas pour un moteur, une lampe LED, une pile...
\end{itemize}
\end{attention}

\begin{exemplebox}[Mesure de la puissance d'une lampe]
On alimente une lampe sous $U = \SI{12}{V}$. L'ampèremètre indique $I = \SI{0,5}{A}$.
\begin{enumerate}
\item Puissance : $P = U \cdot I = 12 \times 0,5 = \SI{6}{W}$.
\item Résistance (si ohmique) : $R = \frac{U}{I} = \frac{12}{0,5} = \SI{24}{\ohm}$.
\item Vérification : $P = R \cdot I^2 = 24 \times 0,25 = \SI{6}{W}$.
\end{enumerate}
\end{exemplebox}

\courssub{Puissance électrique d'un résistor : formules et applications}

\begin{methode}[Calculer la puissance d'un résistor]
\begin{enumerate}
\item Identifie les grandeurs connues ($U$, $I$, $R$).
\item Choisis la formule adaptée : $P = U \cdot I$ (toujours valable), $P = R \cdot I^2$ (si $R$ et $I$ connus), $P = U^2/R$ (si $U$ et $R$ connus).
\item Vérifie les unités : $U$ en V, $I$ en A, $R$ en $\Omega$ $\rightarrow$ $P$ en W.
\item Calcule et exprime le résultat avec son unité.
\end{enumerate}
\end{methode}

\begin{exoresolu}[Bouilloire et résistance chauffante]
Une résistance chauffante de bouilloire a une résistance $R = \SI{22}{\ohm}$. Elle est branchée sur le secteur $U = \SI{220}{V}$.
\begin{enumerate}
\item Calcule l'intensité $I$ du courant.
\item Calcule la puissance $P$ de deux façons différentes.
\end{enumerate}
\tcblower
\textbf{Solution.}
\begin{enumerate}
\item Loi d'Ohm : $I = \frac{U}{R} = \frac{220}{22} = \SI{10}{A}$.
\item Méthode 1 : $P = U \cdot I = 220 \times 10 = \SI{2200}{W}$.
Méthode 2 : $P = \frac{U^2}{R} = \frac{220^2}{22} = \frac{48400}{22} = \SI{2200}{W}$.
(Méthode 3 : $P = R \cdot I^2 = 22 \times 10^2 = \SI{2200}{W}$).
\end{enumerate}
\end{exoresolu}

\courssub{Énergie électrique, compteur et facture ENEO}

\begin{definition}[Énergie électrique consommée]
L'énergie électrique $E$ (en joules \SI{}{J}) consommée par un appareil de puissance $P$ (en watts \SI{}{W}) pendant une durée $t$ (en secondes \SI{}{s}) est :
\[
E = P \cdot t
\]
Pour la facturation domestique, on utilise le \cle{kilowattheure} (\SI{}{kWh}) :
\[
E\,(\text{kWh}) = P\,(\text{kW}) \times t\,(\text{h})
\]
Le \cle{compteur électrique} (compteur ENEO) totalise l'énergie consommée en kWh depuis sa mise en service.
\end{definition}

\begin{methode}[Calculer le coût de l'électricité]
\begin{enumerate}
\item Liste les appareils : puissance $P_i$ (en kW) et durée d'utilisation quotidienne $t_i$ (en h).
\item Énergie quotidienne : $E_{\text{jour}} = \sum P_i \cdot t_i$ (en kWh).
\item Énergie mensuelle : $E_{\text{mois}} = E_{\text{jour}} \times \text{nombre de jours}$.
\item Coût : $\text{Coût} = E_{\text{mois}} \times \text{tarif du kWh}$ (en FCFA).
\end{enumerate}
\end{methode}

\begin{savaistu}[Tarification ENEO au Cameroun]
ENEO applique un tarif progressif par tranches de consommation (kWh/mois) pour les abonnés basse tension (ménages). En 2024-2025, le tarif de base (première tranche) est d'environ \SI{110}{FCFA/kWh}. Il existe aussi un tarif \og heures creuses \fg{} (souvent \SI{70}{FCFA/kWh} entre 22h et 6h) pour les gros consommateurs (chauffe-eau, climatisation) équipés d'un compteur double tarif. Bien lire sa facture : elle comprend aussi la redevance d'abonnement, la TVA (19,25\%) et parfois des taxes communales.
\end{savaistu}

\courssub{Sécurité électrique : protections et prévention}

\begin{propriete}[Protections contre les surintensités]
\begin{itemize}
\item \textbf{Fusible / Disjoncteur divisionnaire (calibre $I_n$) :} Protège les \emph{fils} (câbles) contre l'échauffement dû à un courant trop fort (surcharge ou court-circuit). Il doit être choisi \textbf{inférieur ou égal} à l'intensité maximale admissible par le câble (ex: câble 2,5 mm² $\rightarrow$ 20 A max, disjoncteur 16 A ou 20 A).
\item \textbf{Règle :} $I_{\text{charge max}} \le I_n \le I_{\text{câble max}}$.
\end{itemize}
\end{propriete}

\begin{propriete}[Protection différentielle (Disjoncteur différentiel 30 mA)]
Le \cle{disjoncteur différentiel} (interrupteur différentiel) compare le courant aller et le courant retour. Si une différence (fuite vers la terre) dépasse \SI{30}{mA} (seuil de sensibilité pour la protection des personnes), il coupe \emph{instantanément} l'alimentation. \textbf{Obligatoire} en tête d'installation domestique (norme NF C 15-100 / norme camerounaise).
\end{propriete}

\begin{attention}[Danger des rallonges et multiprises]
\begin{itemize}
\item Une rallonge a un calibre maximal (souvent \SI{10}{A} ou \SI{16}{A} pour du 1,5 ou 2,5 mm²).
\item Brancher fer (\SI{1200}{W}) + bouilloire (\SI{2000}{W}) + four (\SI{1500}{W}) sur une même rallonge sous \SI{220}{V} :
$P_{\text{tot}} = \SI{4700}{W} \Rightarrow I = \frac{4700}{220} \approx \SI{21,4}{A}$.
\item Si la rallonge est prévue pour \SI{16}{A} max : \textbf{SURCHARGE GRAVE} $\rightarrow$ échauffement, fonte de l'isolant, risque d'incendie. \textbf{Interdit.}
\end{itemize}
\end{attention}

\begin{experience}[Analyse d'un schéma de mesure complet]
\textbf{Matériel :} Générateur réglable, résistor $R$, ampèremètre A, voltmètre V, fils.
\textbf{Protocole :}
\begin{enumerate}
\item Monte le circuit : Générateur $\rightarrow$ A $\rightarrow$ R $\rightarrow$ retour Générateur. V en dérivation aux bornes de R.
\item Varie la tension $U$ (ex: 6, 12, 18 V). Relevé $I$ correspondant.
\item Calcule $U/I$ et $P = U \times I$ pour chaque point.
\end{enumerate}
\textbf{Observations attendues :} $U/I$ constant $\rightarrow$ conducteur ohmique. $P$ croît avec $U$ (et $I$).
\textbf{Interprétation :} Vérification de la loi d'Ohm et des formules de puissance.
\end{experience}

\courssub{Entraînement et préparation au BEPC}

\courssub{Je vérifie mes connaissances}

\begin{exercice}[QCM : les bonnes réponses]
Pour chaque question, une seule réponse est exacte. Recopie la lettre correspondante.
\begin{enumerate}
\item L'unité de la puissance électrique est :
\begin{enumerate}
\item le joule (J) ; \item le watt (W) ; \item le kilowattheure (kWh).
\end{enumerate}
\item L'appareil qui mesure directement la puissance électrique reçue par un dipôle est :
\begin{enumerate}
\item le voltmètre ; \item l'ampèremètre ; \item le wattmètre.
\end{enumerate}
\item Un résistor de résistance $R = \SI{10}{\ohm}$ est traversé par un courant $I = \SI{2}{A}$. Sa puissance vaut :
\begin{enumerate}
\item $P = \SI{20}{W}$ ; \item $P = \SI{40}{W}$ ; \item $P = \SI{5}{W}$.
\end{enumerate}
\item Le compteur électrique d'ENEO mesure :
\begin{enumerate}
\item la tension du réseau ; \item l'énergie électrique consommée ; \item l'intensité maximale du courant.
\end{enumerate}
\end{enumerate}
\end{exercice}

\begin{exercice}[Vrai ou faux, en justifiant]
Pour chaque affirmation, réponds par \textbf{vrai} ou \textbf{faux}, puis justifie ta réponse en une phrase ou par un calcul.
\begin{enumerate}
\item La puissance électrique d'un résistor peut s'écrire $P = R \cdot I^2$.
\item Un kilowattheure correspond à l'énergie consommée par un appareil de $\SI{1000}{W}$ fonctionnant pendant une heure.
\item Plus la résistance d'un conducteur ohmique est grande (à intensité constante), plus la puissance dissipée est faible.
\item Le voltmètre se branche en série dans un circuit.
\end{enumerate}
\end{exercice}

\begin{exercice}[Définitions à compléter]
Recopie et complète les phrases suivantes avec les mots ou expressions : \emph{watt}, \emph{énergie}, \emph{wattmètre}, \emph{compteur}, \emph{$U \cdot I$}, \emph{joule}.
\begin{enumerate}
\item La puissance électrique reçue par un dipôle soumis à une tension $U$ et traversé par un courant d'intensité $I$ est $P = \ldots$
\item Dans le système international, la puissance s'exprime en \ldots{} et l'énergie en \ldots
\item Le \ldots{} mesure directement la puissance électrique ; le \ldots{} mesure l'\ldots{} électrique consommée par une installation.
\end{enumerate}
\end{exercice}

\begin{exercice}[Calculs directs]
\begin{enumerate}
\item Une lampe de chevet est alimentée sous $U = \SI{220}{V}$ et traversée par un courant $I = \SI{0,25}{A}$. Calcule sa puissance électrique $P$.
\item Un résistor de résistance $R = \SI{50}{\ohm}$ est traversé par un courant $I = \SI{0,4}{A}$. Calcule la puissance dissipée $P$ en utilisant $P = R \cdot I^2$.
\item Une bouilloire de puissance $P = \SI{2000}{W}$ fonctionne pendant $t = \SI{3}{min}$. Calcule l'énergie consommée en joules, puis en wattheures.
\end{enumerate}
\end{exercice}

\courssub{Je m'entraîne}

\begin{exercice}[Le fer à repasser de Maman]
Une famille de Douala utilise un fer à repasser de puissance $P = \SI{1200}{W}$ pendant $30$ minutes chaque jour.
\begin{enumerate}
\item Convertis la durée quotidienne d'utilisation en heures.
\item Calcule l'énergie électrique consommée chaque jour, en kilowattheures.
\item Calcule l'énergie consommée pendant un mois de $30$ jours.
\item Le kilowattheure est facturé $\SI{110}{FCFA}$ par ENEO. Calcule le coût mensuel du repassage.
\end{enumerate}
\end{exercice}

\begin{exercice}[Lecture d'un schéma de mesure]
Sur le schéma ci-dessous, un résistor est alimenté par un générateur. Un voltmètre et un ampèremètre permettent de mesurer la tension à ses bornes et l'intensité qui le traverse.
\begin{popfigure}
\begin{tikzpicture}[scale=1.0, every node/.style={font=\small}]
% fils du circuit
\draw[thick, popdark] (0,0) -- (6,0) -- (6,3) -- (0,3) -- (0,0);
% générateur
\draw[thick, popdark] (0,1.2) -- (0,1.8);
\draw[thick, popdark] (-0.25,1.5) -- (0.25,1.5);
\node[left, popdark] at (-0.1,1.5) {G};
% résistor
\draw[thick, popblue, fill=popblueL] (2.6,2.75) rectangle (3.4,3.25);
\node[above, popblue] at (3,3.3) {$R$};
% ampèremètre
\draw[thick, poporange, fill=poporangeL] (4.8,0) circle (0.4);
\node[poporange] at (4.8,0) {A};
% voltmètre en dérivation
\draw[thick, popgreen] (2.6,3) -- (2.6,4.2) -- (3.4,4.2) -- (3.4,3);
\draw[thick, popgreen, fill=popgreenL] (3,4.2) circle (0.4);
\node[popgreen] at (3,4.2) {V};
\end{tikzpicture}
\legende{Schéma de mesure : A en série, V en dérivation aux bornes de R.}
\end{popfigure}
L'ampèremètre indique $I = \SI{0,5}{A}$ et le voltmètre indique $U = \SI{12}{V}$.
\begin{enumerate}
\item Pourquoi le voltmètre est-il branché en dérivation aux bornes du résistor ?
\item Calcule la puissance électrique $P$ reçue par le résistor.
\item Déduis-en la valeur de la résistance $R$ du résistor.
\item Vérifie ton résultat en calculant $R \cdot I^2$ et en comparant avec $P$.
\end{enumerate}
\end{exercice}

\begin{exercice}[Le groupe électrogène du quartier]
Pendant une coupure d'ENEO, un commerçant de Yaoundé branche sur son groupe électrogène : un congélateur ($\SI{150}{W}$), trois lampes de $\SI{60}{W}$ chacune et un téléviseur ($\SI{80}{W}$). La tension du groupe est $U = \SI{220}{V}$.
\begin{enumerate}
\item Calcule la puissance totale $P_{\text{tot}}$ demandée au groupe.
\item Calcule l'intensité totale $I$ du courant débité par le groupe.
\item Le groupe fonctionne pendant $4$ heures. Calcule l'énergie fournie en kilowattheures.
\item Le groupe consomme $\SI{0,5}{L}$ d'essence par kilowattheure produit. Quel volume d'essence est nécessaire pour ces $4$ heures ?
\end{enumerate}
\end{exercice}

\begin{exercice}[QCM sécurité électrique]
Pour chaque situation, choisis la bonne réponse et justifie brièvement.
\begin{enumerate}
\item Pour protéger une ligne alimentant des appareils dont le courant total peut atteindre $\SI{20}{A}$, on choisit un fusible (ou disjoncteur) de calibre :
\begin{enumerate}
\item $\SI{10}{A}$ ; \item $\SI{20}{A}$ ; \item $\SI{50}{A}$.
\end{enumerate}
\item Le rôle principal d'un disjoncteur différentiel est de :
\begin{enumerate}
\item augmenter la tension du réseau ;
\item couper le circuit en cas de fuite de courant vers la terre ;
\item réduire la facture d'électricité.
\end{enumerate}
\item Une rallonge sur laquelle on branche un fer à repasser ($\SI{1200}{W}$), une bouilloire ($\SI{2000}{W}$) et un four ($\SI{1500}{W}$) sous $\SI{220}{V}$ est traversée par un courant d'environ :
\begin{enumerate}
\item $\SI{5}{A}$ ; \item $\SI{21}{A}$ ; \item $\SI{47}{A}$.
\end{enumerate}
\item Cette utilisation de la rallonge est-elle sans danger si elle est prévue pour $\SI{16}{A}$ maximum ?
\end{enumerate}
\end{exercice}

\courssub{Je me prépare au BEPC}

\begin{exercice}[Type BEPC : étude complète d'une résistance chauffante]
Une élève de 3ème réalise au laboratoire le montage ci-dessous pour étudier la résistance chauffante d'une bouilloire. Elle dispose d'un générateur de tension réglable, d'un ampèremètre, d'un voltmètre et de fils de connexion.
\begin{popfigure}
\begin{tikzpicture}[scale=1.0, every node/.style={font=\small}]
\draw[thick, popdark] (0,0) -- (7,0) -- (7,3) -- (0,3) -- (0,0);
\draw[thick, popdark] (0,1.2) -- (0,1.8);
\draw[thick, popdark] (-0.25,1.5) -- (0.25,1.5);
\node[left, popdark] at (-0.1,1.5) {G};
\draw[thick, popblue, fill=popblueL] (3.1,2.7) rectangle (3.9,3.3);
\node[above, popblue] at (3.5,3.35) {résistance chauffante};
\draw[thick, poporange, fill=poporangeL] (5.6,0) circle (0.4);
\node[poporange] at (5.6,0) {A};
\draw[thick, popgreen] (3.1,3) -- (3.1,4.3) -- (3.9,4.3) -- (3.9,3);
\draw[thick, popgreen, fill=popgreenL] (3.5,4.3) circle (0.4);
\node[popgreen] at (3.5,4.3) {V};
\end{tikzpicture}
\legende{Montage d'étude d'une résistance : A en série, V en dérivation.}
\end{popfigure}
\textbf{Partie A : mesures.} Pour trois valeurs de la tension, elle relève :
\begin{center}
\begin{tabular}{|c|c|c|c|}
\hline
\rowcolor{popblueL} $U$ (V) & 6 & 12 & 18 \\
\hline
$I$ (A) & 0,27 & 0,55 & 0,82 \\
\hline
\end{tabular}
\end{center}
\begin{enumerate}
\item Légende le schéma : indique le rôle de chaque appareil et le sens du courant.
\item Calcule le quotient $\dfrac{U}{I}$ pour chaque mesure. Que constates-tu ? Quelle loi cela illustre-t-il ?
\item Déduis-en la valeur de la résistance $R$ de la résistance chauffante.
\end{enumerate}
\textbf{Partie B : puissance et énergie.}
\begin{enumerate}
\setcounter{enumi}{3}
\item La bouilloire est maintenant branchée sur le secteur ($U = \SI{220}{V}$). Calcule l'intensité $I$ du courant qui la traverse.
\item Calcule la puissance électrique $P$ de la bouilloire de deux façons différentes.
\item Elle chauffe de l'eau pendant $t = \SI{5}{min}$. Calcule l'énergie électrique consommée en joules.
\item Exprime cette énergie en kilowattheures, puis calcule son coût au tarif de $\SI{110}{FCFA/kWh}$.
\end{enumerate}
\end{exercice}

\begin{exercice}[Type BEPC : le chauffe-eau et les heures creuses]
La courbe ci-dessous donne la puissance $P$ absorbée par le chauffe-eau d'une famille de Bafoussam au cours d'une journée de $24$ heures.
\begin{popfigure}
\begin{tikzpicture}
\begin{axis}[
width=0.85\linewidth, height=5.5cm,
xlabel={temps $t$ (h)}, ylabel={puissance $P$ (W)},
xmin=0, xmax=24, ymin=0, ymax=1600,
xtick={0,2,4,6,8,10,12,14,16,18,20,22,24},
ytick={0,400,800,1200,1600},
grid=both, grid style={popline!50},
axis line style={popdark},
]
\addplot[popcurve, domain=0:24, samples=200, const plot] coordinates {
(0,1200) (4,1200) (4,0) (12,0) (12,1200) (14,1200) (14,0) (20,0) (20,1200) (23,1200) (23,0) (24,0)
};
\end{axis}
\end{tikzpicture}
\legende{Profil de puissance du chauffe-eau sur 24 h (puissance constante 1200 W par paliers).}
\end{popfigure}
Le chauffe-eau fonctionne à puissance constante $P = \SI{1200}{W}$ pendant les périodes : de $0$ h à $4$ h, de $12$ h à $14$ h et de $20$ h à $23$ h.
\begin{enumerate}
\item D'après la courbe, pendant quelles périodes de la journée le chauffe-eau chauffe-t-il ? Quelle est la durée totale de fonctionnement sur $24$ h ?
\item Calcule l'énergie électrique consommée en une journée, en kilowattheures.
\item Au tarif normal de $\SI{110}{FCFA/kWh}$, quel est le coût journalier du chauffe-eau ?
\item ENEO propose un tarif \og heures creuses \fg{} de $\SI{70}{FCFA/kWh}$ entre $22$ h et $6$ h. La famille décide de ne faire fonctionner le chauffe-eau que de $23$ h à $6$ h (soit $7$ h, en supposant que cela suffit à chauffer l'eau). Calcule la nouvelle énergie consommée et le nouveau coût journalier.
\item Calcule l'économie réalisée sur un mois de $30$ jours.
\item Cite deux autres conseils pratiques permettant à cette famille de réduire sa facture d'électricité.
\end{enumerate}
\end{exercice}

\begin{exercice}[Type BEPC : mini-projet, tableau de bord de consommation de la salle de 3ème]
Ton professeur de PCT te confie une mission : suivre la consommation électrique de la salle de classe pendant un mois. La salle contient : $6$ lampes de $\SI{40}{W}$ chacune, $2$ ventilateurs de $\SI{75}{W}$ chacun et un ordinateur de $\SI{200}{W}$. La tension du réseau est $U = \SI{220}{V}$. Les appareils fonctionnent en moyenne $6$ heures par jour, $22$ jours par mois.
\begin{enumerate}
\item Calcule la puissance totale $P_{\text{tot}}$ des appareils de la salle lorsque tout fonctionne en même temps.
\item Calcule l'intensité totale $I$ du courant dans la ligne de la salle. Un disjoncteur de $\SI{10}{A}$ est-il suffisant ? Justifie.
\item Calcule l'énergie consommée par la salle en une journée, puis en un mois, en kilowattheures.
\item Au tarif de $\SI{110}{FCFA/kWh}$, calcule le coût mensuel de la consommation de la salle.
\item Le relevé du compteur chaque lundi donne : $12$ kWh, $14$ kWh, $11$ kWh, $13$ kWh. Calcule la consommation hebdomadaire moyenne et compare-la à ton estimation de la question 3. Propose une explication à l'écart observé.
\item Rédige un court texte (5 à 8 lignes) présentant deux mesures concrètes que la classe peut adopter pour réduire sa consommation (extinction des appareils, lampes LED, etc.) et estime l'économie réalisable si l'on remplace les $6$ lampes de $\SI{40}{W}$ par des lampes LED de $\SI{9}{W}$.
\end{enumerate}
\end{exercice}

\newpage

\coursec{Corrigés des exercices}

\begin{corrige}[Exercice 1 — QCM : les bonnes réponses]
\begin{enumerate}
\item \textbf{Réponse b : le watt (W).} Le joule (\si{\joule}) est l'unité de l'énergie, le kilowattheure (\si{\kilo\watt\hour}) est une unité d'énergie utilisée pour la facturation. La puissance s'exprime en watts (\si{\watt}).
\item \textbf{Réponse c : le wattmètre.} Le voltmètre mesure la tension, l'ampèremètre mesure l'intensité. Seul le wattmètre mesure directement la puissance.
\item \textbf{Réponse b : $P = \SI{40}{W}$.} Pour un résistor : $P = R \cdot I^2 = \SI{10}{\ohm} \times (\SI{2}{A})^2 = 10 \times 4 = \SI{40}{W}$.
\item \textbf{Réponse b : l'énergie électrique consommée.} Le compteur totalise l'énergie en kilowattheures (\si{\kilo\watt\hour}) depuis sa mise en service.
\end{enumerate}
\end{corrige}

\begin{corrige}[Exercice 2 — Vrai ou faux, en justifiant]
\begin{enumerate}
\item \textbf{Vrai.} Pour un résistor, la loi d'Ohm donne $U = R \cdot I$ ; en remplaçant dans $P = U \cdot I$, on obtient $P = R \cdot I \cdot I = R \cdot I^2$.
\item \textbf{Vrai.} $E = P \cdot t = \SI{1}{kW} \times \SI{1}{h} = \SI{1}{kWh}$. C'est précisément la définition du kilowattheure.
\item \textbf{Faux.} À intensité constante, $P = R \cdot I^2$ : la puissance est \emph{proportionnelle} à $R$. Plus $R$ est grande, plus la puissance dissipée est \emph{grande} (et non petite).
\item \textbf{Faux.} Le voltmètre se branche \emph{en dérivation} (en parallèle) aux bornes du dipôle. C'est l'ampèremètre qui se branche en série.
\end{enumerate}
\end{corrige}

\begin{corrige}[Exercice 3 — Définitions à compléter]
\begin{enumerate}
\item La puissance électrique reçue par un dipôle soumis à une tension $U$ et traversé par un courant d'intensité $I$ est $P = \mathbf{U \cdot I}$.
\item Dans le système international, la puissance s'exprime en \textbf{watt} (\si{\watt}) et l'énergie en \textbf{joule} (\si{\joule}).
\item Le \textbf{wattmètre} mesure directement la puissance électrique ; le \textbf{compteur} mesure l'\textbf{énergie} électrique consommée par une installation.
\end{enumerate}
\end{corrige}

\begin{corrige}[Exercice 4 — Calculs directs]
\begin{enumerate}
\item $P = U \cdot I = \SI{220}{V} \times \SI{0,25}{A} = \SI{55}{W}$.
\item $P = R \cdot I^2 = \SI{50}{\ohm} \times (\SI{0,4}{A})^2 = 50 \times \num{0,16} = \SI{8}{W}$.
\item Conversion de la durée : $t = \SI{3}{min} = 3 \times 60 = \SI{180}{s}$.
En joules : $E = P \cdot t = \SI{2000}{W} \times \SI{180}{s} = \SI{360000}{J} = \SI{360}{kJ}$.
En wattheures : $t = \SI{3}{min} = \dfrac{3}{60} = \SI{0,05}{h}$, donc $E = \SI{2000}{W} \times \SI{0,05}{h} = \SI{100}{Wh}$.
Vérification : $\SI{100}{Wh} = 100 \times 3600 = \SI{360000}{J}$. Les deux résultats concordent.
\end{enumerate}
\end{corrige}

\begin{corrige}[Exercice 5 — Le fer à repasser de Maman]
\begin{enumerate}
\item $t = \SI{30}{min} = \dfrac{30}{60} = \SI{0,5}{h}$.
\item Puissance en kilowatts : $P = \SI{1200}{W} = \SI{1,2}{kW}$.
$E_{\text{jour}} = P \cdot t = 1,2 \times 0,5 = \SI{0,6}{kWh}$.
\item $E_{\text{mois}} = 0,6 \times 30 = \SI{18}{kWh}$.
\item $\text{Coût} = 18 \times 110 = \SI{1980}{\text{FCFA}}$.
\end{enumerate}
\textbf{Conclusion :} le repassage coûte environ \SI{1980}{\text{FCFA}} par mois à cette famille.
\end{corrige}

\begin{corrige}[Exercice 6 — Lecture d'un schéma de mesure]
\begin{enumerate}
\item Le voltmètre mesure la \emph{différence de potentiel} entre deux points : il doit donc être branché \textbf{en dérivation} (en parallèle) aux bornes du résistor, pour être soumis à la même tension que lui. Sa résistance interne très grande l'empêche de perturber le circuit.
\item $P = U \cdot I = \SI{12}{V} \times \SI{0,5}{A} = \SI{6}{W}$.
\item D'après la loi d'Ohm : $R = \dfrac{U}{I} = \dfrac{12}{0,5} = \SI{24}{\ohm}$.
\item Vérification : $R \cdot I^2 = 24 \times (\num{0,5})^2 = 24 \times \num{0,25} = \SI{6}{W} = P$. Le résultat est cohérent : le dipôle se comporte bien comme un conducteur ohmique.
\end{enumerate}
\end{corrige}

\begin{corrige}[Exercice 7 — Le groupe électrogène du quartier]
\begin{enumerate}
\item Puissance des trois lampes : $3 \times 60 = \SI{180}{W}$.
$P_{\text{tot}} = 150 + 180 + 80 = \SI{410}{W}$.
\item $I = \dfrac{P_{\text{tot}}}{U} = \dfrac{410}{220} \approx \SI{1,86}{A}$.
\item $P_{\text{tot}} = \SI{0,41}{kW}$ ; $E = P \cdot t = 0,41 \times 4 = \SI{1,64}{kWh}$.
\item Volume d'essence : $V = 1,64 \times 0,5 = \SI{0,82}{L}$ (consommation spécifique admise : \SI{0,5}{L/kWh} pour ce groupe).
\end{enumerate}
\textbf{Conclusion :} il faut environ \SI{0,82}{L} d'essence pour alimenter ces appareils pendant 4 heures.
\end{corrige}

\begin{corrige}[Exercice 8 — QCM sécurité électrique]
\begin{enumerate}
\item \textbf{Réponse b : \SI{20}{A}.} Le calibre du fusible doit être au moins égal au courant maximal de la ligne (sinon il coupe sans raison) et inférieur ou égal au courant admissible par le câble. \SI{10}{A} couperait dès qu'on dépasse \SI{10}{A} ; \SI{50}{A} ne protégerait plus le câble.
\item \textbf{Réponse b : couper le circuit en cas de fuite de courant vers la terre.} Le différentiel compare le courant aller et le courant retour ; dès qu'une fuite dépasse \SI{30}{mA} (par exemple à travers le corps d'une personne), il coupe l'alimentation et protège les personnes contre l'électrisation.
\item \textbf{Réponse b : \SI{21}{A}.} $P_{\text{tot}} = 1200 + 2000 + 1500 = \SI{4700}{W}$ ; $I = \dfrac{4700}{220} \approx \SI{21,4}{A}$.
\item \textbf{Non, c'est dangereux.} Le courant réel ($\approx \SI{21}{A}$) dépasse largement le calibre de la rallonge (\SI{16}{A}). Les fils vont surchauffer, l'isolant peut fondre : risque d'incendie. Il faut brancher ces appareils puissants sur des prises murales distinctes.
\end{enumerate}
\end{corrige}

\begin{corrige}[Exercice 9 — Type BEPC : étude complète d'une résistance chauffante]
\textbf{Partie A : mesures.}
\begin{enumerate}
\item Le générateur G fournit la tension ; l'ampèremètre A, branché \emph{en série}, mesure l'intensité $I$ traversant la résistance ; le voltmètre V, branché \emph{en dérivation}, mesure la tension $U$ à ses bornes. Le courant conventionnel sort de la borne $+$ du générateur, traverse A puis la résistance, et revient à la borne $-$.
\item Calculs des quotients :
\[
\frac{6}{\num{0,27}} \approx 22,2 \ ; \qquad \frac{12}{\num{0,55}} \approx 21,8 \ ; \qquad \frac{18}{\num{0,82}} \approx 22,0.
\]
Le quotient $\dfrac{U}{I}$ est \textbf{constant} (environ $22$, aux erreurs de mesure près). Cela illustre la \textbf{loi d'Ohm} : le dipôle est un conducteur ohmique.
\item $R = \dfrac{U}{I} \approx \SI{22}{\ohm}$.
\end{enumerate}
\textbf{Partie B : puissance et énergie.}
\begin{enumerate}
\setcounter{enumi}{3}
\item Loi d'Ohm : $I = \dfrac{U}{R} = \dfrac{220}{22} = \SI{10}{A}$.
\item Première méthode : $P = U \cdot I = 220 \times 10 = \SI{2200}{W}$.
Deuxième méthode : $P = R \cdot I^2 = 22 \times 10^2 = \SI{2200}{W}$.
(On peut aussi écrire $P = \dfrac{U^2}{R} = \dfrac{48400}{22} = \SI{2200}{W}$.) Les trois formules donnent le même résultat.
\item $t = \SI{5}{min} = 5 \times 60 = \SI{300}{s}$ ; $E = P \cdot t = 2200 \times 300 = \SI{660000}{J}$.
\item En kilowattheures : $E = \dfrac{660000}{3\,600\,000} \approx \SI{0,183}{kWh}$ (vérification : $P = \SI{2,2}{kW}$, $t = \frac{5}{60}\,\text{h}$, $E = 2,2 \times \frac{5}{60} \approx \SI{0,183}{kWh}$).
Coût : $0,183 \times 110 \approx \SI{20,2}{\text{FCFA}}$.
\end{enumerate}
\textbf{Barème indicatif (sur 10) :} légende et sens du courant : 1,5 pt ; calculs des quotients et constat : 2 pts ; valeur de $R$ : 1 pt ; intensité : 1 pt ; puissance par deux méthodes : 2 pts ; énergie en J : 1 pt ; conversion en kWh et coût : 1,5 pt.
\textbf{Points de vigilance :} citer explicitement la loi d'Ohm ; convertir les minutes en secondes \emph{avant} d'appliquer $E = P \cdot t$ ; préciser que $P = R \cdot I^2$ n'est valable que pour un résistor ; arrondir proprement (\SI{20}{\text{FCFA}} environ).
\end{corrige}

\begin{corrige}[Exercice 10 — Type BEPC : le chauffe-eau et les heures creuses]
\begin{enumerate}
\item D'après la courbe, le chauffe-eau chauffe de $0$ h à $4$ h ($4$ h), de $12$ h à $14$ h ($2$ h) et de $20$ h à $23$ h ($3$ h). Durée totale : $4 + 2 + 3 = \SI{9}{h}$.
\item $P = \SI{1200}{W} = \SI{1,2}{kW}$ ; $E = 1,2 \times 9 = \SI{10,8}{kWh}$ par jour.
\item Coût journalier : $10,8 \times 110 = \SI{1188}{\text{FCFA}}$.
\item Nouvelle durée : $7$ h, toutes en heures creuses. $E' = 1,2 \times 7 = \SI{8,4}{kWh}$. Nouveau coût : $8,4 \times 70 = \SI{588}{\text{FCFA}}$.
\item Économie journalière : $1188 - 588 = \SI{600}{\text{FCFA}}$. Sur $30$ jours : $600 \times 30 = \SI{18000}{\text{FCFA}}$.
\item Conseils possibles (deux au choix) : régler le thermostat du chauffe-eau à $55$--$60\,^\circ\text{C}$ et non au maximum ; éteindre les lampes et appareils en veille ; remplacer les lampes à incandescence par des LED ; utiliser le fer à repasser en une seule session (le réchauffage répété gaspille de l'énergie) ; dégivrer régulièrement le congélateur.
\end{enumerate}
\textbf{Barème indicatif (sur 10) :} lecture de la courbe et durée totale : 2 pts ; énergie journalière : 1,5 pt ; coût tarif normal : 1 pt ; nouvelle énergie et nouveau coût : 2 pts ; économie mensuelle : 1,5 pt ; deux conseils pertinents : 2 pts.
\textbf{Points de vigilance :} bien convertir $P$ en kW pour obtenir $E$ en kWh ; lire correctement les paliers de la courbe (fonction en escalier) ; vérifier que la plage $23$ h--$6$ h est entièrement en heures creuses ($22$ h--$6$ h) ; justifier les conseils par un lien avec la puissance ou la durée d'utilisation.
\end{corrige}

\begin{corrige}[Exercice 11 — Type BEPC : mini-projet, tableau de bord de consommation]
\begin{enumerate}
\item Lampes : $6 \times 40 = \SI{240}{W}$ ; ventilateurs : $2 \times 75 = \SI{150}{W}$ ; ordinateur : \SI{200}{W}.
$P_{\text{tot}} = 240 + 150 + 200 = \SI{590}{W}$.
\item $I = \dfrac{P_{\text{tot}}}{U} = \dfrac{590}{220} \approx \SI{2,68}{A}$.
Comme $2,68 < 10$, le disjoncteur de \SI{10}{A} est \textbf{largement suffisant} : il ne se déclenchera pas en fonctionnement normal et protège correctement la ligne.
\item Énergie journalière : $E_{\text{jour}} = 0,59 \times 6 = \SI{3,54}{kWh}$.
Énergie mensuelle : $E_{\text{mois}} = 3,54 \times 22 = \SI{77,88}{kWh} \approx \SI{78}{kWh}$.
\item Coût mensuel : $77,88 \times 110 \approx \SI{8567}{\text{FCFA}}$ (environ \SI{8570}{\text{FCFA}}).
\item Consommation hebdomadaire moyenne relevée : $\dfrac{12 + 14 + 11 + 13}{4} = \dfrac{50}{4} = \SI{12,5}{kWh}$.
Estimation théorique hebdomadaire (5 jours de cours) : $3,54 \times 5 = \SI{17,7}{kWh}$.
Le relevé réel (\SI{12,5}{kWh}) est inférieur à l'estimation (\SI{17,7}{kWh}). Explication possible : les appareils ne fonctionnent pas tous en même temps ni pendant 6 h pleines (ventilateurs éteints par temps frais, ordinateur peu utilisé, lampes éteintes quand la lumière naturelle suffit).
\item \textbf{Texte attendu (exemple) :} Pour réduire la consommation de la salle, la classe peut d'abord éteindre systématiquement les lampes et les ventilateurs pendant les récréations et à la fin des cours, en désignant un élève responsable chaque semaine. Ensuite, elle peut profiter au maximum de la lumière naturelle en ouvrant les volets le matin. Enfin, le remplacement des lampes classiques par des lampes LED est très rentable.
\textbf{Estimation de l'économie LED :} nouvelle puissance des lampes : $6 \times 9 = \SI{54}{W}$ au lieu de \SI{240}{W}, soit une puissance économisée de $240 - 54 = \SI{186}{W} = \SI{0,186}{kW}$.
Énergie économisée par mois : $0,186 \times 6 \times 22 \approx \SI{24,6}{kWh}$.
Économie en FCFA : $24,6 \times 110 \approx \SI{2700}{\text{FCFA}}$ par mois.
\end{enumerate}
\textbf{Barème indicatif (sur 10) :} puissance totale : 1,5 pt ; intensité et avis justifié sur le disjoncteur : 2 pts ; énergies journalière et mensuelle : 2 pts ; coût mensuel : 1 pt ; moyenne hebdomadaire, comparaison et explication : 2 pts ; texte et calcul de l'économie LED : 1,5 pt.
\textbf{Points de vigilance :} additionner les puissances \emph{avant} de calculer l'intensité ; justifier le choix du disjoncteur par une comparaison chiffrée ($I < I_n$) ; convertir les watts en kilowatts pour obtenir des kWh ; pour la question 5, comparer des durées cohérentes (semaine de 5 jours de classe) ; rédiger le texte avec des mesures concrètes et chiffrées.
\end{corrige}

\newpage

\coursec{Fiche bilan}

\begin{popfigure}
\begin{center}\tcbox[colback=popblue,colframe=popblue,coltext=white,arc=9pt,boxsep=4pt,left=6pt,right=6pt]{\bfseries\small Énergie et puissance électriques}\end{center}
\vspace{-2pt}
\begin{tcbraster}[raster columns=3,raster equal height=rows,raster column skip=6pt,raster row skip=6pt,enhanced,blankest]
\begin{tcolorbox}[enhanced,colback=white,colframe=poporange,boxrule=0.9pt,arc=3pt,title={Grandeurs et mesures},fonttitle=\bfseries\scriptsize,coltitle=white,colbacktitle=poporange,left=4pt,right=4pt,top=2pt,bottom=2pt,boxsep=1pt,toptitle=1.5pt,bottomtitle=1.5pt,halign=left]
\begin{itemize}[nosep,leftmargin=*,itemsep=1pt,font=\scriptsize]
\item Voltmètre : $U$ en V
\item Ampèremètre : $I$ en A
\end{itemize}
\end{tcolorbox}
\begin{tcolorbox}[enhanced,colback=white,colframe=popgreen,boxrule=0.9pt,arc=3pt,title={Puissance $P$},fonttitle=\bfseries\scriptsize,coltitle=white,colbacktitle=popgreen,left=4pt,right=4pt,top=2pt,bottom=2pt,boxsep=1pt,toptitle=1.5pt,bottomtitle=1.5pt,halign=left]
\begin{itemize}[nosep,leftmargin=*,itemsep=1pt,font=\scriptsize]
\item $P=U\cdot I$
\item $P=R\cdot I^2$
\end{itemize}
\end{tcolorbox}
\begin{tcolorbox}[enhanced,colback=white,colframe=poppurple,boxrule=0.9pt,arc=3pt,title={Énergie $E$},fonttitle=\bfseries\scriptsize,coltitle=white,colbacktitle=poppurple,left=4pt,right=4pt,top=2pt,bottom=2pt,boxsep=1pt,toptitle=1.5pt,bottomtitle=1.5pt,halign=left]
\begin{itemize}[nosep,leftmargin=*,itemsep=1pt,font=\scriptsize]
\item $E=P\cdot t$
\item Unité : kWh
\end{itemize}
\end{tcolorbox}
\begin{tcolorbox}[enhanced,colback=white,colframe=poppink,boxrule=0.9pt,arc=3pt,title={Compteur et wattmètre},fonttitle=\bfseries\scriptsize,coltitle=white,colbacktitle=poppink,left=4pt,right=4pt,top=2pt,bottom=2pt,boxsep=1pt,toptitle=1.5pt,bottomtitle=1.5pt,halign=left]
\begin{itemize}[nosep,leftmargin=*,itemsep=1pt,font=\scriptsize]
\item Facture ENEO
\end{itemize}
\end{tcolorbox}
\begin{tcolorbox}[enhanced,colback=white,colframe=popgold,boxrule=0.9pt,arc=3pt,title={Sécurité},fonttitle=\bfseries\scriptsize,coltitle=white,colbacktitle=popgold,left=4pt,right=4pt,top=2pt,bottom=2pt,boxsep=1pt,toptitle=1.5pt,bottomtitle=1.5pt,halign=left]
\begin{itemize}[nosep,leftmargin=*,itemsep=1pt,font=\scriptsize]
\item Fusible, disjoncteur
\item Terre, pas d'eau
\end{itemize}
\end{tcolorbox}
\begin{tcolorbox}[enhanced,colback=white,colframe=popmuted,boxrule=0.9pt,arc=3pt,title={Applications},fonttitle=\bfseries\scriptsize,coltitle=white,colbacktitle=popmuted,left=4pt,right=4pt,top=2pt,bottom=2pt,boxsep=1pt,toptitle=1.5pt,bottomtitle=1.5pt,halign=left]
\begin{itemize}[nosep,leftmargin=*,itemsep=1pt,font=\scriptsize]
\item Groupe électrogène
\item Lampes, fer à repasser
\end{itemize}
\end{tcolorbox}
\end{tcbraster}

\legende{Carte mentale de la leçon : les six idées clés autour de l'énergie et de la puissance électriques.}
\end{popfigure}

\begin{bilan}
\textbf{Définitions clés}
\begin{itemize}
  \item La \cle{tension électrique} $U$ se mesure en \cle{volts} (V) avec un \cle{voltmètre} branché en \textbf{dérivation}.
  \item L'\cle{intensité} $I$ se mesure en \cle{ampères} (A) avec un \cle{ampèremètre} branché en \textbf{série}.
  \item La \cle{puissance électrique} $P$ d'un appareil est l'énergie qu'il consomme (ou transforme) par unité de temps ; elle se mesure en \cle{watts} (W) avec un \cle{wattmètre}.
  \item L'\cle{énergie électrique} $E$ consommée se mesure en joules (J) ou en \cle{kilowattheures} (kWh) grâce au \cle{compteur électrique} (compteur ENEO).
  \item Un \cle{résistor} (conducteur ohmique) est un dipôle qui transforme toute l'énergie électrique reçue en chaleur : c'est l'\cle{effet Joule}.
\end{itemize}

\textbf{Formules à connaître par cœur}
\begin{center}
\begin{tabularx}{\linewidth}{|l|X|l|}\hline
\rowcolor{popblueL} \textbf{Formule} & \textbf{Signification} & \textbf{Unités} \\ \hline
$P = U \cdot I$ & Puissance reçue par un dipôle & $P$ en W, $U$ en V, $I$ en A \\ \hline
$P = R \cdot I^2$ & Puissance dissipée par un résistor (effet Joule) & $R$ en $\Omega$ \\ \hline
$E = P \cdot t$ & Énergie consommée pendant la durée $t$ & $E$ en J si $t$ en s ; $E$ en kWh si $P$ en kW et $t$ en h \\ \hline
\SI{1}{kWh} $=$ \SI{3 600 000}{J} & Conversion énergie & \SI{1}{kW} $=$ \SI{1000}{W} \\ \hline
\end{tabularx}
\end{center}

\textbf{Méthodes express}
\begin{enumerate}
  \item \textbf{Calculer une puissance} : relever $U$ et $I$, puis appliquer $P = U \cdot I$. Pour un résistor, on peut aussi utiliser $P = R \cdot I^2$.
  \item \textbf{Calculer une énergie} : convertir $P$ en kW et $t$ en heures, puis $E = P \cdot t$ en kWh.
  \item \textbf{Estimer une facture} : multiplier l'énergie (kWh) par le prix du kWh.
  \item \textbf{Vérifier une installation} : additionner les puissances des appareils branchés et comparer à la puissance maximale supportée par la rallonge ou le disjoncteur.
\end{enumerate}

\textbf{Sécurité}
\begin{itemize}
  \item Ne jamais toucher une prise ou un appareil avec les mains mouillées.
  \item Ne pas surcharger une rallonge : risque d'échauffement et d'incendie.
  \item Le \cle{fusible} et le \cle{disjoncteur} coupent le circuit en cas de surintensité ; la \cle{prise de terre} protège des fuites de courant.
\end{itemize}
\end{bilan}

\begin{aretenir}[Checklist avant le Bac]
\begin{itemize}
  \item $\square$ Je sais définir tension, intensité, puissance et énergie électriques avec leurs unités.
  \item $\square$ Je sais quel instrument mesure chaque grandeur (voltmètre, ampèremètre, wattmètre, compteur).
  \item $\square$ Je sais que le voltmètre se branche en dérivation et l'ampèremètre en série.
  \item $\square$ Je connais par cœur $P = U \cdot I$ et je sais l'appliquer avec les bonnes unités.
  \item $\square$ Je connais $P = R \cdot I^2$ et je sais qu'elle ne s'applique qu'à un résistor.
  \item $\square$ Je sais calculer une énergie avec $E = P \cdot t$ en joules et en kilowattheures.
  \item $\square$ Je sais convertir : \SI{1}{kWh} $=$ \SI{3 600 000}{J} et \SI{1}{kW} $=$ \SI{1000}{W}.
  \item $\square$ Je sais calculer le coût d'une consommation électrique (facture ENEO).
  \item $\square$ Je sais lire la puissance nominale indiquée sur un appareil (lampe, fer, groupe électrogène).
  \item $\square$ Je connais le rôle du fusible, du disjoncteur et de la prise de terre.
  \item $\square$ Je sais citer au moins trois consignes de sécurité électrique à la maison.
  \item $\square$ Je vérifie toujours les unités et les conversions avant de conclure un calcul.
\end{itemize}
\end{aretenir}

\findecours

\end{document}
